% Le problème des minorités — Allemand Tle A — Cours signé OnBuch+
% Programme officiel MINESEC. Compiler avec : tectonic AL04-le-probleme-des-minorites.tex
\newcommand{\DOCMATIERE}{Allemand}
\newcommand{\DOCNIVEAU}{Tle A}
\newcommand{\DOCMODULE}{Chapitre 1 : Vie familiale et sociale}
\newcommand{\DOCLECON}{AL04}
\newcommand{\DOCTITRE}{Le problème des minorités}
% OnBuch+ — préambule « cours signés » ENRICHI (compilé avec tectonic / XeLaTeX)
% Système visuel « Pop » d'OnBuch+ : fond crème (jamais blanc), bordures encre
% épaisses, ombres « dures » décalées, titres Archivo Black en capitales.
% Version MATHÉMATIQUES : graphes (pgfplots/tikz), boîtes pédagogiques
% (définition, propriété, méthode, attention, le savais-tu, exercice résolu,
% exercice, TP, bilan), page de garde et signature OnBuch+ ancrée en bas.
%
% Macros attendues dans main.tex AVANT \input{preamble} :
%   \DOCMATIERE  \DOCNIVEAU  \DOCMODULE  \DOCLECON (numéro)  \DOCTITRE
\documentclass[11pt]{article}

\usepackage{fontspec}
\usepackage[ngerman,french]{babel} % français = langue principale ; l'allemand via \all{..} / textebox
\usepackage[a4paper,margin=2cm,top=3.4cm,bottom=2.8cm,headheight=1.4cm,footskip=1.3cm]{geometry}
\usepackage{amsmath,amssymb,amsfonts}
\usepackage{mathtools} % \xmapsto, \coloneqq... souvent utilisés par les modèles
\usepackage{cancel} % simplifications barrées (bilans de forces)
% \abs{z} (module, valeur absolue) et \norm{u} (norme), utilisés par les modèles
\providecommand{\abs}{}\let\abs\relax\DeclarePairedDelimiter{\abs}{\lvert}{\rvert}
\providecommand{\norm}{}\let\norm\relax\DeclarePairedDelimiter{\norm}{\lVert}{\rVert}
\usepackage[table]{xcolor} % [table] : fournit aussi \rowcolors (alternance de lignes)
\usepackage{pagecolor}
\usepackage{tikz}
\usetikzlibrary{arrows.meta,positioning,calc,shapes.geometric,shapes.misc,shapes.arrows,decorations.pathmorphing,decorations.pathreplacing,decorations.markings,patterns,shadows,mindmap,trees,fit,backgrounds,matrix,intersections,angles,quotes}
\usepackage{pgfplots}
\usepackage[version=4]{mhchem} % équations nucléaires \ce{^{235}_{92}U}
\usepackage[european,siunitx]{circuitikz} % schémas électriques
\pgfplotsset{compat=1.18}
\usepgfplotslibrary{fillbetween,groupplots} % aires entre courbes (calcul intégral)
\usepackage{enumitem}
\usepackage{fancyhdr}
% [most] charge toutes les bibliothèques tcolorbox (listings, minted, raster,
% documentation...) et gonfle inutilement la mémoire TeX sur un document long
% (~300 boîtes) : on ne charge que ce qui est réellement utilisé.
\usepackage[skins,breakable]{tcolorbox}
\usepackage{graphicx}
\usepackage{colortbl}
\usepackage{booktabs}
\usepackage{tabularx}
\usepackage{diagbox} % case d'en-tête diagonale des tableaux à double entrée
% Colonnes extensibles centrée / gauche / droite (C, L, R), souvent utilisées par les modèles
\newcolumntype{C}{>{\centering\arraybackslash}X}
\newcolumntype{L}{>{\raggedright\arraybackslash}X}
\newcolumntype{R}{>{\raggedleft\arraybackslash}X}
\usepackage{array}
\usepackage{multicol}
\usepackage{multirow}
\usepackage{siunitx}
% parse-numbers=false : accepte aussi \SI{2,5 \times 10^{-3}}{...}
\sisetup{locale=FR,output-decimal-marker={,},group-minimum-digits=4,parse-numbers=false}
\DeclareSIUnit\ohm{\ensuremath{\symup{\Omega}}} % U+2126 absent de la police
\DeclareSIUnit\division{div} % réglages d'oscilloscope (ms/div, V/div)
\DeclareSIUnit\tr{tr} % tours (vitesse de rotation, ex. \SI{1500}{\tr\per\minute})
\DeclareSIUnit\molar{M} % concentration molaire (ex. \SI{3}{\molar}, \SI{1}{\micro\molar})
\DeclareSIUnit\an{an} % année (le modèle confond parfois avec \year, un compteur TeX) — ex. \SI{5}{\kilo\meter\per\an}
\DeclareSIUnit\FCFA{FCFA} % franc CFA (ex. \SI{500}{\FCFA\per\kilo\gram})
\DeclareSIUnit\degreeBrix{\textdegree Brix} % teneur en sucre d'un sirop/jus (réfractométrie)
\DeclareSIUnit\month{mois} % le modèle utilise parfois \month, un compteur TeX, comme unité de durée
\usepackage{qrcode}
% \degree : commande fréquemment utilisée par les modèles pour un angle ou une
% température en dehors de \SI{}{} (non fournie par les paquets chargés).
\providecommand{\degree}{^{\circ}}
% \qed (fin de démonstration), utilisé par les modèles sans amsthm
\providecommand{\qed}{\hfill$\square$}
% \barre : double barre des valeurs interdites dans un tableau de variations (tkz-tab)
\providecommand{\barre}{\ensuremath{\|}}
% environnement proof (amsthm non chargé) utilisé par les modèles
\ifdefined\proof\else\newenvironment{proof}[1][Démonstration]{\par\noindent\textit{#1.}\ }{\qed\par}\fi
\usepackage{lastpage}
\usepackage{hyperref}
\hypersetup{hidelinks}

% ============================================================
% Palette « Pop » OnBuch+ — mêmes hex que la classe P (plus_ui.dart)
% ============================================================
\definecolor{poporange}{HTML}{F59321}
\definecolor{popdark}{HTML}{E07A0C}
\definecolor{popcream}{HTML}{FFF6E8}
\definecolor{popink}{HTML}{1C1714}
\definecolor{poppurple}{HTML}{7A5AE0}
\definecolor{popgreen}{HTML}{2E9E6A}
\definecolor{poppink}{HTML}{E0457B}
\definecolor{popblue}{HTML}{2F7DE1}
\definecolor{popgold}{HTML}{C98A12}
\definecolor{popmuted}{HTML}{8A7F76}
\definecolor{popline}{HTML}{EADFD2}
\definecolor{popgray}{HTML}{8A7F76}
\colorlet{popgrey}{popgray}
% Alias tolérants (le modèle nomme parfois les couleurs différemment)
\colorlet{popwhite}{white}
\colorlet{popblack}{popink}
\colorlet{popred}{poppink}
\colorlet{popyellow}{popgold}
% Teintes claires (fonds de tableaux/encadrés) — le modèle nomme parfois
% les couleurs avec un suffixe L (ex. popblueL) sans les avoir définies.
\colorlet{poporangeL}{poporange!20}
\colorlet{popdarkL}{popdark!20}
\colorlet{poppurpleL}{poppurple!20}
\colorlet{popgreenL}{popgreen!20}
\colorlet{poppinkL}{poppink!20}
\colorlet{popblueL}{popblue!20}
\colorlet{popgoldL}{popgold!20}
\colorlet{popredL}{popred!20}
\colorlet{popgrayL}{popgray!20}
% Teintes claires pour fonds de tableaux / zones
\colorlet{popblueL}{popblue!10!white}
\colorlet{poporangeL}{poporange!14!white}
\colorlet{popgreenL}{popgreen!12!white}
\colorlet{poppurpleL}{poppurple!12!white}
\colorlet{poppinkL}{poppink!12!white}
\colorlet{popgoldL}{popgold!14!white}

\pagecolor{popcream}

% ============================================================
% Typographie
% ============================================================
\defaultfontfeatures{Ligatures=TeX}
\setmainfont{PlusJakartaSans-Medium.ttf}[Path=../../../fonts/,BoldFont=PlusJakartaSans-Bold.ttf]
% Maths en sans-serif (Fira Math) avec chiffres et lettres droites de Plus
% Jakarta, pour que les formules se fondent dans le texte.
\usepackage[math-style=ISO,bold-style=ISO]{unicode-math}
\setmathfont{FiraMath-Regular.otf}[Path=../../../fonts/]
\setmathfont{PlusJakartaSans-Medium.ttf}[Path=../../../fonts/,range={up/num,up/{Latin,latin},\mathup}]
\usepackage{icomma} % virgule décimale sans espace en mode math
% Caractères mathématiques Unicode absents des polices de texte (Plus Jakarta,
% Archivo Black) : rendus par leur équivalent mathématique, en texte comme en
% formule — sinon ils s'affichent en carré vide (ex. « Arithmétique dans ℤ »).
\usepackage{newunicodechar}
\newunicodechar{ℝ}{\ensuremath{\mathbb{R}}}
\newunicodechar{ℤ}{\ensuremath{\mathbb{Z}}}
\newunicodechar{ℂ}{\ensuremath{\mathbb{C}}}
\newunicodechar{ℕ}{\ensuremath{\mathbb{N}}}
\newunicodechar{ℚ}{\ensuremath{\mathbb{Q}}}
\newunicodechar{𝔻}{\ensuremath{\mathbb{D}}}
\newunicodechar{α}{\ensuremath{\alpha}}
\newunicodechar{β}{\ensuremath{\beta}}
\newunicodechar{γ}{\ensuremath{\gamma}}
\newunicodechar{δ}{\ensuremath{\delta}}
\newunicodechar{ε}{\ensuremath{\varepsilon}}
\newunicodechar{θ}{\ensuremath{\theta}}
\newunicodechar{λ}{\ensuremath{\lambda}}
\newunicodechar{μ}{\ensuremath{\mu}}
\newunicodechar{σ}{\ensuremath{\sigma}}
\newunicodechar{τ}{\ensuremath{\tau}}
\newunicodechar{φ}{\ensuremath{\varphi}}
\newunicodechar{ψ}{\ensuremath{\psi}}
\newunicodechar{ω}{\ensuremath{\omega}}
\newunicodechar{Σ}{\ensuremath{\Sigma}}
% majuscules produites par \MakeUppercase dans les titres (α-aminés -> Α)
\newunicodechar{Α}{\ensuremath{\alpha}}
\newunicodechar{Β}{\ensuremath{\beta}}
\newunicodechar{Γ}{\ensuremath{\Gamma}}
\newunicodechar{Δ}{\ensuremath{\Delta}}
\newunicodechar{Ε}{\ensuremath{\varepsilon}}
\newunicodechar{Θ}{\ensuremath{\Theta}}
\newunicodechar{Λ}{\ensuremath{\Lambda}}
\newunicodechar{Μ}{\ensuremath{\mu}}
\newunicodechar{Τ}{\ensuremath{\tau}}
\newunicodechar{Φ}{\ensuremath{\Phi}}
\newunicodechar{Ψ}{\ensuremath{\Psi}}
\newunicodechar{Ω}{\ensuremath{\Omega}}
\newunicodechar{↦}{\ensuremath{\mapsto}}
\newunicodechar{∈}{\ensuremath{\in}}
\newunicodechar{∉}{\ensuremath{\notin}}
\newunicodechar{∧}{\ensuremath{\wedge}}
\newunicodechar{⇔}{\ensuremath{\Leftrightarrow}}
\newunicodechar{⟺}{\ensuremath{\Longleftrightarrow}}
\newunicodechar{⇒}{\ensuremath{\Rightarrow}}
\newunicodechar{⟹}{\ensuremath{\Longrightarrow}}
\newunicodechar{∘}{\ensuremath{\circ}}
\newunicodechar{∪}{\ensuremath{\cup}}
\newunicodechar{∩}{\ensuremath{\cap}}
\newunicodechar{≡}{\ensuremath{\equiv}}
\newunicodechar{⊂}{\ensuremath{\subset}}
\newunicodechar{⊥}{\ensuremath{\perp}}
\newunicodechar{∃}{\ensuremath{\exists}}
\newunicodechar{∀}{\ensuremath{\forall}}
\newunicodechar{∛}{\ensuremath{\sqrt[3]{\,}}}
\newunicodechar{∜}{\ensuremath{\sqrt[4]{\,}}}
\newunicodechar{⃗}{\ensuremath{{}^{\to}}}
\newunicodechar{ⁿ}{\ifmmode^{n}\else\textsuperscript{\ensuremath{n}}\fi}
\newunicodechar{ˣ}{\ifmmode^{x}\else\textsuperscript{\ensuremath{x}}\fi}
\newunicodechar{ᵇ}{\ifmmode^{b}\else\textsuperscript{\ensuremath{b}}\fi}
\newunicodechar{ʳ}{\ifmmode^{r}\else\textsuperscript{\ensuremath{r}}\fi}
\newunicodechar{ᵘ}{\ifmmode^{u}\else\textsuperscript{\ensuremath{u}}\fi}
\newunicodechar{ʸ}{\ifmmode^{y}\else\textsuperscript{\ensuremath{y}}\fi}
\newunicodechar{ᵃ}{\ifmmode^{a}\else\textsuperscript{\ensuremath{a}}\fi}
\newunicodechar{ᵉ}{\ifmmode^{e}\else\textsuperscript{\ensuremath{e}}\fi}
\newunicodechar{ᵏ}{\ifmmode^{k}\else\textsuperscript{\ensuremath{k}}\fi}
\newunicodechar{ᵖ}{\ifmmode^{p}\else\textsuperscript{\ensuremath{p}}\fi}
\newunicodechar{ᴺ}{\ifmmode^{N}\else\textsuperscript{\ensuremath{N}}\fi}
\newunicodechar{ᶜ}{\ifmmode^{c}\else\textsuperscript{\ensuremath{c}}\fi}
\newunicodechar{ʰ}{\ifmmode^{h}\else\textsuperscript{\ensuremath{h}}\fi}
\newunicodechar{ᵠ}{\ifmmode^{q}\else\textsuperscript{\ensuremath{q}}\fi}
\newunicodechar{ₙ}{\ifmmode_{n}\else\textsubscript{\ensuremath{n}}\fi}
\newunicodechar{ₐ}{\ifmmode_{a}\else\textsubscript{\ensuremath{a}}\fi}
\newunicodechar{ᵢ}{\ifmmode_{i}\else\textsubscript{\ensuremath{i}}\fi}
\newunicodechar{ᵣ}{\ifmmode_{r}\else\textsubscript{\ensuremath{r}}\fi}
\newunicodechar{ₖ}{\ifmmode_{k}\else\textsubscript{\ensuremath{k}}\fi}
\newunicodechar{ᵦ}{\ifmmode_{\beta}\else\textsubscript{\ensuremath{\beta}}\fi}
\newunicodechar{ₓ}{\ifmmode_{x}\else\textsubscript{\ensuremath{x}}\fi}
\newfontfamily\popdisplay{ArchivoBlack-Regular.ttf}[Path=../../../fonts/]
\newfontfamily\poplogo{SpaceGrotesk-Bold.ttf}[Path=../../../fonts/]
\newfontfamily\popsemi{PlusJakartaSans-SemiBold.ttf}[Path=../../../fonts/]
\newfontfamily\popxbold{PlusJakartaSans-ExtraBold.ttf}[Path=../../../fonts/]
\newfontfamily\popxboldls{PlusJakartaSans-ExtraBold.ttf}[Path=../../../fonts/,LetterSpace=3.0]

\setlist[enumerate]{leftmargin=1.7em,itemsep=5pt,topsep=4pt}
\setlist[itemize]{leftmargin=1.7em,itemsep=4pt,topsep=4pt,label={\color{poporange}\textbullet}}
\setlength{\parindent}{0pt}
\setlength{\parskip}{5pt}
\renewcommand{\arraystretch}{1.25}

% pgfplots : style Pop par défaut
\pgfplotsset{
  every axis/.append style={axis line style={popink,line width=1pt},tick style={popink},
    grid style={popline,line width=0.5pt},label style={font=\small},tick label style={font=\footnotesize}},
  popcurve/.style={poporange,line width=1.8pt,no markers,smooth},
}
% Même style disponible dans un \draw TikZ simple (hors pgfplots)
\tikzset{popcurve/.style={poporange,line width=1.8pt}}
\tikzset{popgrid/.style={popline,very thin}}
\colorlet{popcurve}{poporange} % le nom du style est parfois utilisé comme couleur

% ============================================================
% Pill / squiggle
% ============================================================
\newcommand{\poppill}[3]{%
  \tikz[baseline=-0.55ex]\node[draw=popink,line width=1.2pt,rounded corners=2.6pt,%
    fill=#2,inner xsep=6.5pt,inner ysep=3pt]{%
    \popxboldls\fontsize{7}{7}\selectfont\color{#3}\MakeUppercase{#1}};%
}
\newcommand{\popsquiggle}[1]{%
  \tikz[baseline=-0.4ex]{
    \draw[#1,line width=2.6pt,line cap=round]
      plot [smooth] coordinates {(0,0.09) (0.34,0.01) (0.68,0.21) (1.02,0.01) (1.36,0.10)};
  }%
}
% Mot-clé surligné (marqueur orange sous le texte)
\newcommand{\cle}[1]{{\popxbold\color{popdark}#1}}

% ============================================================
% En-tête / pied
% ============================================================
\pagestyle{fancy}
\fancyhf{}
\renewcommand{\headrulewidth}{0pt}
\renewcommand{\footrulewidth}{0pt}
\lhead{\raisebox{-0.15ex}{\poplogo\fontsize{13}{13}\selectfont\color{popink}OnBuch\color{poporange}+}}
\rhead{\poppill{\DOCMATIERE\ \textbullet\ \DOCNIVEAU\ \textbullet\ Leçon \DOCLECON}{white}{popink}}
\lfoot{\popxbold\fontsize{7.6}{7.6}\selectfont\color{popmuted}\MakeUppercase{Cours signé OnBuch+}\ \textbullet\ {\popsemi\DOCTITRE}}
\rfoot{\tikz[baseline=-0.6ex]\node[rounded corners=8.5pt,fill=popink,minimum height=17pt,minimum width=17pt,inner xsep=5pt,inner sep=0]{\popxbold\fontsize{8}{8}\selectfont\color{popcream}\thepage\,/\,\pageref*{LastPage}};}

% ============================================================
% Titres
% ============================================================
\newcommand{\chaptitre}[2]{%
  \begin{tcolorbox}[enhanced,colback=poporange,colframe=popink,boxrule=2.6pt,arc=11pt,
      drop shadow=popink,left=15pt,right=15pt,top=12pt,bottom=11pt,before skip=3pt,after skip=18pt]
    \poppill{#2}{popink}{popcream}\par\vspace{9pt}
    {\raggedright\hyphenpenalty=10000\exhyphenpenalty=10000\popdisplay\fontsize{22}{24}\selectfont\color{popcream}\MakeUppercase{#1}\par}
  \end{tcolorbox}
}
% Section numérotée : pastille orange avec le numéro + titre Archivo
\newcounter{popsec}
\newcommand{\coursec}[1]{%
  \stepcounter{popsec}\setcounter{popsub}{0}%
  \par\vspace{16pt}\noindent
  \begin{minipage}{\linewidth}
  \tikz[baseline=-0.7ex]\node[draw=popink,line width=1.6pt,fill=poporange,rounded corners=4pt,minimum size=22pt,inner sep=0]{\popdisplay\fontsize{12}{12}\selectfont\color{popcream}\thepopsec};%
  \hspace{8pt}{\popdisplay\fontsize{15}{17}\selectfont\color{popink}\MakeUppercase{#1}}\par
  \vspace{4pt}\noindent{\color{poporange}\rule{36pt}{3.6pt}}
  \end{minipage}\par\nopagebreak\vspace{8pt}
}
\newcounter{popsub}
\newcommand{\courssub}[1]{%
  \stepcounter{popsub}%
  \par\vspace{9pt}\noindent{\popxbold\fontsize{11.5}{13}\selectfont\color{popink}{\color{poporange}\thepopsec.\thepopsub}\hspace{6pt}#1}\par\nopagebreak\vspace{4pt}
}

% ============================================================
% Boîtes pédagogiques — même recette : fond blanc, bordure épaisse colorée,
% ombre dure encre, badge plein. Titre optionnel [#1].
% ============================================================
\tcbset{popbase/.style={breakable,enhanced,colback=white,boxrule=2.3pt,arc=9pt,
  shadow={3pt}{-3pt}{0pt}{popink},left=14pt,right=14pt,top=11pt,bottom=11pt,before skip=11pt,after skip=15pt,
  fontupper=\popsemi\fontsize{10.6}{14.5}\selectfont\color{popink},
  fontlower=\popsemi\fontsize{10.6}{14.5}\selectfont\color{popink}}}
% Coche dessinée (le glyphe ✓ n'existe pas dans les polices du document)
\newcommand{\popcheck}[1]{\tikz[baseline=-0.5ex]\draw[#1,line width=1.6pt,line cap=round,line join=round](-0.13,0.0)--(-0.04,-0.09)--(0.13,0.1);}
\providecommand{\coche}{\popcheck{popgreen}}
\providecommand{\resultat}[1]{\par\smallskip\noindent\fcolorbox{popgreen}{popgreenL}{\parbox{\dimexpr\linewidth-2\fboxsep-2\fboxrule\relax}{\textbf{Résultat.} #1}}\par\smallskip}
\newcommand{\popboxhead}[3]{\poppill{#1}{#2}{white}\ifx\relax#3\relax\else\hspace{6pt}{\popxbold\fontsize{10.6}{12}\selectfont\color{popink}#3}\fi\par\vspace{7pt}}

\newtcolorbox{aretenir}[1][]{popbase,colframe=popgreen,before upper={\popboxhead{À retenir}{popgreen}{#1}}}
% Texte en allemand (document de lecture, dialogue, extrait) : cadre
% d'étude. Un titre optionnel (source, auteur) s'affiche dans l'en-tête.
\newtcolorbox{textebox}[1][]{popbase,colframe=popblue,before upper={\popboxhead{Text}{popblue}{#1}\begin{otherlanguage}{ngerman}},after upper={\end{otherlanguage}}}
% Marques ✓ / ✗ (forme correcte / fautive) souvent utilisées par les modèles
\providecommand{\cross}{\textcolor{poppink}{\ensuremath{\times}}}
\providecommand{\xmark}{\cross}\providecommand{\crossmark}{\cross}\providecommand{\cmark}{\ensuremath{\checkmark}}
% Ligne de réponse / justification à compléter (inventée par les modèles)
\providecommand{\justification}[1]{\par\noindent\textit{Justification~:} \dotfill\par}
% \textipa (package tipa non chargé) : la police ne couvre pas l'API -> simple passage
\providecommand{\textipa}[1]{#1}
% Soulignement (ulem non chargé) : \uline, \ul
\providecommand{\uline}[1]{\underline{#1}}\providecommand{\ul}[1]{\underline{#1}}
% variantes ASCII d'ordinaux écrites en macro
\providecommand{\iere}{\textsuperscript{re}}\providecommand{\ieme}{\textsuperscript{e}}\providecommand{\ere}{\textsuperscript{re}}\providecommand{\eme}{\textsuperscript{e}}\providecommand{\er}{\textsuperscript{er}}\providecommand{\ie}{\textsuperscript{e}}
% Ordinaux français écrits en macro par les modèles : 1\ière, 3\ième
\newcommand{\ière}{\textsuperscript{re}}\newcommand{\ième}{\textsuperscript{e}}
% \exemple (inventée par les modèles) : étiquette « Exemple : »
\providecommand{\exemple}{\textbf{Exemple~:}\ }
% \square utilisable aussi hors mode math (cases à cocher)
\AtBeginDocument{\let\squareMath\square\renewcommand{\square}{\ensuremath{\squareMath}}}
% Mot ou phrase en allemand à l'intérieur d'un paragraphe français
\newcommand{\all}[1]{\ifmmode\text{\textit{#1}}\else\begingroup\selectlanguage{ngerman}\itshape #1\endgroup\fi}
\let\esp\all % alias toléré
\newtcolorbox{exemplebox}[1][]{popbase,colframe=poporange,before upper={\popboxhead{Exemple}{poporange}{#1}}}
\newtcolorbox{definition}[1][]{popbase,colframe=popblue,before upper={\popboxhead{Définition}{popblue}{#1}}}
\newtcolorbox{propriete}[1][]{popbase,colframe=poppurple,before upper={\popboxhead{Propriété}{poppurple}{#1}}}
\newtcolorbox{methode}[1][]{popbase,colframe=popgold,before upper={\popboxhead{Méthode}{popgold}{#1}}}
\newtcolorbox{attention}[1][]{popbase,colframe=poppink,before upper={\popboxhead{Attention}{poppink}{#1}}}
\newtcolorbox{experience}[1][]{popbase,colframe=popblue,colback=popblueL,before upper={\popboxhead{Expérience}{popblue}{#1}}}
% « Le savais-tu ? » : carte violette pleine (contexte, culture, Cameroun)
\newtcolorbox{savaistu}[1][]{popbase,colback=poppurple,colframe=popink,
  fontupper=\popsemi\fontsize{10.4}{14.2}\selectfont\color{white},
  before upper={\poppill{Le savais-tu\,?}{white}{poppurple}\ifx\relax#1\relax\else\hspace{6pt}{\popxbold\color{white}#1}\fi\par\vspace{7pt}}}
% Exercice résolu : énoncé en haut, \tcblower puis la solution sur fond crème
\newcounter{popexo}\newcounter{popexr}
\newtcolorbox{exoresolu}[1][]{popbase,colframe=poporange,colbacklower=popcream,
  segmentation style={popink,line width=1.2pt,dashed},
  before upper={\stepcounter{popexr}\popboxhead{Exercice résolu \thepopexr}{poporange}{#1}},
  before lower={\poppill{Solution}{popink}{popcream}\par\vspace{6pt}}}
% Exercice (non résolu en place) — la correction est dans la partie Corrigés
\newtcolorbox{exercice}[1][]{popbase,colframe=popink,
  before upper={\stepcounter{popexo}\popboxhead{Exercice \thepopexo}{popink}{#1}}}
% Corrigé
\newtcolorbox{corrige}[1][]{popbase,colframe=popgreen,colback=popgreenL,
  before upper={\popboxhead{Corrigé}{popgreen}{#1}}}
% Objectifs / prérequis / bilan
\newtcolorbox{objectifs}{popbase,colframe=popink,colback=poporangeL,
  before upper={\popboxhead{Objectifs de la leçon}{popink}{}}}
\newtcolorbox{prerequis}{popbase,colframe=popink,colback=popblueL,
  before upper={\popboxhead{Prérequis}{popblue}{}}}
\newtcolorbox{bilan}{popbase,colframe=popink,colback=white,boxrule=2.8pt,
  before upper={\popboxhead{Fiche bilan}{poporange}{L'essentiel à connaître pour l'examen}}}
% Figure encadrée avec légende
\newtcolorbox{popfigure}[1][]{enhanced,breakable=false,colback=white,colframe=popline,boxrule=1.4pt,arc=8pt,
  left=10pt,right=10pt,top=10pt,bottom=8pt,before skip=10pt,after skip=12pt,halign=center,
  lower separated=false,halign lower=center,
  fontlower=\popsemi\fontsize{9}{11.5}\selectfont\color{popmuted},
  #1}
\newcommand{\legende}[1]{\par\vspace{6pt}{\popsemi\fontsize{9}{11.5}\selectfont\color{popmuted}#1}}

% ============================================================
% Signature OnBuch+
% ============================================================
\newcommand{\poplogoinline}[1]{{\poplogo\fontsize{#1}{#1}\selectfont\color{popink}OnBuch\color{poporange}+}}
\newcommand{\signatureonbuch}{%
  \begin{tcolorbox}[enhanced,colback=popink,colframe=popink,boxrule=2.2pt,arc=10pt,
      drop shadow=poporange,left=14pt,right=14pt,top=10pt,bottom=10pt,before skip=0pt,after skip=0pt]
    \tikz[baseline=-0.7ex]\node[circle,fill=popgreen,draw=popcream,line width=1.3pt,minimum size=22pt,inner sep=0]{\popcheck{white}};%
    \hspace{9pt}\begin{minipage}[c]{0.62\linewidth}
      {\popxbold\fontsize{12}{13}\selectfont\color{popcream}Signé\ \textbullet\ {\poplogo OnBuch\color{poporange}+}}\par\vspace{2pt}
      {\raggedright\popsemi\fontsize{8}{10.5}\selectfont\color{popline}\MakeUppercase{Équipe pédagogique OnBuch+}\\\MakeUppercase{Conforme au programme officiel MINESEC}\par}
    \end{minipage}\hfill
    {\poplogo\fontsize{22}{22}\selectfont\color{popcream}OnBuch\color{poporange}+}
  \end{tcolorbox}%
}

% ============================================================
% Page de garde
% #1 titre · #2 matière · #3 niveau · #4 module · #5 n° leçon · #6 durée ·
% #7 description / objectifs (texte ou itemize)
% ============================================================
\newcommand{\pagegarde}[7]{%
  \thispagestyle{empty}
  \newgeometry{a4paper,margin=1.7cm,top=1.6cm,bottom=1.4cm}%
  \begin{tikzpicture}[remember picture,overlay]
    \fill[poporange!18] ([xshift=-3.2cm,yshift=-3.6cm]current page.north east) circle (4.6cm);
    \fill[poppurple!14] ([xshift=2.4cm,yshift=7.6cm]current page.south west) circle (3.8cm);
  \end{tikzpicture}%
  {\poplogo\fontsize{30}{30}\selectfont\color{popink}OnBuch\color{poporange}+}\hfill
  \raisebox{0.6ex}{\poppill{Cours signé \textbullet\ Premium}{popink}{popcream}}\par
  \vspace{1pt}\popsquiggle{poppurple}\par
  \vspace{8pt}
  \begin{tcolorbox}[enhanced,colback=poporange,colframe=popink,boxrule=2.8pt,arc=13pt,
      drop shadow=popink,left=18pt,right=18pt,top=12pt,bottom=13pt,after skip=10pt]
    \poppill{#2 \textbullet\ #3}{popink}{popcream}\hspace{5pt}\poppill{#4}{white}{popink}\par\vspace{8pt}
    {\popdisplay\fontsize{14}{14}\selectfont\color{popink}LEÇON #5}\par\vspace{4pt}
    {\raggedright\hyphenpenalty=10000\exhyphenpenalty=10000\popdisplay\fontsize{25}{27}\selectfont\color{popcream}\MakeUppercase{#1}\par}
    \vspace{8pt}
    {\popxbold\fontsize{9.5}{9.5}\selectfont\color{popink}\MakeUppercase{Durée conseillée : #6}}
  \end{tcolorbox}
  \begin{tcolorbox}[enhanced,colback=white,colframe=popink,boxrule=2.2pt,arc=10pt,
      drop shadow=popink,left=16pt,right=16pt,top=10pt,bottom=10pt,after skip=8pt]
    \poppill{Dans cette leçon}{poppurple}{white}\par\vspace{5pt}
    \popsemi\fontsize{9.3}{12.6}\selectfont\color{popink}#7
  \end{tcolorbox}
  \noindent\begin{minipage}[t]{0.487\linewidth}
    \begin{tcolorbox}[enhanced,colback=popgreen,colframe=popink,boxrule=2.2pt,arc=10pt,
        drop shadow=popink,left=11pt,right=11pt,top=10pt,bottom=10pt,height=3.1cm,valign=center]
      \tcbox[colback=white,colframe=popink,boxrule=1.3pt,arc=5pt,boxsep=0pt,nobeforeafter,
        left=3pt,right=3pt,top=3pt,bottom=3pt,valign=center]{\qrcode[height=1.9cm]{https://play.google.com/store/apps/details?id=com.luvvix.onbuch&hl=fr}}%
      \hspace{8pt}\begin{minipage}[c]{0.5\linewidth}
        {\popdisplay\fontsize{11}{12.5}\selectfont\color{white}\MakeUppercase{Télécharge OnBuch}}\par\vspace{4pt}
        {\raggedright\popsemi\fontsize{8.4}{11}\selectfont\color{white}Gratuit \textbullet\ Google Play\\Tuteur IA, annales, concours, résultats\par}
      \end{minipage}
    \end{tcolorbox}
  \end{minipage}\hfill
  \begin{minipage}[t]{0.487\linewidth}
    \begin{tcolorbox}[enhanced,colback=poppurple,colframe=popink,boxrule=2.2pt,arc=10pt,
        drop shadow=popink,left=11pt,right=11pt,top=10pt,bottom=10pt,height=3.1cm,valign=center]
      \tikz[baseline=-0.7ex]\node[circle,fill=white,draw=popink,line width=1.3pt,minimum size=1.6cm,inner sep=0]{\popdisplay\fontsize{24}{24}\selectfont\color{poppurple}+};%
      \hspace{8pt}\begin{minipage}[c]{0.6\linewidth}
        {\popdisplay\fontsize{11}{12.5}\selectfont\color{white}\MakeUppercase{Passe à OnBuch+}}\par\vspace{4pt}
        {\raggedright\popsemi\fontsize{8.4}{11}\selectfont\color{white}Tous les cours signés, Tuteur IA illimité, fiches PDF — onglet OnBuch+ de l'app\par}
      \end{minipage}
    \end{tcolorbox}
  \end{minipage}
  \vspace{8pt}
  \popcontenu
  \vfill
  \signatureonbuch
  \restoregeometry
  \newpage
}

% Rangée « Ce cours contient » (page de garde)
\newcommand{\poptile}[3]{% #1 couleur #2 gros texte #3 légende
  \begin{tcolorbox}[enhanced,colback=#1,colframe=popink,boxrule=2pt,arc=9pt,shadow={2.5pt}{-2.5pt}{0pt}{popink},
      left=6pt,right=6pt,top=9pt,bottom=9pt,halign=center,valign=center,height=1.75cm,width=\linewidth,nobeforeafter]
    {\popdisplay\fontsize{12.5}{13}\selectfont\color{white}\MakeUppercase{#2}}\par\vspace{3pt}
    {\centering\popsemi\fontsize{7.8}{9.5}\selectfont\color{white}#3\par}
  \end{tcolorbox}}
\newcommand{\popcontenu}{%
  \par\noindent{\popxboldls\fontsize{7.5}{7.5}\selectfont\color{popmuted}CE COURS SIGNÉ CONTIENT}\par\vspace{7pt}
  \noindent\begin{minipage}[t]{0.235\linewidth}\poptile{popblue}{Cours}{complet, illustré, pas à pas}\end{minipage}\hfill
  \begin{minipage}[t]{0.235\linewidth}\poptile{poporange}{Méthodes}{et exercices résolus}\end{minipage}\hfill
  \begin{minipage}[t]{0.235\linewidth}\poptile{poppink}{Exercices}{type examen + corrigés}\end{minipage}\hfill
  \begin{minipage}[t]{0.235\linewidth}\poptile{popgreen}{Fiche bilan}{l'essentiel à retenir}\end{minipage}\par}

% Sommaire
\newtcolorbox{sommaire}{enhanced,colback=white,colframe=popink,boxrule=2.2pt,arc=10pt,
  drop shadow=popink,left=18pt,right=18pt,top=13pt,bottom=13pt,before skip=6pt,after skip=20pt,
  before upper={\poppill{Sommaire}{poppurple}{white}\par\vspace{9pt}\popsemi\fontsize{10.6}{14.5}\selectfont\color{popink}}}

% Signature de fin de cours — poussée tout en bas de la dernière page
\newcommand{\findecours}{%
  \par\vfill\nopagebreak
  \begin{center}\popsquiggle{poporange}\end{center}\vspace{4pt}
  \signatureonbuch
}
\AtBeginDocument{\def\llbracket{\mathopen{[\mkern-3.5mu[}}\def\rrbracket{\mathclose{]\mkern-3.5mu]}}} % intervalles d'entiers (glyphe ⟦ absent de la police)
% Glyphes absents de Fira Math : redéfinis à partir de glyphes disponibles
\AtBeginDocument{%
  \def\setminus{\mathbin{\backslash}}\let\smallsetminus\setminus
  \def\vdots{\mathord{\vbox{\baselineskip3.2pt\lineskiplimit0pt\kern4pt\hbox{.}\hbox{.}\hbox{.}}}}%
  \def\ddots{\mathinner{\mkern1mu\raise6.5pt\hbox{.}\mkern2mu\raise3.6pt\hbox{.}\mkern2mu\raise0.7pt\hbox{.}\mkern1mu}}%
  \def\triangle{\mathord{\scalebox{0.9}{$\symup{\Delta}$}}}%
  \def\nabla{\mathord{\raisebox{\depth}{\scalebox{1}[-1]{$\symup{\Delta}$}}}}%
  \def\checkmark{\popcheck{popdark}}%
  \def\star{\mathbin{\ast}}\def\bigstar{\mathord{\ast}}%
  \def\diamond{\mathbin{\scalebox{0.6}{$\Diamond$}}}%
  \def\bigcup{\mathop{\mathchoice{\raisebox{-0.3em}{\scalebox{1.7}{$\cup$}}}{\scalebox{1.25}{$\cup$}}{\cup}{\cup}}\displaylimits}%
  \def\bigcap{\mathop{\mathchoice{\raisebox{-0.3em}{\scalebox{1.7}{$\cap$}}}{\scalebox{1.25}{$\cap$}}{\cap}{\cap}}\displaylimits}%
  \def\lesseqqgtr{\mathrel{\lessgtr}}\def\bigoplus{\mathop{\oplus}\displaylimits}%
}
\providecommand{\ssi}{si et seulement si} % abréviation fréquente
\AtBeginDocument{\providecommand{\wideparen}{\overparen}} % arc de cercle

\begin{document}

\pagegarde{Le problème des minorités}{Allemand}{Tle A}{Chapitre 1 : Vie familiale et sociale}{AL04}{2 h}{Cette leçon de type « texte » propose la lecture et le commentaire d'une lettre de lecteur sur la situation des minorités dans les pays germanophones. Elle permet d'acquérir le vocabulaire thématique (Minderheit, Diskriminierung, Integration...), de réviser la déclinaison de l'adjectif, la concession (obwohl, trotzdem) et le subjonctif II du souhait, avant une production guidée.
\vspace{4pt}
\begin{multicols}{2}\small
\begin{itemize}
\item À la fin de cette leçon, je sais comprendre globalement puis en détail un texte sur le problème des minorités.
\item À la fin de cette leçon, je sais employer correctement le vocabulaire thématique (die Minderheit, die Diskriminierung, die Rechte, die Sprache, das Vorurteil, die Integration) avec articles et pluriels.
\item À la fin de cette leçon, je sais décliner l'adjectif épithète après article défini, indéfini et sans article.
\item À la fin de cette leçon, je sais exprimer la concession avec obwohl (subordonnée) et trotzdem (connecteur) sans les confondre.
\item À la fin de cette leçon, je sais exprimer un souhait au subjonctif II (wäre, hätte, gäbe, könnte).
\item À la fin de cette leçon, je sais répondre à des questions de compréhension en phrases complètes.
\end{itemize}\end{multicols}}

\begin{sommaire}
\begin{multicols}{2}
\begin{enumerate}
\item Texte support : « Ein Brief aus Bautzen »
\item Compréhension du texte
\item Lexique thématique : Minderheiten und Integration
\item Grammaire 1 : la déclinaison de l'adjectif
\item Grammaire 2 : la concession et le subjonctif II du souhait
\item Commentaire modèle et production guidée
\item Activité d'intégration
\item Méthodes et exercices résolus
\item Exercices
\item Corrigés des exercices
\item Fiche bilan
\end{enumerate}
\end{multicols}
\end{sommaire}

\begin{objectifs}
\begin{itemize}
\item À la fin de cette leçon, je sais comprendre globalement puis en détail un texte sur le problème des minorités.
\item À la fin de cette leçon, je sais employer correctement le vocabulaire thématique (die Minderheit, die Diskriminierung, die Rechte, die Sprache, das Vorurteil, die Integration) avec articles et pluriels.
\item À la fin de cette leçon, je sais décliner l'adjectif épithète après article défini, indéfini et sans article.
\item À la fin de cette leçon, je sais exprimer la concession avec obwohl (subordonnée) et trotzdem (connecteur) sans les confondre.
\item À la fin de cette leçon, je sais exprimer un souhait au subjonctif II (wäre, hätte, gäbe, könnte).
\item À la fin de cette leçon, je sais répondre à des questions de compréhension en phrases complètes.
\item À la fin de cette leçon, je sais rédiger un court paragraphe argumenté sur l'intégration des minorités au Cameroun ou en Allemagne.
\item À la fin de cette leçon, je sais résumer un texte en quelques phrases avec mes propres mots.
\end{itemize}
\end{objectifs}

\begin{prerequis}
\begin{itemize}
\item La déclinaison des articles (der/die/das, ein/eine) et les quatre cas (classe de Première)
\item La place du verbe : verbe en 2e position dans la principale, verbe à la fin dans la subordonnée
\item Le présent et le prétérit des verbes forts et faibles
\item Le vocabulaire de la vie en société vu en Première (die Gesellschaft, zusammenleben)
\item Notion de base sur le subjonctif II (formes avec würde + infinitif)
\end{itemize}
\end{prerequis}

\newpage

\coursec{Texte support : « Ein Brief aus Bautzen »}

Au Cameroun, plus de 250 groupes ethniques cohabitent et le pays compte deux langues officielles : le français et l'anglais. Comment vit-on quand on appartient à une minorité linguistique ou culturelle ? En Allemagne vivent des Sorabes en Saxe et au Brandebourg, des Danois au Schleswig-Holstein, des Sinti et Roma dans tout le pays ; l'Autriche reconnaît les Croates du Burgenland et la Suisse protège ses quatre langues nationales. Un jeune Sorabe de Bautzen écrit une lettre au journal : il se sent parfois étranger dans son propre pays. Que dit-il ? Quels mots utilise-t-il pour décrire les préjugés et l'intégration ? C'est ce que nous allons découvrir.

\courssub{Présentation du texte et lecture globale}

\begin{definition}[Le genre : le courrier des lecteurs (\all{der Leserbrief})]
Un \cle{Leserbrief} est une lettre qu'un lecteur adresse à la rédaction d'un journal ou d'une revue pour réagir à un article, témoigner d'une expérience personnelle ou exprimer une opinion. Il est publié dans la rubrique \all{Leserbriefe} ou \all{Leserforum}.
\textbf{Caractéristiques :}
\begin{itemize}
    \item Ton personnel, subjectif, souvent émotionnel (\all{Ich-Perspektive}).
    \item Structure formelle : salutation (\all{Liebe Redaktion,}), corps du texte, formule de politesse finale (\all{Mit freundlichen Grüßen}), signature (prénom ou nom complet + lieu).
    \item Argumentation concise : constat $\rightarrow$ exemple $\rightarrow$ revendication / souhait.
\end{itemize}
\end{definition}

\begin{textebox}[Ein Brief aus Bautzen — Leserbrief dans la \all{Sächsische Zeitung}]
„Liebe Redaktion,

ich heiße Tomas und ich bin 17 Jahre alt. Ich komme aus Bautzen und ich bin Sorabe. Die sorbische Minderheit lebt seit Jahrhunderten in der Oberlausitz, obwohl viele Deutsche in Sachsen gar nicht wissen, dass es uns gibt. In meiner Schule lernen wir Sorbisch als Muttersprache, trotzdem sprechen die meisten jungen Leute untereinander nur Deutsch. Manche Menschen haben Vorurteile gegen uns: sie denken, wir seien keine richtigen Deutschen, weil wir eine andere Sprache und andere Traditionen haben. Das ist eine Form von Diskriminierung, die wehtut.

Unsere Rechte sind zwar im Grundgesetz und im Sächsischen Gesetz über die Rechte der Sorben geschützt, aber die Integration im Alltag ist nicht immer leicht. Manchmal fühle ich mich fremd in meinem eigenen Land. Ich wünschte, alle Menschen hätten mehr Respekt vor anderen Kulturen und Sprachen. Wenn es keine Vorurteile gäbe, wäre das Zusammenleben viel einfacher.

Mit freundlichen Grüßen
Tomas (17), Bautzen“
\end{textebox}

\begin{exemplebox}[Identification du texte]
\textbf{Source :} \all{Sächsische Zeitung} (quotidien régional de Saxe), rubrique \all{Leserforum}.\\
\textbf{Auteur :} Tomas, 17 ans, lycéen, \all{Sorabe} (appartient à la minorité sorabe).\\
\textbf{Destinataire :} La rédaction (\all{die Redaktion}) $\rightarrow$ public large (lecteurs du journal).\\
\textbf{Sujet :} La situation de la minorité sorabe en Saxe : méconnaissance, préjugés, droits légaux vs réalité vécue, souhait de respect.\\
\textbf{Problématique implicite :} Comment concilier identité minoritaire et citoyenneté allemande quand la société majoritaire ignore votre existence ?
\end{exemplebox}

\courssub{Lecture détaillée par paragraphe et champ lexical}

\begin{propriete}[Méthode : lecture détaillée en trois temps]
\begin{enumerate}
    \item \textbf{Découper} le texte en unités de sens (paragraphes logiques).
    \item \textbf{Annoter} : souligner les mots-clés du thème (minorité, droits, préjugés...), repérer les connecteurs logiques (\all{obwohl, trotzdem, aber, weil}).
    \item \textbf{Réformuler} en français chaque paragraphe en une phrase pour vérifier la compréhension.
\end{enumerate}
\end{propriete}

\begin{center}
\begin{tabularx}{\linewidth}{|X|X|X|}
\hline
\rowcolor{popblueL}
\textbf{Paragraphe (Allemagne)} & \textbf{Idée principale} & \textbf{Mots-clés du champ lexical} \\
\hline
\all{Ich heiße Tomas... gibt.} (Présentation + méconnaissance) & Identité + invisibilité de la minorité & \all{die Minderheit, seit Jahrhunderten, obwohl, wissen} \\
\hline
\all{In meiner Schule... nur Deutsch.} (École + pratique linguistique) & Droit à la langue maternelle vs réalité monolingue & \all{die Muttersprache, trotzdem, die meisten, untereinander} \\
\hline
\all{Manche Menschen... wehtut.} (Préjugés + discrimination) & Stéréotypes excluant + sentiment d'injustice & \all{das Vorurteil, die Diskriminierung, richtig, wehtun} \\
\hline
\all{Unsere Rechte... nicht immer leicht.} (Cadre légal + intégration difficile) & Protection juridique insuffisante dans le quotidien & \all{die Rechte, schützen, das Gesetz, die Integration, der Alltag} \\
\hline
\all{Manchmal fühle ich... einfacher.} (Sentiment + souhait) & Étranger chez soi + vœu d'une société inclusive & \all{fremd, das eigene Land, wünschen, der Respekt, das Zusammenleben} \\
\hline
\end{tabularx}
\end{center}

\begin{textebox}[Glossaire bilingue — Mots difficiles du texte (12 items)]
\begin{itemize}
    \item \all{die Minderheit, die Minderheiten} = la minorité
    \item \all{das Vorurteil, die Vorurteile} = le préjugé
    \item \all{die Diskriminierung, die Diskriminierungen} = la discrimination
    \item \all{das Recht, die Rechte} = le droit
    \item \all{die Integration} = l'intégration (pas de pluriel usuel)
    \item \all{die Muttersprache, die Muttersprachen} = la langue maternelle
    \item \all{schützen} (regelmäßig) = protéger
    \item \all{der Respekt} = le respect
    \item \all{das Zusammenleben} = la vie commune, la cohabitation
    \item \all{fremd} = étranger, étrange ; \all{sich fremd fühlen} = se sentir étranger
    \item \all{wehtun} (unregelmäßig: tut weh, hat wehgetan) = faire mal (physique ou moral)
    \item \all{die Oberlausitz} = la Haute-Lusace (région de Bautzen/Budyšin)
\end{itemize}
\end{textebox}

\courssub{Repérage des structures grammaticales du programme}

Le texte support est une mine d'exemples pour les trois points de grammaire/conjugaison de cette leçon. Observons-les \emph{in situ} avant d'en énoncer les règles (sections 4 et 5).

\begin{exemplebox}[1. Déclinaison de l'adjectif (repérage dans le texte)]
\begin{itemize}
    \item \all{die \textbf{sorbische} Minderheit} (féminin, nominatif/accusatif, article défini $\rightarrow$ déclinaison faible : \textbf{-e})
    \item \all{viele \textbf{Deutsche}} (pluriel, nominatif/accusatif, article indéfini/quantificatif $\rightarrow$ déclinaison forte : \textbf{-e})
    \item \all{die \textbf{sorbische} Sprache} (féminin, accusatif, article défini $\rightarrow$ faible : \textbf{-e})
    \item \all{die \textbf{meisten} \textbf{jungen} Leute} (pluriel, nominatif, article défini + adjectif $\rightarrow$ faible : \textbf{-en} / forte : \textbf{-e})
    \item \all{keine \textbf{richtigen} Deutschen} (pluriel, accusatif, article négatif $\rightarrow$ forte : \textbf{-en})
    \item \all{vor \textbf{anderen} Kulturen} (pluriel, datif après préposition \all{vor}, article indéfini/possessif absent $\rightarrow$ forte : \textbf{-n})
\end{itemize}
\emph{Règle détaillée en Section 4.}
\end{exemplebox}

\begin{exemplebox}[2. La concession : \all{obwohl} vs \all{trotzdem}]
\begin{itemize}
    \item \all{\textbf{Obwohl} viele Deutsche das nicht wissen,} (subordonnée, verbe à la fin) \all{lebt die Minderheit hier.} (principale, verbe en 2\textsuperscript{e} position).
    \item \all{In der Schule lernen wir Sorbisch,} \textbf{\all{trotzdem}} \all{sprechen die meisten nur Deutsch.} (adverbe de liaison, place libre, ici en tête de principale $\rightarrow$ inversion sujet/verbe).
\end{itemize}
\emph{Règle détaillée en Section 5.}
\end{exemplebox}

\begin{exemplebox}[3. Subjonctif II pour exprimer un souhait (irréel du présent)]
\begin{itemize}
    \item \all{Ich \textbf{wünschte}, alle Menschen \textbf{hätten} mehr Respekt.} (\all{wünschen} + \all{dass}-phrase au Subjonctif II : \all{hätten} = Subj. II de \all{haben}).
    \item \all{\textbf{Wenn} es keine Vorurteile \textbf{gäbe}, \textbf{wäre} das Zusammenleben einfacher.} (Phrase conditionnelle \all{wenn} + Subj. II \all{gäbe} / \all{wäre}).
\end{itemize}
\emph{Règle détaillée en Section 5.}
\end{exemplebox}

\begin{exoresolu}[Repérage et analyse grammaticale]
\textbf{Énoncé :} Relevez dans le texte \textbf{un autre adjectif décliné} (autre que ceux cités ci-dessus), \textbf{la phrase introduite par \all{obwohl}} et \textbf{la forme du Subjonctif II du verbe \all{sein}} dans la phrase conditionnelle. Justifiez le cas et le genre de l'adjectif choisi.

\textbf{Solution détaillée :}
\begin{enumerate}
    \item \textbf{Adjectif :} \all{das \textbf{eigene} Land} (dernier paragraphe : \all{in meinem eigenen Land}).
        \begin{itemize}
            \item Genre : neutre (\all{das Land}).
            \item Cas : datif (préposition \all{in} + lieu $\rightarrow$ datif).
            \item Article : possessif \all{meinem} (terminaison \textbf{-m} = datif neutre).
            \item Déclinaison : \textbf{faible} (article défini/possessif présent) $\rightarrow$ terminaison \textbf{-n} pour le datif neutre. Forme : \all{eigen\textbf{n}}.
        \end{itemize}
    \item \textbf{Phrase \all{obwohl} :} \all{Obwohl viele Deutsche das nicht wissen, lebt die sorbische Minderheit seit Jahrhunderten in der Oberlausitz.}
    \item \textbf{Subjonctif II de \all{sein} :} \all{wäre} (3\textsuperscript{e} personne du singulier). Phrase : \all{Wenn es keine Vorurteile gäbe, \textbf{wäre} das Zusammenleben viel einfacher.}
\end{enumerate}
\end{exoresolu}

\courssub{Méthode : Lire un texte au Bac (compréhension écrite)}

\begin{methode}[Réussir l'épreuve de compréhension écrite]
\begin{enumerate}
    \item \textbf{Lire le titre, la source et l'auteur.} Identifier le genre (\all{Leserbrief, Artikel, Kurzgeschichte}), le thème probable et le point de vue.
    \item \textbf{Première lecture globale (5 min).} Ne pas s'arrêter aux mots inconnus. Chercher : Qui parle ? À qui ? De quoi ? Quel ton ? (colère, espoir, ironie...).
    \item \textbf{Deuxième lecture active (10 min).} Souligner : \textcolor{popblue}{mots-clés du thème}, \textcolor{popgreen}{connecteurs logiques}, \textcolor{poporange}{structures grammaticales ciblées}. Noter en marge l'idée-force de chaque paragraphe.
    \item \textbf{Traiter les questions.} \textbf{Répondre en phrases complètes en allemand.} Reprendre le vocabulaire du texte et de la question. Citer le texte (numéro de ligne ou guillemets) pour justifier.
    \item \textbf{Relire ses réponses.} Vérifier : orthographe (majuscules noms, Umlaute), grammaire (cas, ordre des mots), cohérence avec le texte.
\end{enumerate}
\end{methode}

\courssub{Le savais-tu ? : Les Sorabes, une minorité slave au cœur de l'Allemagne}

\begin{savaistu}[Les Sorabes (\all{die Sorben} / \all{Serbja})]
\begin{itemize}
    \item Peuple slave occidental installé depuis le VI\textsuperscript{e} siècle en \all{Lausitz} (Lusace) : aujourd'hui Saxe (\all{Sachsen}) et Brandebourg (\all{Brandenburg}).
    \item Effectif : environ \textbf{60 000 personnes} (estimation), dont 20 000 à 30 000 locuteurs actifs du sorabe (haut-sorabe et bas-sorabe).
    \item Statut : minorité nationale reconnue (\all{anerkanntes Volk}) avec droits constitutionnels (écoles bilingues, panneaux bilingues, usage devant les tribunaux).
    \item \textbf{Bautzen = Budyšin} : centre culturel sorabe (théâtre, maison d'édition, gymnase sorabe). Panneaux de ville bilingues \all{Bautzen / Budyšin}.
    \item Fête nationale : \all{Vogelhochzeit} (25 janvier), \all{Zapust} (carnaval), costumes traditionnels (\all{Trachten}) portés fièrement.
    \item Parallèle camerounais : comme les Bamiléké, les Bassa ou les Peuls ont leur langue et culture au sein de la nation camerounaise, les Sorabes sont Allemands \textbf{et} Sorabes.
\end{itemize}
\end{savaistu}

\begin{popfigure}
\begin{tikzpicture}[
    mindmap,
    grow cyclic,
    every node/.style={concept, execute at begin node=\hskip0pt},
    concept color=popblue,
    level 1/.append style={level distance=4.5cm, sibling angle=72},
    level 2/.append style={level distance=3cm, sibling angle=45},
    Sorben/.style={concept color=poporange},
    Daenen/.style={concept color=popgreen},
    SintiRoma/.style={concept color=poppurple},
    Kroaten/.style={concept color=poppink},
    Schweiz/.style={concept color=popgold}
]
\node[concept] {die Minderheiten}
    child[Sorben] { node[concept] {Soraben} }
    child[Daenen] { node[concept] {Dänen} }
    child[SintiRoma] { node[concept] {Sinti und Roma} }
    child[Kroaten] { node[concept] {Kroaten} }
    child[Schweiz] { node[concept] {Sprachen der Schweiz} };
\end{tikzpicture}
\legende{Carte mentale : les minorités reconnues dans l'espace germanophone (Allemagne, Autriche, Suisse).}
\end{popfigure}

\begin{attention}[Pièges fréquents pour francophones — Vocabulaire \& Faux amis]
\begin{itemize}
    \item \all{die Minderheit} (féminin) $\neq$ \all{der Minderheit} (masculin n'existe pas). Attention : \all{Minderheit} prend \textbf{die} au singulier.
    \item \all{das Vorurteil} (neutre, composé : \all{vor} + \all{Urteil}) $\rightarrow$ pluriel \all{die Vorurteil\textbf{e}}. Ne pas confondre avec \all{das Urteil} (le jugement).
    \item \all{die Diskriminierung} (féminin, suffixe \all{-ierung}) $\rightarrow$ pluriel \all{-en}. Verbe : \all{diskriminieren} (faible). Ne pas dire \all{*diskriminieren gegen} mais \all{jemanden diskriminieren} (accusatif) ou \all{gegen jemanden diskriminieren}.
    \item \all{die Integration} (féminin, pas de pluriel usuel) $\neq$ \all{die Integrität} (l'intégrité, l'honnêteté). Verbe : \all{integrieren} $\rightarrow$ \all{sich integrieren} (réfléchi) ou \all{jemanden integrieren}.
    \item \all{fremd} (adjectif : étranger) vs \all{der Fremde} (nom : l'étranger, l'inconnu). \all{sich fremd fühlen} = se sentir dépaysé / étranger.
    \item Majuscules : \all{Sorbe, Deutscher, Deutsche, Sorbisch, Deutsch} (noms de langues, peuples, adjectifs de nationalité substantivés).
\end{itemize}
\end{attention}

\coursec{Compréhension du texte}

Après la lecture du texte support \emph{« Ein Brief aus Bautzen »} (section 1), nous passons maintenant à l'exploitation méthodique de ce document. La compréhension écrite suit une progression classique : du global vers le détaillé, puis la reformulation, l'évaluation par vrai/faux justifié, le résumé et enfin l'ouverture interculturelle. Chaque étape est l'occasion de réactiver le lexique du thème (Minderheiten, Integration, Vorurteile) et de s'entraîner aux exercices types du Baccalauréat.

\begin{textebox}{Ein Brief aus Bautzen — Lettre de Tomas (texte support)}
Lieber Lukas,

ich schreibe dir aus Bautzen, einer schönen Stadt in der Oberlausitz in Sachsen. Hier lebe ich mit meiner Familie, und ich möchte dir von meinem Alltag als Sorbe berichten.

Die Sorben sind eine slawische Minderheit in Deutschland. Wir haben unsere eigene Sprache, das Sorbische (oben- und niedersorbisch), unsere eigenen Traditionen, Feste und Trachten. Obwohl die sorbische Sprache durch Gesetze geschützt ist und es sorbische Schulen, Kindergärten und einen sorbischen Rundfunk gibt, erlebe ich im Alltag oft Vorurteile.

In der Schule zum Beispiel: Wenn ich mit meinen Freunden auf Sorbisch spreche, schauen manche Mitschüler komisch. Einmal hat ein Lehrer sogar gesagt: „Hier wird Deutsch gesprochen!“ Ich finde, das ist eine Form von Diskriminierung. Meine Sprache ist ein Teil meiner Identität. Warum soll ich sie verstecken?

Meine Eltern kämpfen für unsere Rechte. Mein Vater ist in der Domowina, dem Dachverband der sorbischen Vereine, aktiv. Er sagt: „Ohne Sprache stirbt ein Volk.“ Ich wünschte, mehr Deutsche würden verstehen, dass Vielfalt eine Bereicherung ist, keine Bedrohung.

Trotzdem bleibe ich optimistisch. Nächste Woche feiern wir das Vogelhochzeit-Fest (Ptači kwas). Das ist ein altes sorbisches Fest für die Kinder. Ich hoffe, du besuchst mich einmal, damit ich dir unsere Kultur zeigen kann.

Dein Freund
Tomas
\end{textebox}

\begin{definition}[Glossaire du texte — Vokabeln]
\begin{itemize}
\item \cle{die Oberlausitz} (f., kein Plural) — la Haute-Lusace (région de Saxe)
\item \cle{die Sorben} (Pl.) — les Sorabes (peuple slave minoritaire en Allemagne)
\item \cle{slawisch} — slave
\item \cle{die Minderheit, -en} — la minorité
\item \cle{das Sorbische} — le sorabe (langue)
\item \cle{obersorbisch / niedersorbisch} — haut-sorabe / bas-sorabe
\item \cle{die Tradition, -en} — la tradition
\item \cle{das Fest, -e} — la fête
\item \cle{die Tracht, -en} — le costume traditionnel
\item \cle{geschützt} — protégé
\item \cle{der Kindergarten, die Kindergärten} — la maternelle
\item \cle{der Rundfunk} — la radio / radiodiffusion
\item \cle{das Vorurteil, -e} — le préjugé
\item \cle{die Diskriminierung, -en} — la discrimination
\item \cle{die Identität, -en} — l'identité
\item \cle{der Dachverband, -e} — l'organisation faîtière
\item \cle{aktiv} — actif / engagé
\item \cle{sterben (stirbt, starb, ist gestorben)} — mourir / s'éteindre
\item \cle{die Vielfalt} — la diversité
\item \cle{die Bereicherung, -en} — l'enrichissement
\item \cle{die Bedrohung, -en} — la menace
\item \cle{optimistisch} — optimiste
\item \cle{die Vogelhochzeit} — la fête des oiseaux (fête sorabe pour enfants, 25 janvier)
\end{itemize}
\end{definition}

\courssub{Compréhension globale — Gesamtverständnis}

\begin{methode}[Comment répondre aux questions de compréhension (type Bac)]
\begin{enumerate}
\item \textbf{Lisez la question} et soulignez les mots-clés (Wer? Was? Wo? Warum? Wie?).
\item \textbf{Repérez l'information} dans le texte (paragraphe, ligne).
\item \textbf{Reprenez les mots de la question} pour construire votre réponse en allemand.
\item \textbf{Répondez par une phrase complète} (sujet + verbe conjugué + complément).
\item \textbf{Citez le texte si la consigne l'exige} : \emph{„Zitieren Sie!“} → copiez le segment exact entre guillemets allemands „…“.
\item \textbf{Ne traduisez pas mot à mot} du français : le correcteur attend une phrase allemande correcte, pas un calque.
\end{enumerate}
\end{methode}

\begin{exemplebox}[Questions de compréhension globale — modèles corrigés]
\textbf{1. Um was für einen Text handelt es sich? Wer ist der Autor?} \\
\emph{Antwort:} Es handelt sich um einen privaten Brief. Der Autor ist Tomas, ein sorbischer Schüler.

\textbf{2. An wen schreibt Tomas?} \\
\emph{Antwort:} Er schreibt an seinen Freund Lukas.

\textbf{3. Was ist das Ziel / der Zweck des Briefes?} \\
\emph{Antwort:} Tomas möchte Lukas über seinen Alltag als Sorbe informieren und ihn für die Situation der Minderheit sensibilisieren. Er möchte auch, dass Lukas ihn besucht.
\end{exemplebox}

\begin{attention}[Piège fréquent : calque du français]
Ne répondez pas : „\emph{Es ist ein Brief von Tomas an Lukas.}“ (phrase nominale, style télégraphique). \\
Le correcteur attend : „\emph{Es handelt sich um einen Brief, den Tomas an seinen Freund Lukas schreibt.}“ (phrase complète, subordonnée relative). \\
De même, évitez : „\emph{Der Brief ist über Sorben.}“ → Correct : „\emph{In dem Brief berichtet Tomas über sein Leben als Sorbe.}"
\end{attention}

\courssub{Compréhension détaillée — Detailverständnis}

\begin{exoresolu}[Exercice type Bac : questions détaillées]
\textbf{Énoncé :} Répondez en allemand par des phrases complètes. Citez le texte si nécessaire.

\textbf{1. Wo wohnen die Sorben? (Où vivent les Sorabes ?)} \\
\textbf{2. Welches Problem hat Tomas in der Schule? (Quel problème Tomas rencontre-t-il à l'école ?)} \\
\textbf{3. Nennen Sie zwei Vorurteile, die Tomas beschreibt. (Citez deux préjugés que Tomas décrit.)} \\
\textbf{4. Was macht Tomas' Vater für die sorbischen Rechte? (Que fait le père de Tomas pour les droits sorabes ?)} \\
\textbf{5. Was wünscht sich Tomas? (Que souhaite Tomas ?)}

\tcblower
\textbf{Solution détaillée :}

\textbf{1.} Die Sorben wohnen in der Oberlausitz in Sachsen. (Z. 3–4: „Hier lebe ich mit meiner Familie… in der Oberlausitz in Sachsen.“)

\textbf{2.} In der Schule schauen Mitschüler komisch, wenn er Sorbisch spricht, und ein Lehrer hat gesagt: „Hier wird Deutsch gesprochen!“ (Z. 10–12)

\textbf{3.} Erstens: Die Mitschüler reagieren ablehnend („schauen komisch“), wenn sie Sorbisch hören. Zweitens: Ein Lehrer verbietet implizit die sorbische Sprache im Unterricht („Hier wird Deutsch gesprochen!“). Beide Reaktionen zeigen Vorurteile gegen die Minderheitensprache.

\textbf{4.} Er ist in der Domowina, dem Dachverband der sorbischen Vereine, aktiv und kämpft für die Rechte der Sorben. (Z. 14–15)

\textbf{5.} Er wünscht sich, dass mehr Deutsche \textbf{verstünden} (oder: \textbf{würden verstehen}), dass Vielfalt eine Bereicherung ist, und er hofft, dass Lukas ihn besucht, um die sorbische Kultur kennenzulernen. (Z. 17–19)
\end{exoresolu}

\courssub{Reformulation — Erklären mit eigenen Worten}

\begin{propriete}[Consigne type : „Erklären Sie … mit eigenen Worten“]
Cette consigne demande de \textbf{reformuler} un passage du texte \textbf{sans le copier}, en utilisant votre propre vocabulaire allemand. Montrez que vous avez compris le sens profond.
\end{propriete}

\begin{exemplebox}[Reformulation modèle]
\textbf{Consigne :} Erklären Sie den Satz: „Das ist eine Form von Diskriminierung“ mit eigenen Worten.

\textbf{Antwort:} Wenn Tomas in der Schule Sorbisch spricht und dafür negativ beurteilt oder sogar zurechtgewiesen wird, obwohl seine Sprache gesetzlich geschützt ist, wird er wegen seiner Herkunft und Sprache benachteiligt. Das nennt man Diskriminierung: die Ungleichbehandlung einer Person aufgrund ihrer Zugehörigkeit zu einer Minderheit.
\end{exemplebox}

\begin{exoresolu}[Entraînement reformulation]
\textbf{Énoncé :} Erklären Sie mit eigenen Worten: „Ohne Sprache stirbt ein Volk.“

\tcblower
\textbf{Solution :} Die Sprache ist der Träger der Kultur, der Geschichte und der Identität eines Volkes. Wenn eine Minderheitssprache nicht mehr gesprochen und weitergegeben wird, gehen das Wissen, die Traditionen und das Zusammengehörigkeitsgefühl verloren. Das Volk als kulturelle Einheit hört auf zu existieren, auch wenn die Menschen biologisch weiterleben.
\end{exoresolu}

\courssub{Vrai ou faux justifié — Richtig oder falsch? Begründen Sie mit einem Zitat!}

\begin{methode}[Méthode Vrai/Faux justifié (entraînement Bac)]
\begin{enumerate}
\item Lisez l'affirmation attentivement.
\item Cherchez dans le texte l'information qui la confirme ou la contredit.
\item Écrivez \textbf{Richtig} ou \textbf{Falsch}.
\item \textbf{Citez le passage exact} du texte (avec guillemets „…“) qui justifie votre choix.
\item Si c'est faux, donnez la correction brève.
\end{enumerate}
\end{methode}

\begin{exoresolu}[Série Vrai/Faux justifié]
\textbf{Énoncé :} Entscheiden Sie: Richtig oder Falsch? Begründen Sie mit einem Zitat aus dem Text.

1. Die sorbische Sprache ist in Deutschland nicht geschützt. \\
2. Tomas spricht in der Schule nie Sorbisch. \\
3. Die Domowina ist ein sorbischer Dachverband. \\
4. Tomas ist pessimistisch über die Zukunft der Sorben. \\
5. Das Vogelhochzeit-Fest ist ein Fest für Erwachsene.

\tcblower
\textbf{Solution :}

1. \textbf{Falsch.} Zitat: „Obwohl die sorbische Sprache durch Gesetze geschützt ist…“ (Z. 7–8)

2. \textbf{Falsch.} Zitat: „Wenn ich mit meinen Freunden auf Sorbisch spreche…“ (Z. 10) — il parle sorabe avec ses amis.

3. \textbf{Richtig.} Zitat: „Mein Vater ist in der Domowina, dem Dachverband der sorbischen Vereine, aktiv.“ (Z. 14–15)

4. \textbf{Falsch.} Zitat: „Trotzdem bleibe ich optimistisch.“ (Z. 17)

5. \textbf{Falsch.} Zitat: „Das ist ein altes sorbisches Fest für die Kinder.“ (Z. 18) — c'est une fête pour les enfants, pas pour les adultes.
\end{exoresolu}

\courssub{Résumé guidé — Zusammenfassung schreiben}

\begin{methode}[Comment rédiger un résumé en 3–4 phrases (tronc commun)]
\begin{enumerate}
\item \textbf{Identifiez l'idée principale} : de quoi parle le texte en une phrase ? (Sujet + thèse).
\item \textbf{Supprimez les exemples et détails} (noms propres, anecdotes, chiffres) — gardez l'essentiel.
\item \textbf{Reformulez} avec vos propres mots (pas de copier-coller).
\item \textbf{Reliez les idées} avec des connecteurs logiques (\all{zuerst, dann, außerdem, schließlich, weil, obwohl, deshalb}).
\item \textbf{Vérifiez} : 3–4 phrases, temps cohérent (Präsens ou Präteritum), allemand correct.
\end{enumerate}
\end{methode}

\begin{popfigure}
\begin{tikzpicture}[
  node distance=1.2cm and 1.5cm,
  box/.style={rectangle, rounded corners, draw=popblue, thick, fill=popblueL, text width=3.2cm, align=center, font=\small},
  arrow/.style={-Stealth, thick, popdark}
]
\node[box, align=center] (id){Identifier l'idée\\principale};
\node[box, right=of id, align=center] (sup){Supprimer les\\exemples/détails};
\node[box, right=of sup, align=center] (ref){Reformuler\\avec ses mots};
\node[box, right=of ref, align=center] (lien){Relier avec\\des connecteurs};
\draw[arrow] (id) -- (sup);
\draw[arrow] (sup) -- (ref);
\draw[arrow] (ref) -- (lien);
\end{tikzpicture}
\legende{Organigramme : les 4 étapes du résumé}
\end{popfigure}

\begin{exemplebox}[Résumé modèle (3 phrases)]
Tomas, ein sorbischer Schüler aus Bautzen, beschreibt die Lage seiner Minderheit in Sachsen. Obwohl die sorbische Sprache und Kultur gesetzlich geschützt sind, erlebt er im Schulalltag Vorurteile und Diskriminierung. Er wünscht sich, dass die deutsche Gesellschaft Vielfalt als Bereicherung anerkennt, und lädt seinen Freund ein, das Vogelhochzeit-Fest kennenzulernen.
\end{exemplebox}

\begin{exercice}[À vous de jouer : résumé personnel]
Rédigez votre propre résumé du texte en 3–4 phrases en allemand. Utilisez les connecteurs : \all{zuerst, außerdem, aber, deshalb, hofft}.
\end{exercice}

\courssub{Discussion orale : Le Cameroun et les minorités — Mündliche Diskussion}

\begin{experience}[Débat guidé : Minorités au Cameroun et en Allemagne]
\textbf{Consigne :} En binôme, préparez une courte intervention orale (1–2 min) en allemand sur le thème : \emph{„Gibt es vergleichbare Situationen wie bei den Sorben in Kamerun?“} Utilisez le vocabulaire de la leçon et les structures de concession (\all{obwohl, trotzdem}) et de souhait (\all{ich wünschte, …}).

\textbf{Pistes de réflexion (mots-clés allemands) :}
\begin{itemize}
\item \cle{die Minderheiten in Kamerun} — minorités au Cameroun (ex. : Pygmées/Baka, Mbororo, Montagnards)
\item \cle{die Sprachenvielfalt} — diversité linguistique (plus de 250 langues)
\item \cle{die Amtssprachen} — langues officielles (français, anglais)
\item \cle{die Diskriminierung / Benachteiligung} — discrimination / désavantage
\item \cle{die Integration / Assimilation} — intégration / assimilation
\item \cle{die Rechte der indigenen Völker} — droits des peuples autochtones
\end{itemize}

\textbf{Modèle de phrase de départ :}
\all{„In Kamerun gibt es auch Minderheiten, zum Beispiel die Baka. Obwohl sie Ureinwohner sind, werden sie oft benachteiligt. Ich wünschte, ihre Rechte würden besser respektiert.“}
\end{experience}

\begin{savaistu}[Le saviez-vous ? — Les Sorabes et le Cameroun]
Les Sorabes (\all{die Sorben}) sont une minorité slave reconnue en Allemagne (environ 60 000 personnes) vivant en Lusace (\all{die Lausitz}), à cheval sur la Saxe et le Brandebourg. Ils bénéficient d'écoles bilingues, de médias en sorabe et d'une autonomie culturelle via la \all{Domowina}. C'est un modèle souvent cité pour la protection des minorités en Europe.

Au Cameroun, la diversité linguistique est immense (plus de 250 langues), mais les langues officielles restent le français et l'anglais (héritage colonial). Certaines communautés (Baka, Mbororo) militent pour la reconnaissance de leurs droits linguistiques et territoriaux, dans le cadre de la Déclaration ONU sur les droits des peuples autochtones (2007).
\end{savaistu}

\begin{aretenir}[Check-list compréhension texte — Le problème des minorités]
\begin{itemize}
\item \textbf{Type de texte :} identifier (Brief, Artikel, Interview…) + auteur + destinataire + but.
\item \textbf{Global :} idée principale en 1 phrase.
\item \textbf{Détaillé :} réponses complètes en allemand, citations exactes („Zitat“).
\item \textbf{Reformulation :} expliquer \emph{mit eigenen Worten}, sans copier.
\item \textbf{Vrai/Faux :} décision + justification par citation.
\item \textbf{Résumé :} 3–4 phrases, connecteurs, reformulation, temps cohérent.
\item \textbf{Oral :} transfert vers contexte camerounais, vocabulaire thématique, structures ciblées (concession, souhait).
\end{itemize}
\end{aretenir}

\coursec{Lexique thématique : Minderheiten und Integration}

Après avoir analysé le texte \emph{Ein Brief aus Bautzen} et répondu aux questions de compréhension, nous allons maintenant construire systématiquement le vocabulaire nécessaire pour parler des minorités, de leurs droits, des problèmes qu'elles rencontrent et des solutions pour une coexistence harmonieuse. Ce lexique est la base de toute expression écrite ou orale sur ce thème au Baccalauréat.

\courssub{Champ 1 — Les personnes et les groupes : die Menschen und Gruppen}

Observons d'abord comment le texte introduit les acteurs : \all{„Die Sorben sind eine anerkannte Minderheit in Deutschland.“} (Les Sorabes sont une minorité reconnue en Allemagne.) / \all{„Die Mehrheit der Bevölkerung spricht Deutsch.“} (La majorité de la population parle allemand.)

\begin{definition}[Les acteurs sociaux]
Les noms suivants désignent les groupes humains au cœur du thème. \textbf{Mémorisez impérativement l'article, le genre et le pluriel.}
\end{definition}

\begin{center}
\begin{tabularx}{\linewidth}{|X|X|X|}
\hline
\rowcolor{popblueL}
\textbf{Deutsch (Singular / Plural)} & \textbf{Français} & \textbf{Remarque / Exemple} \\
\hline
\cle{die Minderheit, -en} & la minorité & \all{Eine nationale Minderheit.} (Une minorité nationale.) \\
\cle{die Mehrheit, -en} & la majorité & \all{Die absolute Mehrheit.} (La majorité absolue.) Antonyme de \cle{Minderheit}. \\
\cle{der Einwanderer, -} & l'immigré (homme) & Pluriel inchangé : \all{die Einwanderer}. Feminin : \cle{die Einwanderin, -nen}. \\
\cle{die Bevölkerung, -en} & la population & Toujours au singulier pour un groupe donné (\cle{die Bevölkerungen} = plusieurs peuples). \all{Die Bevölkerung wächst.} \\
\cle{der Ausländer, -} & l'étranger (homme) & Pluriel : \all{die Ausländer}. \textbf{Attention : connotation souvent administrative/négative.} \\
\cle{der Migrant, -en} & le migrant & Terme plus neutre, descriptif du mouvement. Feminin : \cle{die Migrantin, -nen}. \\
\cle{der Flüchtling, -e} & le réfugié & Personne fuyant la guerre/persécution. Feminin : \cle{die Flüchtling, -e} (inchangé). \\
\cle{der Staatsbürger, -} / \cle{die Staatsbürgerin, -nen} & le citoyen / la citoyenne & Droit de vote, passeport. \all{Deutsche Staatsbürger.} \\
\hline
\end{tabularx}
\end{center}

\begin{attention}[Ausländer vs. Menschen mit Migrationshintergrund]
Dans un texte formel, journalistique ou administratif moderne, on évite \cle{der Ausländer} (litt. « celui qui vient de l'extérieur », perçu comme excluant). On préfère \textbf{\cle{Menschen mit Migrationshintergrund}} (personnes ayant un passé migratoire) ou \cle{der Migrant / die Migrantin}. Le terme \cle{Ausländer} reste courant dans le langage juridique (\all{Ausländergesetz}).
\end{attention}

\courssub{Champ 2 — Les droits et la protection : Rechte und Schutz}

Le texte mentionne : \all{„Ihre Rechte sind durch das Grundgesetz geschützt.“} (Leurs droits sont protégés par la Loi fondamentale.) / \all{„Diskriminierung ist verboten.“} (La discrimination est interdite.)

\begin{propriete}[Noms en \texorpdfstring{-ung}{-ung} : toujours féminins, pluriel en -en]
Tous les noms dérivés de verbes par le suffixe \cle{-ung} sont de genre \textbf{féminin} (\cle{die}) et forment leur pluriel en \textbf{-en}.
\all{diskriminieren} $\rightarrow$ \cle{die Diskriminierung, -en} \\
\all{integrieren} $\rightarrow$ \cle{die Integration, -en} \\
\all{bevorzugen} $\rightarrow$ \cle{die Bevorzugung, -en} (préférence / favoritisme) \\
\all{benachteiligen} $\rightarrow$ \cle{die Benachteiligung, -en} (désavantage / discrimination indirecte)
\end{propriete}

\begin{center}
\begin{tabularx}{\linewidth}{|X|X|X|}
\hline
\rowcolor{popgreenL}
\textbf{Deutsch} & \textbf{Français} & \textbf{Collocations / Exemples} \\
\hline
\cle{das Recht, -e} & le droit & \cle{Menschenrechte} (Pl., droits de l'homme), \cle{Grundrechte} (droits fondamentaux), \cle{Wahlrecht} (droit de vote). \\
\cle{die Menschenrechte (Pl.)} & les droits de l'homme & \all{Die Menschenrechte sind universell.} \\
\cle{schützen} (schützte, hat geschützt) & protéger & \all{Das Gesetz schützt die Minderheiten.} \\
\cle{verbieten} (verbot, verboten) & interdire & Verbe à voyelle changeante (e $\rightarrow$ a/o). \all{Das Gesetz verbietet Rassismus.} \\
\cle{erlauben} & permettre, autoriser & \all{Das Grundgesetz erlaubt freie Meinungsäußerung.} \\
\cle{die Gleichberechtigung} & l'égalité des droits & \all{Gleichberechtigung von Mann und Frau.} (Égalité H/F.) \\
\cle{die Gleichstellung} & l'égalité (statutaire) & Souvent synonyme, plus administratif. \\
\cle{der Schutz} & la protection & \cle{Minderheitenschutz} (protection des minorités). \\
\hline
\end{tabularx}
\end{center}

\begin{exemplebox}[Droits et verbes]
\all{Die Verfassung garantiert die Rechte der Minderheiten.} \\
« La constitution garantit les droits des minorités. » \\
\all{Man darf niemanden wegen seiner Herkunft benachteiligen.} \\
« On ne doit désavantager personne à cause de son origine. » (\cle{wegen} + génitif).
\end{exemplebox}

\courssub{Champ 3 — Les problèmes : die Probleme}

Le texte évoque : \all{„Vorurteile führen oft zu Diskriminierung.“} (Les préjugés mènent souvent à la discrimination.) / \all{„Niemand soll ausgegrenzt werden.“} (Personne ne doit être mis à l'écart.)

\begin{center}
\begin{tabularx}{\linewidth}{|X|X|X|}
\hline
\rowcolor{poporangeL}
\textbf{Deutsch} & \textbf{Français} & \textbf{Construction / Exemple} \\
\hline
\cle{die Diskriminierung, -en} & la discrimination & \all{opfer von Diskriminierung sein} (être victime de...). \\
\cle{das Vorurteil, -e} & le préjugé & Composé : \cle{vor} (avant) + \cle{das Urteil} (le jugement). Voir \textbf{Définition} ci-dessous. \\
\cle{der Rassismus} & le racisme & Pas de pluriel usuel. \all{gegen Rassismus kämpfen.} \\
\cle{ausgrenzen} & exclure, marginaliser & Verbe à particule séparable. \all{Man grenzt sie aus.} (On les exclut.) \\
\cle{benachteiligen} & désavantager, discriminer & \all{Benachteiligung am Arbeitsmarkt.} (Discrimination à l'embauche.) \\
\cle{der Hass} & la haine & \all{Hassrede} (discours de haine / \cle{Hatespeech}). \\
\cle{die Ausgrenzung, -en} & l'exclusion (sociale) & Nominalisation de \cle{ausgrenzen} (féminin -ung). \\
\cle{die Intoleranz} & l'intolérance & Antonyme : \cle{die Toleranz}. \\
\hline
\end{tabularx}
\end{center}

\begin{definition}[\texorpdfstring{\cle{das Vorurteil, -e}}{das Vorurteil}]
\cle{vor} (devant, avant) + \cle{das Urteil} (le jugement) = un jugement porté \textbf{avant} de connaître la personne ou les faits. C'est une opinion préconçue, souvent négative, basée sur des stéréotypes (\cle{das Klischee, -s / -s}).
\all{Er hat viele Vorurteile gegen Ausländer.} (Il a beaucoup de préjugés contre les étrangers.)
\end{definition}

\begin{aretenir}[Famille de mots : \texorpdfstring{urteilen}{urteilen}]
\begin{itemize}
    \item \all{urteilen} (juger) $\rightarrow$ \cle{das Urteil, -e} (le jugement, le verdict).
    \item \cle{vor} + \cle{Urteil} $\rightarrow$ \cle{das Vorurteil, -e} (le préjugé).
    \item \all{vorurteilsfrei} (adj., sans préjugé, objectif).
    \item \all{urteilen über jmdn/etw} (juger qqn/qqch).
\end{itemize}
\end{aretenir}

\courssub{Champ 4 — Les solutions : die Lösungen}

Le texte propose : \all{„Integration gelingt durch gegenseitigen Respekt.“} (La réussite de l'intégration passe par le respect mutuel.) / \all{„Wir müssen toleranter werden.“} (Nous devons devenir plus tolérants.)

\begin{center}
\begin{tabularx}{\linewidth}{|X|X|X|}
\hline
\rowcolor{poppurpleL}
\textbf{Deutsch} & \textbf{Français} & \textbf{Grammaire / Exemple} \\
\hline
\cle{die Integration} & l'intégration & Nom en \cle{-ung} (fém., pl. \cle{-en}). \all{Die Integration gelingt.} (L'intégration réussit.) \\
\cle{integrieren} / \cle{sich integrieren} & intégrer / s'intégrer & Verbe régulier. \textbf{Réfléchi} pour « s'intégrer » : \all{Er integriert sich gut.} \\
\cle{akzeptieren} & accepter & \all{Andere Meinungen akzeptieren.} \\
\cle{respektieren} & respecter & \all{Die Kultur anderer respektieren.} \\
\cle{die Toleranz} & la tolérance & \all{Toleranz zeigen.} (Faire preuve de tolérance.) Adj. : \cle{tolerant}. \\
\cle{das Zusammenleben} & la coexistence, le vivre-ensemble & Nom composé : \cle{zusammen} + \cle{leben}. \all{friedliches Zusammenleben.} \\
\cle{die Teilhabe} & la participation (sociale) & \all{Teilhabe am gesellschaftlichen Leben.} \\
\cle{fördern} & favoriser, encourager & \all{Integration fördern.} (Favoriser l'intégration.) \\
\hline
\end{tabularx}
\end{center}

\begin{exemplebox}[Verbes clés : transitif vs réfléchi]
\all{Die Schule integriert die Kinder.} (L'école intègre les enfants. \textbf{Transitif}, accusatif.) \\
\all{Die Kinder integrieren sich schnell.} (Les enfants s'intègrent vite. \textbf{Réfléchi}, accusatif \cle{sich}.)
\end{exemplebox}

\courssub{Champ 5 — Langue et culture : Sprache und Kultur}

Le texte souligne : \all{„Die sorbische Sprache ist ein Schatz.“} (La langue sorabe est un trésor.) / \all{„Zweisprachigkeit ist ein Vorteil.“} (Le bilinguisme est un atout.)

\begin{center}
\begin{tabularx}{\linewidth}{|X|X|X|}
\hline
\rowcolor{poppinkL}
\textbf{Deutsch} & \textbf{Français} & \textbf{Remarque} \\
\hline
\cle{die Sprache, -n} & la langue & \cle{Muttersprache} (langue maternelle), \cle{Fremdsprache} (langue étrangère), \cle{die Amtssprache} (langue officielle). \\
\cle{die Muttersprache, -n} & la langue maternelle & Lit. « langue de la mère ». Synonyme : \cle{Erstsprache}. \\
\cle{die Kultur, -en} & la culture & \cle{kulturelle Vielfalt} (diversité culturelle). Adj. : \cle{kulturell}. \\
\cle{die Tradition, -en} & la tradition & \cle{Traditionen pflegen} (entretenir des traditions). \\
\cle{zweisprachig} & bilingue & Adj. invariable en prédicat : \all{Er ist zweisprachig.} (Il est bilingue.) \\
\cle{die Zweisprachigkeit} & le bilinguisme & Nom fém. en \cle{-keit}. \\
\cle{die Identität, -en} & l'identité & \cle{kulturelle Identität}. \\
\cle{der Dialekt, -e} & le dialecte & \all{Regionale Dialekte.} \\
\hline
\end{tabularx}
\end{center}

\begin{savaistu}[La Suisse : un modèle d'intégration linguistique]
La Suisse (\cle{die Schweiz}) possède \textbf{quatre langues nationales} : l'allemand (\cle{Deutsch}, parlé par ~63\%), le français (\cle{Französisch}, ~23\%), l'italien (\cle{Italienisch}, ~8\%) et le romanche (\cle{Rätoromanisch}, <1\%). Le romanche est langue officielle au niveau fédéral pour la communication avec les romanches. Ce quadrilinguisme institutionnel est souvent cité en exemple (au Bac !) de \cle{Zusammenleben} pacifique et de \cle{Minderheitenschutz} par la loi (\cle{Sprachengesetz}).
\end{savaistu}

\courssub{Synthèse visuelle : Carte mentale « Integration »}

\begin{popfigure}
\begin{tikzpicture}[
    mindmap,
    grow cyclic,
    every node/.style=concept,
    concept color=popblue,
    level 1/.append style={level distance=4.5cm, sibling angle=90},
    level 2/.append style={level distance=3cm, sibling angle=45},
    level 3/.append style={level distance=2.5cm},
    text=white,
    font=\small\bfseries
]
\node[concept] {Integration}
    child[concept color=popgreen] { node {Rechte}
        child { node {Menschenrechte} }
        child { node {Gleichberechtigung} }
        child { node {Minderheitenschutz} }
        child { node {Teilhabe} }
    }
    child[concept color=poporange] { node {Probleme}
        child { node {Diskriminierung} }
        child { node {Rassismus} }
        child { node {Vorurteile} }
        child { node {Ausgrenzung} }
        child { node {Hass} }
    }
    child[concept color=poppurple] { node {Lösungen}
        child { node {Toleranz} }
        child { node {Respekt} }
        child { node {Akzeptanz} }
        child { node {Förderung} }
        child { node {Zusammenleben} }
    }
    child[concept color=poppink] { node {Sprache und Kultur}
        child { node {Muttersprache} }
        child { node {Zweisprachigkeit} }
        child { node {Kulturelle Vielfalt} }
        child { node {Identität} }
        child { node {Traditionen} }
    };
\end{tikzpicture}
\legende{Carte mentale lexicale : les quatre piliers du thème « Minorités et Intégration ».}
\end{popfigure}

\courssub{Texte d'approfondissement : « Leben in Vielfalt »}

Le texte suivant (niveau B1/B2) réemploie le vocabulaire des 5 champs. Lisez-le, repérez les mots appris et répondez aux questions.

\begin{textebox}[Leben in Vielfalt — Ein Plädoyer für Toleranz]
Deutschland ist ein Einwanderungsland. Seit Jahrzehnten kommen Menschen aus aller Welt hierher: als Arbeitsmigranten in den 1960er Jahren, als Flüchtlinge vor Krieg und Verfolgung oder als Fachkräfte für die Wirtschaft. Heute hat fast jeder vierte Einwohner einen Migrationshintergrund. Diese Vielfalt bereichert das Land kulturell und ökonomisch. Doch die Realität ist nicht immer einfach. Viele Menschen mit Migrationshintergrund erleben Diskriminierung auf dem Wohnungsmarkt, bei der Jobsuche oder im Alltag. Vorurteile und rassistische Klischees führen zu Ausgrenzung. Das Grundgesetz verbietet Benachteiligung wegen der Herkunft (Artikel 3). Dennoch reicht Rechtsschutz allein nicht. Integration ist ein Prozess, der von beiden Seiten Anstrengung verlangt: Die Gesellschaft muss Vielfalt akzeptieren und Respekt zeigen. Zugleich sollen Zuwanderer die deutsche Sprache lernen und die Gesetze achten. Zweisprachigkeit ist dabei kein Hindernis, sondern ein Gewinn. Schulen und Kitas fördern Mehrsprachigkeit. Ein gelungenes Zusammenleben braucht Toleranz, Neugier auf das Fremde und den Mut, Vorurteile abzubauen. Nur so wird aus einer Gesellschaft nebeneinander eine Gemeinschaft miteinander.
\end{textebox}

\begin{center}
\begin{tabularx}{\linewidth}{|X|X|}
\hline
\rowcolor{popblueL}
\textbf{Allemand (Glossaire)} & \textbf{Français} \\
\hline
\cle{das Einwanderungsland, -er} & le pays d'immigration \\
\cle{der Arbeitsmigrant, -en} & le travailleur migrant (Gastarbeiter) \\
\cle{die Verfolgung} & la persécution \\
\cle{die Fachkraft, -kräfte} & le/la spécialiste, main-d'œuvre qualifiée \\
\cle{bereichern} & enrichir \\
\cle{ökonomisch} & économique / économiquement \\
\cle{der Wohnungsmarkt, -\"{a}rkte} & le marché du logement \\
\cle{die Jobsuche} & la recherche d'emploi \\
\cle{das Klischee, -s} & le cliché, le stéréotype \\
\cle{reichen} & suffire \\
\cle{der Rechtsschutz} & la protection juridique \\
\cle{beide Seiten} & les deux parties / les deux côtés \\
\cle{die Anstrengung, -en} & l'effort \\
\cle{der Zuwanderer, -} & l'immigré (celui qui vient s'installer) \\
\cle{achten} & respecter (lois, règles) \\
\cle{das Hindernis, -se} & l'obstacle \\
\cle{der Gewinn, -e} & le gain, l'avantage \\
\cle{die Kita} (=\cle{Kindertagesstätte}) & la crèche / garderie \\
\cle{fördern} & favoriser, encourager, soutenir \\
\cle{nebeneinander} & côte à côte (sans mélange) \\
\cle{miteinander} & ensemble (en interaction) \\
\cle{abbauen} & démanteler, réduire, abattre \\
\hline
\end{tabularx}
\end{center}

\begin{exercice}[Questions sur le texte]
\begin{enumerate}
    \item Selon le texte, quelle proportion de la population a un passé migratoire ?
    \item Citez trois domaines où la discrimination se manifeste concrètement.
    \item Quel article du Grundgesetz interdit la discrimination ? Que dit-il (selon le texte) ?
    \item Pourquoi l'intégration est-elle décrite comme un processus des « deux côtés » (\cle{beide Seiten}) ?
    \item Quel est l'avantage du bilinguisme (\cle{Zweisprachigkeit}) selon l'auteur ?
    \item Expliquez la dernière phrase : \all{„Nur so wird aus einer Gesellschaft nebeneinander eine Gemeinschaft miteinander.“}
\end{enumerate}
\end{exercice}

\begin{corrige}[Exercice : Questions sur « Leben in Vielfalt »]
\begin{enumerate}
    \item Presque un habitant sur quatre (\all{fast jeder vierte Einwohner}).
    \item Marché du logement (\cle{Wohnungsmarkt}), recherche d'emploi (\cle{Jobsuche}), vie quotidienne (\cle{Alltag}).
    \item Article 3. Il interdit la discrimination/désavantage (\cle{Benachteiligung}) en raison de l'origine (\cle{wegen der Herkunft}).
    \item Parce que la société doit accepter la diversité et montrer du respect, ET les immigrants doivent apprendre l'allemand et respecter les lois. C'est une responsabilité partagée.
    \item Le bilinguisme n'est pas un obstacle (\cle{Hindernis}), mais un gain/avantage (\cle{Gewinn}).
    \item \cle{nebeneinander} = coexistence passive, sans liens (côte à côte). \cle{miteinander} = communauté active, solidaire (ensemble). Le texte dit que la tolérance et la lutte contre les préjugés transforment la simple coexistence en vraie communauté.
\end{enumerate}
\end{corrige}

\courssub{Familles de mots : Dérivation et enrichissement lexical}

Travailler par familles permet de multiplier le vocabulaire actif à partir d'une racine unique. C'est une stratégie gagnante pour l'expression écrite.

\begin{center}
\begin{tabularx}{\linewidth}{|X|X|X|X|}
\hline
\rowcolor{popgoldL}
\textbf{Verbe (Infinitif)} & \textbf{Nom (die/der/das)} & \textbf{Adjectif / Participe} & \textbf{Autres dérivés} \\
\hline
\all{integrieren} & \cle{die Integration, -en} & \cle{integriert} / \cle{integrativ} & \cle{der Integrationskurs, -e} (cours d'intégration) \\
\all{diskriminieren} & \cle{die Diskriminierung, -en} & \cle{diskriminierend} & \cle{diskriminierungsfrei} (sans discrimination) \\
\all{urteilen} & \cle{das Urteil, -e} & \cle{urteilsfähig} & \cle{das Vorurteil, -e}, \cle{vorurteilsfrei} \\
\all{assimilieren} & \cle{die Assimilation} & \cle{assimiliert} & \cle{Assimilationsdruck} (pression d'assimilation) \\
\all{segregeren} & \cle{die Segregation} & \cle{segregiert} & \cle{getrennt} (séparé, part. de \cle{trennen}) \\
\all{integrieren} $\neq$ \all{assimilieren} & & & \textbf{Distinction clé :} Intégration = participation en gardant son identité ; Assimilation = abandon de sa culture d'origine. \\
\hline
\end{tabularx}
\end{center}

\begin{exoresolu}[Transformation lexicale (Nominalisation)]
\textbf{Énoncé :} Transformez les verbes entre parenthèses en noms (avec article et majuscule) pour compléter les phrases.
\begin{enumerate}
    \item Die \_\_\_\_\_\_\_\_\_\_\_\_ (integrieren) der Minderheiten ist eine gesamtgesellschaftliche Aufgabe.
    \item Durch \_\_\_\_\_\_\_\_\_\_\_\_ (diskriminieren) werden Chancen vertan.
    \item Wir müssen \_\_\_\_\_\_\_\_\_\_\_\_ (ausgrenzen) abbauen.
    \item \_\_\_\_\_\_\_\_\_\_\_\_ (teilhaben) am politischen Leben ist ein Grundrecht.
\end{enumerate}

\textbf{Solution détaillée :}
\begin{enumerate}
    \item \cle{die Integration} (Nom fém. en \cle{-ung} $\rightarrow$ \cle{die Integration}). \all{Die Integration der Minderheiten...}
    \item \cle{die Diskriminierung} (Nom fém. en \cle{-ung}). \all{Durch die Diskriminierung...} (après préposition \cle{durch} + accusatif $\rightarrow$ \cle{die Diskriminierung}).
    \item \cle{die Ausgrenzung} (Nom fém. en \cle{-ung} de \cle{ausgrenzen}). \all{Wir müssen die Ausgrenzung abbauen.} (Accusatif singulier).
    \item \cle{die Teilhabe} (Nom fém. en \cle{-e} ; attention : \cle{teilhaben} $\rightarrow$ \cle{die Teilhabe}, pas de \cle{-ung}). \all{Die Teilhabe am politischen Leben...}
\end{enumerate}
\end{exoresolu}

\courssub{Expressions utiles pour l'expression écrite/orale (Redemittel)}

Ces blocs phraséologiques (chunks) sont à apprendre par cœur pour fluidifier vos productions (résumé, commentaire, débat).

\begin{center}
\begin{tabularx}{\linewidth}{|X|X|}
\hline
\rowcolor{popmuted}
\textbf{Deutsch} & \textbf{Français / Contexte} \\
\hline
\all{Ein zentrales Problem ist die Diskriminierung.} & Un problème central est la discrimination. (Introduction) \\
\all{Viele Migranten leiden unter Vorurteilen.} & Beaucoup de migrants souffrent de préjugés. (Constat) \\
\all{Die Gesetzgebung schützt die Rechte der Minderheiten.} & La législation protège les droits des minorités. (Argument droit) \\
\all{Integration gelingt nur durch gegenseitigen Respekt.} & L'intégration ne réussit que par le respect mutuel. (Thèse) \\
\all{Man muss zwischen Integration und Assimilation unterscheiden.} & Il faut distinguer intégration et assimilation. (Nuance) \\
\all{Zweisprachigkeit ist eine Bereicherung, kein Defizit.} & Le bilinguisme est un enrichissement, pas un déficit. (Argument fort) \\
\all{Wir brauchen mehr Toleranz und weniger Ausgrenzung.} & Nous avons besoin de plus de tolérance et moins d'exclusion. (Appel) \\
\all{Ein friedliches Zusammenleben ist möglich, wenn...} & Une coexistence pacifique est possible si... (Condition) \\
\hline
\end{tabularx}
\end{center}

\begin{methode}[Apprendre le vocabulaire thématique efficacement]
\begin{enumerate}
    \item \textbf{Par champs sémantiques :} Apprenez les 5 tableaux ci-dessus par cœur (allemand $\rightarrow$ français ET français $\rightarrow$ allemand).
    \item \textbf{Avec articles et pluriels :} Ne séparez jamais le nom de son article (\cle{die Minderheit, -en}). Utilisez des fiches recto-verso.
    \item \textbf{Par familles :} Construisez l'arbre \cle{integrieren $\rightarrow$ Integration $\rightarrow$ integriert $\rightarrow$ Integrationskurs}.
    \item \textbf{En contexte :} Écrivez 5 phrases personnelles avec les \cle{Redemittel} du tableau ci-dessus (ex. : situation au Cameroun, à l'école, dans votre quartier).
    \item \textbf{Révision espacée :} Relisez le texte \emph{Leben in Vielfalt} à J+1, J+3, J+7.
\end{enumerate}
\end{methode}

\begin{attention}[Pièges fréquents pour francophones]
\begin{itemize}
    \item \textbf{Genre de \cle{Minderheit / Mehrheit} :} Tous deux sont \textbf{féminins} (\cle{die}). Ne pas confondre avec \cle{der Minderheit} (datif singulier) ou \cle{der Mehrheit}.
    \item \textbf{\cle{der Ausländer} vs \cle{der Einwanderer} (Pluriel) :} Le pluriel est \cle{die Ausländer} / \cle{die Einwanderer} (pas de \cle{-n} ni \cle{-s}).
    \item \textbf{\cle{integrieren} vs \cle{sich integrieren} :} \all{Die Schule integriert ihn} (accusatif) MAIS \all{Er integriert sich} (réfléchi accusatif). Oublier le \cle{sich} change le sens (il intègre qqn d'autre).
    \item \textbf{\cle{zweisprachig} (adj.) vs \cle{die Zweisprachigkeit} (nom) :} \all{Er ist zweisprachig.} (pas de \cle{-e} à la fin en prédicat) / \all{Die Zweisprachigkeit ist wichtig.}
    \item \textbf{Faux ami \cle{das Vorurteil} :} Ne pas traduire par « préjugé favorable ». En allemand comme en français, \cle{Vorurteil} a une connotation \textbf{négative} (jugement hâtif, injuste). Pour « opinion préalable neutre », on dit \cle{die Voreinstellung} ou \cle{die Annahme}.
\end{itemize}
\end{attention}

\coursec{Grammaire 1 : la déclinaison de l'adjectif}

Après avoir exploré le lexique spécifique aux minorités et à l'intégration, nous devons maintenant maîtriser l'outil grammatical qui permet d'insérer ces noms dans des phrases correctes : \cle{la déclinaison de l'adjectif épithète}. En allemand, l'adjectif placé devant le nom (épithète) \textbf{s'accorde en genre, en nombre et en cas}. C'est la différence majeure avec le français, où l'adjectif ne s'accorde qu'en genre et en nombre. La marque de cet accord est portée soit par l'article (ou l'équivalent d'article), soit par l'adjectif lui-même, selon une logique rigoureuse que nous allons décortiquer.

\courssub{Observation dans le texte support}

Reprenons trois extraits du texte \emph{Ein Brief aus Bautzen} (section 1) pour observer l'adjectif épithète en situation réelle.

\begin{textebox}[Extraits du texte « Ein Brief aus Bautzen »]
„In \textbf{die sorbische Minderheit} wird viel investiert.“ \\
„Es ist \textbf{eine Form von Diskriminierung}, wenn man die Sprache verbietet.“ \\
„Manche denken, wir seien keine \textbf{richtigen Deutschen}.“
\end{textebox}

\noindent
\textbf{Analyse rapide :}
\begin{itemize}
    \item \all{die sorbische Minderheit} : article défini \cle{die} (féminin, nominatif/accusatif singulier) + adjectif en \textbf{-e}.
    \item \all{eine Form} : article indéfini \cle{eine} (féminin, nominatif/accusatif singulier) + adjectif absent ici, mais \all{von Diskriminierung} montre le cas (génitif après \cle{Form}).
    \item \all{keine richtigen Deutschen} : article négatif \cle{keine} (pluriel) + adjectif en \textbf{-en} ; nom \cle{Deutsche} (substantivé, se décline comme un adjectif, pluriel \cle{-n}).
\end{itemize}

\courssub{Principes généraux}

\begin{propriete}[Règle fondamentale : l'adjectif épithète]
L'adjectif épithète (placé entre l'article et le nom) \textbf{prend toujours une terminaison}. Cette terminaison dépend de trois paramètres :
\begin{enumerate}
    \item Le \textbf{genre} du nom (masculin, féminin, neutre).
    \item Le \textbf{nombre} (singulier, pluriel).
    \item Le \textbf{cas} (nominatif, accusatif, datif, génitif).
\end{enumerate}
Le choix de la terminaison dépend surtout de \textbf{ce qui précède l'adjectif} (article défini, indéfini/négatif/possessif, ou rien du tout). On distingue donc \textbf{trois déclinaisons} : faible, mixte, forte.
\end{propriete}

\begin{aretenir}[Logique de la marque de cas]
Le système allemand veut que le \textbf{cas soit visible une fois} dans le groupe nominal.
\begin{itemize}
    \item Si l'article \textbf{porte déjà la marque du cas} (ex. \cle{der}, \cle{die}, \cle{dem}, \cle{des}) $\rightarrow$ l'adjectif \textbf{se repose} : il prend une terminaison \textbf{faible} (\cle{-e} ou \cle{-en}).
    \item Si l'article \textbf{ne montre pas le genre/cas} (ex. \cle{ein}, \cle{mein}, \cle{kein} au masc. nom. / neut. nom./acc.) $\rightarrow$ l'adjectif \textbf{doit porter la marque} : déclinaison \textbf{mixte}.
    \item S'il n'y a \textbf{pas d'article} $\rightarrow$ l'adjectif \textbf{porte toute la charge} : déclinaison \textbf{forte}.
\end{itemize}
\end{aretenir}

\courssub{Les trois tableaux de déclinaison}

Voici les trois tableaux de référence. Apprenez-les par cœur en comprenant la logique (où est la marque du cas ?), pas par cœur bête.

\begin{propriete}[Tableau 1 : Déclinaison FAIBLE (après article défini \cle{der}, \cle{die}, \cle{das}, \cle{den}, \cle{dem}, \cle{des}...)]
L'article montre clairement le genre, le nombre et le cas. L'adjectif n'a que deux formes : \textbf{-e} (nominatif sg. + accusatif fém./neut.) et \textbf{-en} (tous les autres cas).
\begin{center}
\begin{tabularx}{\linewidth}{|>{\bfseries}X|X|X|X|X|}\hline
\rowcolor{popblueL}
Cas & Masculin & Féminin & Neutre & Pluriel \\ \hline
Nominatif & der \all{gute} Mann & die \all{gute} Frau & das \all{gute} Kind & die \all{guten} Kinder \\ \hline
Accusatif & den \all{guten} Mann & die \all{gute} Frau & das \all{gute} Kind & die \all{guten} Kinder \\ \hline
Datif & dem \all{guten} Mann & der \all{guten} Frau & dem \all{guten} Kind & den \all{guten} Kindern \\ \hline
Génitif & des \all{guten} Mannes & der \all{guten} Frau & des \all{guten} Kindes & der \all{guten} Kinder \\ \hline
\end{tabularx}
\end{center}
\emph{Exemple du texte :} \all{die sorbische Minderheit} (Nom./Acc. fém. sg. $\rightarrow$ \cle{-e}) ; \all{mit der sorbischen Sprache} (Dat. fém. sg. $\rightarrow$ \cle{-en}).
\end{propriete}

\begin{propriete}[Tableau 2 : Déclinaison MIXTE (après \cle{ein}, \cle{kein}, \cle{mein}, \cle{dein}, \cle{sein}, \cle{ihr}, \cle{unser}, \cle{euer}, \cle{Ihr})]
L'article \textbf{ne distingue pas} le masculin nominatif (\cle{ein}) ni le neutre nominatif/accusatif (\cle{ein}) des autres genres. L'adjectif \textbf{doit donc prendre la marque forte} (celle de l'article défini) à ces cas précis. Ailleurs, il prend \cle{-en}.
\begin{center}
\begin{tabularx}{\linewidth}{|>{\bfseries}X|X|X|X|X|}\hline
\rowcolor{poporangeL}
Cas & Masculin & Féminin & Neutre & Pluriel \\ \hline
Nominatif & ein \all{guter} Mann & eine \all{gute} Frau & ein \all{gutes} Kind & \cle{keine} \all{guten} Kinder \\ \hline
Accusatif & einen \all{guten} Mann & eine \all{gute} Frau & ein \all{gutes} Kind & \cle{keine} \all{guten} Kinder \\ \hline
Datif & einem \all{guten} Mann & einer \all{guten} Frau & einem \all{guten} Kind & \cle{keinen} \all{guten} Kindern \\ \hline
Génitif & eines \all{guten} Mannes & einer \all{guten} Frau & eines \all{guten} Kindes & \cle{keiner} \all{guten} Kinder \\ \hline
\end{tabularx}
\end{center}
\cle{Attention !} Au masculin nominatif et au neutre nominatif/accusatif, l'adjectif prend la terminaison de l'article défini (\cle{-er}, \cle{-es}). C'est la \textbf{erreur n°1 des francophones} (voir boîte \texttt{attention} plus bas).
\end{propriete}

\begin{propriete}[Tableau 3 : Déclinaison FORTE (sans article, ou après \cle{viel}, \cle{wenig}, \cle{mehrere}, \cle{einige}, \cle{alle} au pluriel)]
L'adjectif est seul porteur de l'information de cas. Il reprend \textbf{presque intégralement les terminaisons de l'article défini} (sauf génitif masc./neut. : \cle{-en} au lieu de \cle{-es}).
\begin{center}
\begin{tabularx}{\linewidth}{|>{\bfseries}X|X|X|X|X|}\hline
\rowcolor{popgreenL}
Cas & Masculin & Féminin & Neutre & Pluriel \\ \hline
Nominatif & \all{guter} Wein & \all{gute} Milch & \all{gutes} Wasser & \all{gute} Freunde \\ \hline
Accusatif & \all{guten} Wein & \all{gute} Milch & \all{gutes} Wasser & \all{gute} Freunde \\ \hline
Datif & \all{gutem} Wein & \all{guter} Milch & \all{gutem} Wasser & \all{guten} Freunden \\ \hline
Génitif & \all{guten} Weines & \all{guter} Milch & \all{guten} Wassers & \all{guter} Freunde \\ \hline
\end{tabularx}
\end{center}
\emph{Note :} Au génitif masc./neut., l'adjectif prend \cle{-en} (pas \cle{-es}) car le nom porte déjà un \cle{-s} ou \cle{-es} (\all{guten Weines}, \all{guten Wassers}).
\end{propriete}

\courssub{Schéma récapitulatif visuel}

\begin{popfigure}
\begin{tikzpicture}[
    node distance=1.2cm and 1.5cm,
    box/.style={draw, thick, rounded corners, minimum width=4.5cm, minimum height=3.2cm, align=center, font=\small, fill=popcream},
    title/.style={font=\bfseries\large, text=popdark, yshift=0.5cm},
    arrow/.style={-Stealth, thick, popline}
]
% Nœuds
\node[box, fill=popblue!20, align=center] (faible){
    \textbf{Article défini}\\
    (der, die, das, dem, den, des)\\
    \vspace{0.2cm}
    $\rightarrow$ Terminaison : \textcolor{popblue}{\textbf{-e / -en}}\\
    \vspace{0.2cm}
    \textit{Ex: \all{die sorbische Sprache}}
};
\node[box, fill=poporange!20, right=of faible, align=center] (mixte){
    \textbf{Ein / Kein / Possessif}\\
    (ein, keine, mein, dein...)\\
    \vspace{0.2cm}
    $\rightarrow$ Terminaison : \textcolor{poporange}{\textbf{Mixte}}\\
    (forte si article neutre, sinon -en)\\
    \vspace{0.2cm}
    \textit{Ex: \all{ein guter Mann}, \all{keine richtigen Deutschen}}
};
\node[box, fill=popgreen!20, right=of mixte, align=center] (forte){
    \textbf{Sans article}\\
    (ou viel, wenig, mehrere...)\\
    \vspace{0.2cm}
    $\rightarrow$ Terminaison : \textcolor{popgreen}{\textbf{Forte}}\\
    (comme article défini, sauf gén. -en)\\
    \vspace{0.2cm}
    \textit{Ex: \all{guter Wein}, \all{kalte Milch}}
};

% Flèches décoratives
\draw[arrow] (faible.east) -- (mixte.west) node[midway, above, font=\tiny, text=popmuted] {Moins de marque sur l'article};
\draw[arrow] (mixte.east) -- (forte.west) node[midway, above, font=\tiny, text=popmuted] {Plus de marque sur l'adjectif};
\end{tikzpicture}
\legende{Les trois types de déclinaison de l'adjectif épithète : le choix dépend de l'article (ou de son absence) qui précède.}
\end{popfigure}

\courssub{Cas particuliers et points de vigilance}

\begin{exemplebox}[Adjectifs invariables (pas de terminaison)]
Certains adjectifs (souvent des emprunts récents ou des mots d'argot) ne se déclinent \textbf{jamais} :
\cle{super, prima, extra, cool, lila, rosa, beige, orange}.
\all{Ein super Film.} \quad \all{Eine prima Idee.} \quad \all{Das lila T-Shirt.}
\end{exemplebox}

\begin{exemplebox}[Adjectif devant nom propre : apposition vs attribut]
\begin{itemize}
    \item \textbf{En apposition (sans article, titre/fonction) :} l'adjectif reste \textbf{invariant} (forme de base, souvent \cle{-e} ou \cle{-er}).
    \all{alt Bundeskanzler Scholz} \quad \all{Stadtpräsidentin Müller}
    \item \textbf{En attribut (avec article) :} l'adjectif se décline \textbf{normalement}.
    \all{der alte Mann} \quad \all{das schöne Berlin} \quad \all{in der schönen Stadt}
\end{itemize}
\end{exemplebox}

\begin{exemplebox}[Après \cle{etwas}, \cle{nichts}, \cle{viel}, \cle{wenig}, \cle{alles} (indéfinis neutres)]
Ces mots fonctionnent comme un neutre. L'adjectif qui suit prend une majuscule (substantivé) et la terminaison \textbf{-es} (fort neutre nominatif/accusatif).
\all{etwas \textbf{Neues}} (quelque chose de nouveau) \\
\all{nichts \textbf{Besonderes}} (rien de spécial) \\
\all{viel \textbf{Wichtiges}} (beaucoup de choses importantes)
\end{exemplebox}

\begin{attention}[Erreur fréquente n°1 : « ein gute Mann » au lieu de « ein guter Mann »]
En français, on dit « un bon homme ». L'article « un » ne change pas. En allemand, \cle{ein} au masculin nominatif \textbf{ne montre pas le genre} (pas de \cle{-r}). \textbf{L'adjectif DOIT le faire.}
\begin{center}
\begin{tabular}{|l|l|l|}\hline
\rowcolor{poppinkL}
\textbf{Faux (francophone)} & \textbf{Correct} & \textbf{Explication} \\ \hline
ein \textbf{gute} Mann & ein \textbf{guter} Mann & Masc. Nom. : article neutre $\rightarrow$ adj. fort \cle{-er} \\ \hline
ein \textbf{kleine} Kind & ein \textbf{kleines} Kind & Neut. Nom./Acc. : article neutre $\rightarrow$ adj. fort \cle{-es} \\ \hline
kein \textbf{schöne} Frau & keine \textbf{schöne} Frau & Fém. : article \cle{keine} montre le genre $\rightarrow$ adj. faible \cle{-e} \\ \hline
\end{tabular}
\end{center}
\emph{Astuce :} Si l'article est \cle{ein} (sans \cle{-r}, \cle{-s}, \cle{-m}, \cle{-n}), l'adjectif \textbf{vole} la terminaison de \cle{der}/\cle{das} (\cle{-er}, \cle{-es}).
\end{attention}

\begin{attention}[Erreur fréquente n°2 : Oublier le \cle{-n} au datif pluriel]
Après article défini (\cle{den}), indéfini (\cle{keinen}) ou sans article, l'adjectif au datif pluriel prend \textbf{toujours -en}, et le nom prend souvent \textbf{-n} (sauf pluriel en \cle{-s}).
\all{mit den \textbf{kleinen} Kind\textbf{ern}} \quad \all{mit \textbf{kleinen} Kind\textbf{ern}} \quad \all{mit \textbf{guten} Freunden}
\end{attention}

\begin{methode}[Méthode en 4 étapes pour décliner correctement]
Pour ne jamais vous tromper dans un exercice ou une rédaction, suivez cet algorithme :
\begin{enumerate}
    \item \textbf{Identifiez le nom :} Quel est son genre ? (M, F, N) Son nombre ? (Sg, Pl)
    \item \textbf{Identifiez le cas :} Quelle est la fonction du groupe nominal ? (Sujet $\rightarrow$ Nom ; COD $\rightarrow$ Acc ; COI/Prép. datif $\rightarrow$ Dat ; Complément du nom $\rightarrow$ Gén)
    \item \textbf{Regardez ce qui précède l'adjectif :} Article défini (\cle{der}/\cle{dem}...) ? Article indéfini/possessif/négatif (\cle{ein}/\cle{mein}/\cle{kein}) ? Rien ?
    \item \textbf{Choisissez le bon tableau} (Faible / Mixte / Forte) et piochez la terminaison à l'intersection [Genre $\times$ Cas].
\end{enumerate}
\end{methode}

\courssub{Comparaison français / allemand}

\begin{center}
\begin{tabularx}{\linewidth}{|>{\bfseries}X|X|X|}\hline
\rowcolor{poppurpleL}
\textbf{Critère} & \textbf{Français} & \textbf{Allemand} \\ \hline
Accord de l'adjectif épithète & Genre + Nombre seulement & \textbf{GENRE + NOMBRE + CAS} \\ \hline
Marque de l'accord & Sur l'adjectif (souvent muette à l'oral) & Sur l'article \textbf{ET/OU} l'adjectif (audible) \\ \hline
Exemple : « le petit garçon » & \textit{le petit garçon} (masc. sg.) & \all{der \textbf{kleine} Junge} (Faible -e) \\ \hline
Exemple : « un petit garçon » & \textit{un petit garçon} & \all{ein \textbf{kleiner} Junge} (Mixte -er) \\ \hline
Exemple : « petits garçons » (sans art.) & \textit{petits garçons} & \all{\textbf{kleine} Jungen} (Forte -e) \\ \hline
Datif « au petit garçon » & \textit{au petit garçon} (pas de cas visible) & \all{dem \textbf{kleinen} Jungen} (Faible -en + -n nom) \\ \hline
\end{tabularx}
\end{center}

\courssub{Exercices d'application sur le thème des minorités}

\begin{exoresolu}[Compléter les terminaisons adjectivales]
Complétez les phrases suivantes en mettant l'adjectif entre parenthèses à la bonne forme. Le cas est indiqué.

\begin{enumerate}
    \item Die Sorben sind \_\_\_\_\_\_\_\_\_\_ (klein) Minderheit in Deutschland. (Nominatif)
    \item Man hört oft \_\_\_\_\_\_\_\_\_\_ (sorbisch) Sprache in der Lausitz. (Accusatif)
    \item Wir sprechen mit \_\_\_\_\_\_\_\_\_\_ (alt) Leuten über Traditionen. (Datif)
    \item Das ist \_\_\_\_\_\_\_\_\_\_ (wichtig) Recht für alle. (Nominatif, après article indéfini)
    \item Ohne \_\_\_\_\_\_\_\_\_\_ (eigen) Sprache verliert man Identität. (Accusatif, sans article)
    \item Viele \_\_\_\_\_\_\_\_\_\_ (jung) Sorben lernen die Sprache in der Schule. (Nominatif pluriel, sans article)
\end{enumerate}

\tcblower
\textbf{Solution détaillée :}
\begin{enumerate}
    \item \textbf{kleine} — \cle{Die} (art. déf. fém. sg. Nom.) $\rightarrow$ Faible \cle{-e}. \textit{« Les Sorbes sont une petite minorité... »}
    \item \textbf{sorbische} — \cle{die} (art. déf. fém. sg. Acc.) $\rightarrow$ Faible \cle{-e}. \textit{« On entend souvent la langue sorabe... »}
    \item \textbf{alten} — \cle{den} (art. déf. plur. Dat.) $\rightarrow$ Faible \cle{-en} + \cle{-n} au nom (\cle{Leuten}). \textit{« Nous parlons avec de vieux gens... »}
    \item \textbf{wichtiges} — \cle{ein} (art. indéf. neutre sg. Nom.) $\rightarrow$ Mixte : article neutre $\rightarrow$ adj. fort \cle{-es}. \textit{« C'est un droit important... »}
    \item \textbf{eigene} — Sans article, fém. sg. Acc. $\rightarrow$ Forte : comme \cle{die} $\rightarrow$ \cle{-e}. \textit{« Sans langue propre... »}
    \item \textbf{junge} — Sans article, plur. Nom. $\rightarrow$ Forte : comme \cle{die} $\rightarrow$ \cle{-e}. \textit{« Beaucoup de jeunes Sorbes... »}
\end{enumerate}
\end{exoresolu}

\begin{exercice}[Mise en situation : Droits des minorités]
Mettez l'adjectif entre parenthèses à la forme correcte (genre, nombre, cas). Attention aux articles !

\begin{enumerate}
    \item In \_\_\_\_\_\_\_\_\_\_ (bundesrepublikanisch) Deutschland gibt es vier anerkannte Minderheiten.
    \item Die \_\_\_\_\_\_\_\_\_\_ (dänisch) Minderheit lebt im Norden.
    \item Man respektiert \_\_\_\_\_\_\_\_\_\_ (friesisch) Kultur und Sprache.
    \item Ein \_\_\_\_\_\_\_\_\_\_ (gut) Integrationspolitik braucht Geld und Zeit.
    \item Ohne \_\_\_\_\_\_\_\_\_\_ (politisch) Wille ändert sich nichts.
    \item Wir helfen \_\_\_\_\_\_\_\_\_\_ (neu) Zugewanderten bei der Integration.
\end{enumerate}
\end{exercice}

\begin{corrige}[Exercice : Mise en situation]
\begin{enumerate}
    \item \textbf{bundesrepublikanischem} — \cle{In} + Dat. sg. neutre (\cle{Deutschland}), \textbf{sans article} (noms de pays) $\rightarrow$ Forte \cle{-em}.
    \item \textbf{dänische} — \cle{Die} (art. déf. fém. sg. Nom.) $\rightarrow$ Faible \cle{-e}.
    \item \textbf{friesische} — \cle{die} (art. déf. fém. sg. Acc.) $\rightarrow$ Faible \cle{-e}.
    \item \textbf{gute} — \cle{Eine} (art. indéf. fém. sg. Nom.) $\rightarrow$ Mixte : art. montre genre $\rightarrow$ Faible \cle{-e}.
    \item \textbf{politischen} — \cle{Ohne} + Acc. sg. masc. (\cle{Wille}) $\rightarrow$ Sans article $\rightarrow$ Forte : comme \cle{den} $\rightarrow$ \cle{-en}.
    \item \textbf{neuen} — \cle{helfen} + Dat. plur. (\cle{Zugewanderten}) $\rightarrow$ Faible \cle{-en} (art. déf. \cle{den} implicite) + \cle{-n} au nom (faible).
\end{enumerate}
\end{corrige}

\begin{savaistu}[Le savais-tu ? : La déclinaison au Baccalauréat]
Au Bac allemand (LV2), la \textbf{qualité de la langue} (note sur 4 à 6 points selon les académies) sanctionne lourdement les erreurs de terminaison adjectivale. Une copie qui maîtrise les trois déclinaisons (faible, mixte, forte) et évite le piège \cle{ein guter Mann} / \cle{ein kleines Kind} gagne \textbf{au moins 1 à 2 points de langue} par rapport à une copie qui laisse des adjectifs « nus » (\cle{ein gut Mann}). Entraînez-vous à décliner \textbf{systématiquement} chaque groupe nominal de vos rédactions : c'est le meilleur investissement temps/points de l'examen !
\end{savaistu}

\coursec{Grammaire 2 : la concession et le subjonctif II du souhait}

Dans la section précédente, nous avons appris à décliner l'adjectif pour décrire précisément les réalités sociales : \all{eine kleine Minderheit}, \all{die eigenen Rechte}, \all{das deutsche Schulsystem}. Mais décrire ne suffit pas : pour commenter le problème des minorités, il faut aussi \cle{nuancer}, \cle{opposer des idées} et \cle{exprimer des souhaits}. C'est exactement ce que permettent les deux outils grammaticaux de cette section : la \cle{concession} (\all{obwohl}, \all{trotzdem}) et le \cle{subjonctif II du souhait} (\all{wäre}, \all{hätte}, \all{gäbe}...). Ce sont deux points de grammaire très fréquents au baccalauréat, aussi bien en compréhension de texte qu'en thème.

\courssub{Observation : deux constructions pour une même idée}

Reprenons deux phrases du texte support \all{Ein Brief aus Bautzen} :

\begin{exemplebox}[Deux phrases du texte]
\all{Die Sorben haben ihre eigene Sprache, obwohl viele Deutsche das nicht wissen.}\\
« Les Sorabes ont leur propre langue, bien que beaucoup d'Allemands l'ignorent. »

\all{Die Sorben haben ihre eigene Sprache. Trotzdem sprechen die meisten nur Deutsch.}\\
« Les Sorabes ont leur propre langue. Pourtant, la plupart ne parlent que l'allemand. »
\end{exemplebox}

Dans les deux cas, l'énonciateur oppose deux faits : un fait principal (les Sorabes ont leur langue) et un fait qui semble le contredire (on l'ignore, on ne la parle pas). C'est le rapport logique de \cle{concession}. Mais observez bien la \textbf{place du verbe} :

\begin{itemize}
\item avec \all{obwohl} : \all{obwohl viele Deutsche das nicht \textbf{wissen}} → le verbe conjugué est \textbf{à la fin} ;
\item avec \all{trotzdem} : \all{\textbf{Trotzdem sprechen} die meisten...} → le verbe conjugué est \textbf{en deuxième position}, et le sujet vient après (inversion).
\end{itemize}

Même idée, deux constructions différentes : c'est le cœur de cette leçon.

\courssub{La concession avec \all{obwohl} : conjonction de subordination}

\begin{definition}[La concession]
La concession est un rapport logique qui présente un obstacle réel mais insuffisant pour empêcher l'action principale. En français : « bien que », « quoique », « malgré », « pourtant », « néanmoins ». En allemand : \all{obwohl} (bien que), \all{trotzdem} / \all{dennoch} (pourtant), \all{trotz} + génitif (malgré).
\end{definition}

\begin{propriete}[\all{obwohl} : verbe conjugué à la fin]
\all{Obwohl} est une \textbf{conjonction de subordination}. Elle introduit une \textbf{proposition subordonnée} : le verbe conjugué se place \textbf{en dernière position} de cette subordonnée.

\all{Obwohl die Minderheit klein \textbf{ist}, hat sie eigene Rechte.}\\
« Bien que la minorité soit petite, elle a ses propres droits. »

Variantes soutenues, identiques en tout point : \all{obgleich} et \all{obschon} (registre littéraire, fréquent dans les textes de presse au bac).
\end{propriete}

Deux configurations sont possibles, comme en français :

\begin{enumerate}
\item \textbf{Subordonnée d'abord} : \all{Obwohl er krank \textbf{ist}, \textbf{geht} er zur Schule.} — la subordonnée entière occupe la position 1, donc le verbe de la principale vient juste après la virgule, \textbf{avant le sujet} (\all{geht er}, inversion).
\item \textbf{Principale d'abord} : \all{Er \textbf{geht} zur Schule, obwohl er krank \textbf{ist}.} — ordre normal dans la principale.
\end{enumerate}

\begin{attention}[\all{obwohl} ne demande PAS de subjonctif]
En français, « bien que » exige le subjonctif : « bien qu'il \emph{soit} malade ». En allemand, \all{obwohl} introduit un fait réel et s'emploie avec l'\textbf{indicatif} : \all{obwohl er krank \textbf{ist}} (et jamais une forme de subjonctif). Ne calquez pas le subjonctif français sur l'allemand !
\end{attention}

\begin{exemplebox}[Exemples traduits]
\all{Obwohl die Sorben nur etwa 60 000 Menschen sind, haben sie eigene Schulen.} « Bien que les Sorabes ne soient qu'environ 60 000, ils ont leurs propres écoles. »

\all{Obgleich sie seit Generationen in Deutschland leben, kennen viele Deutsche ihre Kultur nicht.} « Bien qu'ils vivent en Allemagne depuis des générations, beaucoup d'Allemands ne connaissent pas leur culture. »

\all{Die Sinti und Roma kämpfen für ihre Rechte, obwohl sie oft diskriminiert werden.} « Les Sintis et les Roms luttent pour leurs droits, bien qu'ils soient souvent discriminés. »
\end{exemplebox}

\begin{attention}[Piège de la virgule et de la place du verbe]
Deux erreurs classiques des francophones :
\begin{itemize}
\item Oublier l'inversion quand la subordonnée est en tête : on écrit \all{Obwohl er krank ist, er geht zur Schule.} — FAUX. Il faut : \all{Obwohl er krank ist, \textbf{geht er} zur Schule.} La subordonnée compte comme la position 1.
\item Laisser le verbe en position 2 dans la subordonnée : \all{obwohl er ist krank} — FAUX. Le verbe conjugué va \textbf{à la fin} : \all{obwohl er krank \textbf{ist}}.
\end{itemize}
\end{attention}

\courssub{La concession avec \all{trotzdem} : adverbe connecteur}

\begin{propriete}[\all{trotzdem} / \all{dennoch} : verbe en position 2]
\all{Trotzdem} (pourtant, néanmoins) n'est \textbf{pas} une conjonction : c'est un \textbf{adverbe connecteur}. Il relie deux \textbf{phrases indépendantes} (séparées par un point ou un point-virgule). Placé en tête de phrase, il occupe la position 1 et provoque \textbf{l'inversion} : le verbe conjugué reste en position 2, le sujet passe après.

\all{Er ist diskriminiert worden. \textbf{Trotzdem bleibt er} optimistisch.}\\
« Il a été discriminé. Pourtant, il reste optimiste. »

Variante plus soutenue : \all{dennoch}. On peut aussi placer l'adverbe après le verbe : \all{Er bleibt trotzdem optimistisch.}
\end{propriete}

\begin{attention}[Inversion obligatoire après \all{trotzdem}]
Ne jamais écrire : \all{Trotzdem er spricht Deutsch.} — FAUX.\\
La seule forme correcte : \all{\textbf{Trotzdem spricht er} Deutsch.}\\
Règle d'or valable pour \textbf{tous} les adverbes en tête de phrase (\all{heute}, \all{dann}, \all{deshalb}, \all{leider}...) : position 1 + verbe + sujet.
\end{attention}

\begin{attention}[Jamais \all{obwohl} et \all{trotzdem} dans la même phrase]
Le français dit « bien que... pourtant... ». L'allemand considère cela comme un \textbf{pléonasme fautif}. Il faut \textbf{choisir} l'une des deux structures :

\all{Obwohl es Probleme gibt, bleibt er optimistisch.} (correct)\\
\all{Es gibt Probleme. Trotzdem bleibt er optimistisch.} (correct)\\
\all{Obwohl es Probleme gibt, bleibt er trotzdem optimistisch.} (à éviter absolument)
\end{attention}

\courssub{La transformation \all{obwohl} $\leftrightarrow$ \all{trotzdem}}

C'est un exercice classique du baccalauréat. Voici le mécanisme :

\begin{methode}[Transformer une phrase concessive]
\begin{enumerate}
\item Identifier les deux faits opposés : A (fait principal) et B (obstacle).
\item \textbf{De \all{obwohl} vers \all{trotzdem}} : faire de la subordonnée une phrase indépendante (verbe en position 2), mettre un point, puis commencer la phrase principale par \all{Trotzdem} + verbe + sujet.
\item \textbf{De \all{trotzdem} vers \all{obwohl}} : transformer la phrase de l'obstacle en subordonnée introduite par \all{obwohl} (verbe à la fin), supprimer \all{trotzdem}, joindre par une virgule.
\item Vérifier la place de chaque verbe conjugué : fin dans la subordonnée, position 2 ailleurs.
\end{enumerate}
\end{methode}

\begin{center}
\begin{tabularx}{\linewidth}{|X|X|}
\hline
\rowcolor{popblueL} \textbf{Subordination : \all{obwohl}} & \textbf{Deux phrases : \all{trotzdem}} \\
\hline
\all{Obwohl er krank ist, geht er zur Schule.} & \all{Er ist krank. Trotzdem geht er zur Schule.} \\
\hline
\all{Obwohl sie diskriminiert werden, kämpfen sie weiter.} & \all{Sie werden diskriminiert. Trotzdem kämpfen sie weiter.} \\
\hline
\all{Obwohl die Sprache bedroht ist, lernen Kinder sie noch.} & \all{Die Sprache ist bedroht. Trotzdem lernen Kinder sie noch.} \\
\hline
\end{tabularx}
\end{center}

\begin{popfigure}
\begin{tikzpicture}[
box/.style={draw=popblue, fill=popblueL, rounded corners=4pt, align=center, text width=5.6cm, inner sep=6pt, font=\small},
arr/.style={-{Stealth[length=3mm]}, very thick, poporange}]
\node[box, align=center] (a) at (0,0){\textbf{obwohl}\\[2pt] \all{Obwohl er krank \underline{ist}, geht er zur Schule.}\\[2pt] subordonnée : verbe \textbf{à la fin}};
\node[box, fill=poporangeL, draw=poporange, align=center] (b) at (9.2,0){\textbf{trotzdem}\\[2pt] \all{Er ist krank. \underline{Trotzdem geht} er zur Schule.}\\[2pt] phrase 2 : verbe en \textbf{position 2}, inversion};
\draw[arr] (a.east) -- node[above, font=\small\itshape, text=popdark]{on casse en deux phrases} (b.west);
\draw[arr] (b.south) .. controls (9.2,-2.2) and (0,-2.2) .. node[below, font=\small\itshape, text=popdark]{on fusionne : verbe à la fin} (a.south);
\end{tikzpicture}
\legende{Même idée de concession, deux constructions : subordination avec \all{obwohl} ou juxtaposition avec \all{trotzdem}.}
\end{popfigure}

\textbf{Complément utile pour le bac : la préposition \all{trotz} + génitif.}\par
Quand l'obstacle est un \textbf{nom} (et non une phrase entière), l'allemand utilise la préposition \all{trotz} (« malgré »), suivie du \textbf{génitif} :

\all{Trotz der Probleme bleibt die Gemeinschaft stark.} « Malgré les problèmes, la communauté reste forte. »

\all{Trotz jahrhundertelanger Diskriminierung hat die sorbische Sprache überlebt.} « Malgré des siècles de discrimination, la langue sorabe a survécu. »

Dans la langue courante, on entend parfois \all{trotz} + datif (\all{trotz dem Wetter}), mais au baccalauréat, exigez le génitif : \all{trotz des Wetters}. Attention à ne pas confondre \all{trotz} (préposition, « malgré ») et \all{trotzdem} (adverbe, « pourtant ») : \all{trotz} est toujours suivi d'un nom, \all{trotzdem} introduit une phrase complète.

\courssub{Le subjonctif II du souhait}

Après avoir décrit la situation des minorités, notre auteur exprime un \textbf{souhait} : que les choses soient différentes. En allemand, le souhait irréel s'exprime avec le \cle{Konjunktiv II}.

\begin{methode}[Former le Konjunktiv II]
\begin{enumerate}
\item Partir du \textbf{prétérit} du verbe.
\item Pour les \textbf{verbes forts et irréguliers fréquents}, ajouter un \textbf{tréma} sur la voyelle radicale : \all{war} → \all{wäre} ; \all{hatte} → \all{hätte} ; \all{gab} → \all{gäbe} ; \all{konnte} → \all{könnte} ; \all{wusste} → \all{wüsste}.
\item Pour \textbf{tous les autres verbes}, utiliser \all{würde} + infinitif : \all{ich würde helfen} (j'aiderais).
\item Pour le \textbf{passé} (souhait irréalisable) : \all{hätte / wäre} + participe passé : \all{ich hätte das gewusst} (si j'avais su cela).
\end{enumerate}
\end{methode}

\begin{center}
\begin{tabular}{|l|l|l|}
\hline
\rowcolor{popblueL} \textbf{Infinitif} & \textbf{Präteritum} & \textbf{Konjunktiv II} \\
\hline
\all{sein} (être) & \all{war} & \all{wäre} (serait) \\
\hline
\all{haben} (avoir) & \all{hatte} & \all{hätte} (aurait) \\
\hline
\all{geben} (\all{es gibt}) & \all{gab} & \all{gäbe} (il y aurait) \\
\hline
\all{können} (pouvoir) & \all{konnte} & \all{könnte} (pourrait) \\
\hline
\all{mögen} (aimer) & \all{mochte} & \all{möchte} (aimerait, voudrait) \\
\hline
\all{wissen} (savoir) & \all{wusste} & \all{wüsste} (saurait) \\
\hline
\all{werden} + inf. & \all{wurde} & \all{würde} + inf. (ferait) \\
\hline
\end{tabular}
\end{center}

\begin{propriete}[Les trois formes de la phrase de souhait]
\begin{enumerate}
\item \textbf{Phrase déclarative} avec \all{ich wünschte} : \all{Ich wünschte, alle hätten mehr Respekt.} « Je voudrais que tous aient plus de respect. »
\item \textbf{Phrase en \all{wenn}} avec particule : \all{Wenn es doch mehr Toleranz gäbe!} « S'il y avait seulement plus de tolérance ! » — le verbe conjugué va à la fin (subordonnée).
\item \textbf{Phrase sans \all{wenn}} : le verbe conjugué en \textbf{première position} : \all{Wäre das Leben nur einfacher!} « Si seulement la vie était plus simple ! »
\end{enumerate}
Les particules \all{doch}, \all{nur}, \all{bloß} (« seulement ») renforcent le souhait et sont quasiment obligatoires dans les formes 2 et 3.
\end{propriete}

\begin{exemplebox}[Souhaits traduits]
\all{Wenn es doch keine Vorurteile gäbe!} « S'il n'y avait pas de préjugés ! »

\all{Hätten die Minderheiten nur mehr Rechte!} « Si seulement les minorités avaient plus de droits ! »

\all{Ich wünschte, die Leute würden mehr über andere Kulturen lernen.} « Je voudrais que les gens apprennent davantage sur les autres cultures. »

\all{Könnten wir doch alle in Frieden zusammenleben!} « Si seulement nous pouvions tous vivre ensemble en paix ! »
\end{exemplebox}

\textbf{Le souhait irréalisable au passé.}\par
Pour regretter quelque chose qui ne s'est pas passé, on utilise le Konjunktiv II au passé : \all{hätte / wäre} + participe passé, le participe allant à la fin :

\all{Wenn ich das nur früher gewusst hätte!} « Si seulement j'avais su cela plus tôt ! »

\all{Wäre man nur früher gegen die Diskriminierung eingetreten!} « Si seulement on s'était engagé plus tôt contre la discrimination ! » (\all{eintreten} se conjugue avec \all{sein} → \all{wäre ... eingetreten}).

\textbf{Comparaison français / allemand.}\par

\begin{center}
\begin{tabularx}{\linewidth}{|X|X|X|}
\hline
\rowcolor{popblueL} \textbf{Français} & \textbf{Deutsch} & \textbf{Remarque} \\
\hline
Bien qu'il soit malade, il va à l'école. & \all{Obwohl er krank ist, geht er zur Schule.} & subordonnée : verbe à la fin, indicatif \\
\hline
Il est malade. Pourtant il va à l'école. & \all{Er ist krank. Trotzdem geht er zur Schule.} & inversion après \all{trotzdem} \\
\hline
Malgré les problèmes & \all{trotz der Probleme} & préposition + génitif \\
\hline
S'il y avait plus de tolérance ! & \all{Wenn es doch mehr Toleranz gäbe!} & français : imparfait ; allemand : Konj. II \\
\hline
Je voudrais que tout soit différent. & \all{Ich wünschte, alles wäre anders.} & le français utilise le subjonctif, pas l'allemand \\
\hline
\end{tabularx}
\end{center}

\begin{savaistu}[D'où vient \all{gäbe} ?]
\all{Gäbe} vient de \all{es gab}, le prétérit de \all{es gibt} (« il y a »). La formule \all{Wenn es doch ... gäbe!} est l'une des phrases de souhait les plus fréquentes de la langue allemande, à l'écrit comme à l'oral : \all{Wenn es doch mehr Zeit gäbe!} (s'il y avait plus de temps !), \all{Wenn es doch keinen Krieg gäbe!} (s'il n'y avait pas de guerre !). Retenez-la comme un bloc : elle impressionne toujours un correcteur de production écrite.
\end{savaistu}

\courssub{Texte d'application : souhaits et concessions}

\begin{textebox}[Stimmen aus der Lausitz]
\all{„Ich bin Sorbin, und ich bin stolz darauf“, sagt Maria, 17 Jahre alt, aus Bautzen. „Obwohl meine Familie seit Jahrhunderten hier lebt, fragen mich manche Leute, woher ich ‚wirklich‘ komme. Das verletzt mich. Trotzdem antworte ich immer freundlich: Ich komme von hier!“

Ihr Bruder Jakob sieht die Lage kritischer: „Obwohl das Sorbische offiziell geschützt ist, hört man es in der Stadt kaum noch. Viele junge Sorben sprechen nur Deutsch. Trotz der Probleme lerne ich die Sprache meiner Großeltern weiter. Wenn es doch mehr sorbische Sendungen im Fernsehen gäbe! Und wäre es nicht schön, wenn alle Kinder in der Lausitz zweisprachig aufwachsen könnten?“

Auch ihre Mutter mischt sich ins Gespräch: „Obgleich wir manchmal Vorurteile erleben, ist unsere Lage heute viel besser als früher. Dennoch bleibt viel zu tun. Ich wünschte, alle Deutschen wüssten, dass Minderheiten kein Problem sind, sondern ein Reichtum. Wenn die Menschen nur mehr miteinander reden würden! Hätten wir das nur schon vor zwanzig Jahren gelernt!“

Am Ende lacht Maria: „Obwohl die Welt nicht perfekt ist, gebe ich die Hoffnung nicht auf. Trotzdem — oder gerade deshalb — bleibe ich optimistisch.“}
\end{textebox}

\textbf{Glossaire :} \all{die Sorbin, -nen} (la Sorabe, femme) ; \all{stolz} (fier) ; \all{verletzen} (blesser) ; \all{die Lage, -n} (la situation) ; \all{geschützt} (protégé) ; \all{die Sendung, -en} (l'émission) ; \all{zweisprachig aufwachsen} (grandir bilingue) ; \all{sich ins Gespräch mischen} (se mêler à la conversation) ; \all{erleben} (vivre, éprouver) ; \all{der Reichtum, die Reichtümer} (la richesse) ; \all{die Hoffnung aufgeben} (abandonner l'espoir) ; \all{gerade deshalb} (justement pour cela).

\textbf{Questions de grammaire sur le texte :}
\begin{enumerate}
\item Relevez deux subordonnées avec \all{obwohl} / \all{obgleich} et soulignez le verbe conjugué à la fin.
\item Relevez deux phrases avec \all{trotzdem} / \all{dennoch} et justifiez l'inversion.
\item Relevez trois phrases de souhait et identifiez la forme du Konjunktiv II utilisée.
\end{enumerate}

\courssub{Exercices}

\begin{exoresolu}[Transformation \all{obwohl} $\leftrightarrow$ \all{trotzdem}]
\textbf{Énoncé.} 1) Transformez avec \all{trotzdem} : \all{Obwohl die Minderheit klein ist, hat sie eigene Rechte.} 2) Transformez avec \all{obwohl} : \all{Die Sprache ist bedroht. Trotzdem lernen viele Kinder sie.}
\tcblower
\textbf{Solution.} 1) On casse la subordonnée en phrase indépendante (verbe en position 2), puis \all{Trotzdem} + verbe + sujet : \all{Die Minderheit ist klein. Trotzdem hat sie eigene Rechte.} « La minorité est petite. Pourtant, elle a ses propres droits. » 2) On transforme la première phrase en subordonnée (verbe à la fin), on supprime \all{trotzdem} : \all{Obwohl die Sprache bedroht ist, lernen viele Kinder sie.} « Bien que la langue soit menacée, beaucoup d'enfants l'apprennent. » Vérification : dans la subordonnée, \all{ist} est bien en dernière position ; dans la principale, le verbe est en position 2 avec inversion (\all{lernen viele Kinder}).
\end{exoresolu}

\begin{exoresolu}[Mise au Konjunktiv II du souhait]
\textbf{Énoncé.} Exprimez le souhait avec la particule indiquée : 1) \all{Es gibt mehr Respekt.} (\all{nur}) 2) \all{Das Leben ist einfacher.} (\all{doch}, sans \all{wenn}) 3) \all{Die Leute können miteinander reden.} (\all{doch})
\tcblower
\textbf{Solution.} 1) \all{gab} → \all{gäbe} : \all{Wenn es nur mehr Respekt gäbe!} « S'il y avait seulement plus de respect ! » 2) \all{war} → \all{wäre}, verbe en première position : \all{Wäre das Leben doch einfacher!} « Si seulement la vie était plus simple ! » 3) \all{konnten} → \all{könnten} : \all{Wenn die Leute doch miteinander reden könnten!} « Si seulement les gens pouvaient parler entre eux ! »
\end{exoresolu}

\begin{exercice}[À vous]
\begin{enumerate}
\item Transformez avec \all{trotzdem} : \all{Obwohl er oft diskriminiert wird, bleibt er optimistisch.}
\item Transformez avec \all{obwohl} : \all{Es gibt viele Vorurteile. Trotzdem leben die Menschen friedlich zusammen.}
\item Mettez au Konjunktiv II du souhait : a) \all{Es gibt keine Diskriminierung.} (\all{doch}) b) \all{Alle Kinder lernen zwei Sprachen.} (\all{nur}, avec \all{würden}) c) \all{Ich habe das früher gewusst.} (\all{nur}, au passé)
\item Traduisez : « Bien que les problèmes soient grands, la communauté ne perd pas espoir. » puis : « Malgré les préjugés, elle reste fière de sa langue. »
\end{enumerate}
\end{exercice}

\begin{corrige}[Exercice « À vous »]
1) \all{Er wird oft diskriminiert. Trotzdem bleibt er optimistisch.} — phrase indépendante + inversion après \all{trotzdem}.

2) \all{Obwohl es viele Vorurteile gibt, leben die Menschen friedlich zusammen.} — verbe \all{gibt} à la fin de la subordonnée ; inversion dans la principale (\all{leben die Menschen}).

3) a) \all{Wenn es doch keine Diskriminierung gäbe!} b) \all{Wenn nur alle Kinder zwei Sprachen lernen würden!} (ou \all{Würden nur alle Kinder zwei Sprachen lernen!}) c) \all{Wenn ich das nur früher gewusst hätte!} — passé irréel : \all{hätte} + participe \all{gewusst}.

4) \all{Obwohl die Probleme groß sind, verliert die Gemeinschaft die Hoffnung nicht.} / \all{Trotz der Vorurteile bleibt sie stolz auf ihre Sprache.} — attention au génitif après \all{trotz} : \all{der Vorurteile}.
\end{corrige}

\begin{aretenir}[L'essentiel de la section]
\begin{itemize}
\item \all{Obwohl} = conjonction de subordination → verbe conjugué \textbf{à la fin}, à l'\textbf{indicatif} : \all{Obwohl er krank ist, geht er zur Schule.}
\item \all{Trotzdem} / \all{dennoch} = adverbe connecteur → \textbf{inversion} : \all{Trotzdem geht er zur Schule.}
\item Jamais \all{obwohl} et \all{trotzdem} dans la même phrase ; \all{trotz} + génitif = « malgré » : \all{trotz der Probleme}.
\item Souhait : Konjunktiv II = prétérit + tréma (\all{wäre}, \all{hätte}, \all{gäbe}, \all{könnte}, \all{möchte}) ou \all{würde} + infinitif ; particules \all{doch}, \all{nur}, \all{bloß} ; passé irréel : \all{hätte/wäre} + participe.
\end{itemize}
\end{aretenir}

\coursec{Commentaire modèle et production guidée}

Dans la section précédente, nous avons appris à exprimer la concession avec \all{obwohl} et \all{trotzdem}, et à formuler un souhait au subjonctif II : \all{Ich wünschte, alle Menschen wären toleranter.} « J'aimerais que tous les hommes soient plus tolérants. » Ces outils grammaticaux ne sont pas des exercices isolés : ils vont maintenant nous servir à \cle{commenter} le texte « Ein Brief aus Bautzen » et à \cle{produire} nos propres textes. C'est exactement ce qu'on vous demandera à l'épreuve écrite du Baccalauréat : comprendre, résumer, donner son avis — en allemand correct.

\courssub{La méthode du commentaire de texte}

Un commentaire de texte (\all{ein Kommentar}) n'est ni un résumé pur ni une simple opinion : c'est un texte organisé en trois grandes parties, comme en français, auxquelles on ajoute souvent une prise de position personnelle avant la conclusion.

\begin{methode}[Commenter un texte en trois étapes]
\begin{enumerate}
\item \textbf{Introduction (\all{Einleitung})} : présenter le document — de quoi s'agit-il ? qui parle ? de quel thème ? Utiliser des formules comme \all{Der Text ist ein Leserbrief von...} « Le texte est une lettre au journal de... » ou \all{Es geht um das Problem der Minderheiten.} « Il s'agit du problème des minorités. »
\item \textbf{Développement / Analyse (\all{Hauptteil})} : dégager les idées principales dans l'ordre du texte et montrer \emph{comment} l'auteur s'exprime (exemples, ton personnel, questions directes au lecteur). Connecteurs utiles : \all{zuerst}, \all{dann}, \all{außerdem}, \all{schließlich}.
\item \textbf{Conclusion (\all{Schluss})} : donner son avis personnel argumenté — \all{meiner Meinung nach}, \all{ich finde, dass...} — et éventuellement ouvrir sur la situation au Cameroun.
\end{enumerate}
\end{methode}

\begin{popfigure}
\begin{center}
\begin{tikzpicture}[
  node distance=0.55cm,
  box/.style={draw=popblue, thick, fill=popblueL, rounded corners=3pt, align=center, minimum width=3.1cm, minimum height=1.1cm, font=\small},
  arr/.style={-{Stealth[length=2.5mm]}, thick, popdark}
]
\node[box, align=center] (intro){\textbf{Einleitung}\\ Autor, Quelle, Thema};
\node[box, right=of intro, align=center] (analyse){\textbf{Analyse}\\ Ideen + Verfahren};
\node[box, right=of analyse, align=center] (meinung){\textbf{Meinung}\\ meiner Meinung nach};
\node[box, right=of meinung, align=center] (schluss){\textbf{Schluss}\\ Fazit + Öffnung};
\draw[arr] (intro) -- (analyse);
\draw[arr] (analyse) -- (meinung);
\draw[arr] (meinung) -- (schluss);
\end{tikzpicture}
\end{center}
\legende{Organigramme du commentaire de texte : quatre étapes enchaînées.}
\end{popfigure}

\courssub{Les connecteurs d'opinion}

Pour passer de l'analyse à l'avis personnel, il faut des connecteurs d'opinion. Observons trois modèles :

\begin{exemplebox}[Exprimer son opinion]
\begin{itemize}
\item \all{Ich finde, dass die Soraben mehr Rechte brauchen.} « Je trouve que les Sorabes ont besoin de plus de droits. »
\item \all{Meiner Meinung nach ist Integration sehr wichtig.} « À mon avis, l'intégration est très importante. »
\item \all{Ich bin der Meinung, dass man Vorurteile bekämpfen muss.} « Je suis d'avis qu'il faut combattre les préjugés. »
\end{itemize}
\end{exemplebox}

\begin{propriete}[Règle de construction]
\begin{itemize}
\item \all{ich finde, dass...} et \all{ich bin der Meinung, dass...} introduisent une \textbf{subordonnée} : le verbe conjugué va \textbf{à la fin} (\all{..., dass man Vorurteile bekämpfen muss}).
\item \all{meiner Meinung nach} se place en tête de phrase et le \textbf{verbe reste en 2\textsuperscript{e} position} : \all{Meiner Meinung nach \textbf{ist} Integration wichtig.}
\item La tournure \all{nach meiner Meinung} existe aussi et est correcte, mais la forme la plus fréquente et la plus élégante est \all{meiner Meinung nach} (nom au génitif + \all{nach} postposé).
\end{itemize}
\end{propriete}

\begin{attention}[Piège typique du francophone]
Ne dites jamais \all{Meiner Meinung nach, ich finde, dass...} : c'est un double marqueur, comme « à mon avis, je trouve que » en français. Choisissez \textbf{un seul} connecteur. Autre erreur : oublier l'inversion après \all{meiner Meinung nach} — le verbe suit immédiatement : \all{Meiner Meinung nach \textbf{ist}...}, et non \all{Meiner Meinung nach Integration ist...}. Attention aussi à l'orthographe : \all{die Meinung} prend une majuscule (c'est un nom), et \all{die Diskriminierung} s'écrit avec \textbf{-ierung}, comme \all{die Integration}.
\end{attention}

\courssub{Commentaire modèle : « Ein Brief aus Bautzen »}

\begin{textebox}[Début de commentaire modèle]
Der Text ist ein Leserbrief von Tomas, einem jungen Soraben aus Bautzen. Er beschreibt die Situation seiner Minderheit in Sachsen. Zuerst erklärt er, dass die Soraben seit Jahrhunderten in Deutschland leben und eine eigene Sprache haben. Dann spricht er über die Vorurteile, die viele Deutsche gegen Minderheiten haben. Außerdem zeigt er, dass die sorbische Kultur in Schulen und Medien weiterlebt. Schließlich fordert er mehr Respekt und Toleranz. Meiner Meinung nach ist dieser Brief sehr mutig, weil Tomas offen über Diskriminierung spricht. Ich wünschte, alle Menschen wären so tolerant wie er.
\end{textebox}

\textbf{Glossaire} : der Leserbrief, -e (lettre au journal / courrier des lecteurs) ; der Sorabe, -n (le Sorabe — nom faible : \all{einem jungen Sorab\textbf{en}}, \all{die Sorab\textbf{en}}) ; das Jahrhundert, -e (siècle) ; die Vorurteile (les préjugés, pluriel de das Vorurteil) ; weiterleben (continuer à vivre, verbe à particule inséparable) ; fordern (exiger, réclamer) ; mutig (courageux) ; offen (ouvertement, franchement).

Ce modèle fait environ huit lignes et respecte l'organigramme : introduction (deux premières phrases), analyse avec \all{zuerst / dann / außerdem / schließlich}, opinion avec \all{meiner Meinung nach}, et ouverture avec un souhait au subjonctif II. Remarquez la place des verbes : verbe conjugué en 2\textsuperscript{e} position dans les phrases principales, verbe à la fin après \all{dass} et \all{weil}.

\begin{aretenir}[La boîte à connecteurs du commentaire]
\begin{center}
\begin{tabularx}{\linewidth}{|X|X|X|}
\hline
\rowcolor{popblueL} \textbf{Fonction} & \textbf{Deutsch} & \textbf{Français} \\
\hline
Ordre des idées & zuerst, dann, außerdem, schließlich & d'abord, ensuite, de plus, enfin \\
\hline
Opinion & meiner Meinung nach ; ich finde, dass... & à mon avis ; je trouve que... \\
\hline
Concession & obwohl ; trotzdem & bien que ; pourtant \\
\hline
Souhait & Ich wünschte, ... wären... & J'aimerais que... soient... \\
\hline
Conclusion & Zum Schluss ; Deshalb finde ich... & Pour finir ; C'est pourquoi je trouve... \\
\hline
\end{tabularx}
\end{center}
\end{aretenir}

\courssub{Production guidée 1 : le résumé en quatre phrases}

Résumer (\all{zusammenfassen}) signifie : reprendre les idées essentielles \textbf{avec ses propres mots}, sans citer le texte, sans donner son avis.

\begin{methode}[Réussir un résumé court]
\begin{enumerate}
\item Repérer les quatre idées principales du texte (une par phrase).
\item Remplacer les mots du texte par des synonymes : \all{leben} $\rightarrow$ \all{wohnen} ; \all{sagen} $\rightarrow$ \all{erklären}.
\item Rédiger au Präsens, avec des phrases simples : sujet + verbe en 2\textsuperscript{e} position.
\item Ne pas écrire \all{ich} : le résumé est impersonnel.
\end{enumerate}
\end{methode}

\begin{exoresolu}[Résumé du « Brief aus Bautzen »]
Énoncé : Fasse den Leserbrief von Tomas in vier Sätzen mit deinen eigenen Worten zusammen. \tcblower
Solution possible : \all{Tomas ist ein junger Sorabe und wohnt in Bautzen in Sachsen. Seine Minderheit hat eine eigene Sprache und Kultur, aber viele Menschen kennen sie nicht. In der Schule und im Alltag gibt es manchmal Vorurteile gegen die Soraben. Tomas möchte mehr Respekt und Toleranz für alle Minderheiten in Deutschland.} — « Tomas est un jeune Sorabe et habite à Bautzen en Saxe. Sa minorité a sa propre langue et sa propre culture, mais beaucoup de gens ne la connaissent pas. À l'école et dans la vie quotidienne, il y a parfois des préjugés contre les Sorabes. Tomas souhaite plus de respect et de tolérance pour toutes les minorités en Allemagne. » Remarquez : quatre phrases, quatre idées, aucune opinion personnelle, verbes toujours en 2\textsuperscript{e} position.
\end{exoresolu}

\courssub{Production guidée 2 : paragraphe argumenté sur le Cameroun}

\begin{exercice}[Gibt es Minderheiten in Kamerun?]
Rédige un paragraphe de 8 à 10 lignes : \all{Gibt es Minderheiten in Kamerun? Welche Probleme haben sie?} Contraintes obligatoires :
\begin{itemize}
\item au moins \textbf{5 mots du lexique} du chapitre (\all{die Minderheit, die Diskriminierung, die Rechte, die Sprache, das Vorurteil, die Integration}) ;
\item une phrase avec \all{obwohl} (verbe à la fin) ;
\item une phrase avec \all{trotzdem} (verbe en 2\textsuperscript{e} position) ;
\item un souhait au subjonctif II (\all{Ich wünschte,...}).
\end{itemize}
\end{exercice}

\begin{corrige}[Production guidée 2 — modèle de rédaction]
\all{Ja, in Kamerun gibt es viele Minderheiten, zum Beispiel die Pygmäen im Süden oder kleine Sprachgruppen im Norden. Obwohl Kamerun mehr als 250 Sprachen hat, spielen viele kleine Sprachen in der Schule keine Rolle. Die Minderheiten haben oft Probleme mit Diskriminierung und Vorurteilen, und ihre Rechte werden nicht immer respektiert. Trotzdem kämpfen viele junge Leute für ihre Kultur und ihre Integration. Meiner Meinung nach ist diese Vielfalt ein Reichtum für das Land. Ich wünschte, alle Kameruner hätten die gleichen Chancen in Schule und Beruf.} — « Oui, au Cameroun il existe de nombreuses minorités, par exemple les Pygmées du Sud ou de petits groupes linguistiques du Nord. Bien que le Cameroun ait plus de 250 langues, beaucoup de petites langues ne jouent aucun rôle à l'école. Les minorités ont souvent des problèmes de discrimination et de préjugés, et leurs droits ne sont pas toujours respectés. Pourtant, beaucoup de jeunes luttent pour leur culture et leur intégration. À mon avis, cette diversité est une richesse pour le pays. J'aimerais que tous les Camerounais aient les mêmes chances à l'école et dans la profession. » Vérifiez : 5 mots du lexique employés (\all{Minderheiten, Sprachen, Diskriminierung, Vorurteilen, Rechte, Integration}), \all{obwohl} + verbe final (\all{hat}), \all{trotzdem} + verbe en 2\textsuperscript{e} position (\all{kämpfen}), subjonctif II (\all{hätten}), passif correct (\all{werden respektiert}).
\end{corrige}

\begin{attention}[Vorurteile : cas et préposition]
Dans le corrigé, on lit \all{Probleme mit Diskriminierung und Vorurteil\textbf{en}} : après la préposition \all{mit}, le nom est au \textbf{datif pluriel}, donc avec un \textbf{-n} final (\all{das Vorurteil} $\rightarrow$ dat. pl. \all{den Vorurteilen} ; sans article : \all{mit Vorurteilen}). De même : \all{Vorurteile \textbf{gegen} die Soraben} — le préjugé contre quelqu'un se construit avec \all{gegen} + accusatif, comme en français « contre ».
\end{attention}

\begin{methode}[La production écrite au Bac en quatre gestes]
\begin{enumerate}
\item \textbf{Lire la consigne} et souligner les mots-clés (thème, longueur, contraintes grammaticales).
\item \textbf{Faire un plan rapide} au brouillon : 3 ou 4 idées dans l'ordre logique.
\item \textbf{Rédiger} avec des phrases simples mais correctes — mieux vaut une phrase courte juste qu'une phrase longue fausse.
\item \textbf{Vérifier} : les terminaisons des adjectifs (\all{die kleine Minderheit / kleine Minderheiten}), la place des verbes (2\textsuperscript{e} position, fin après \all{dass/obwohl/weil}), la majuscule aux noms.
\end{enumerate}
\end{methode}

\courssub{Grille d'auto-évaluation}

Après avoir rédigé, évaluez-vous honnêtement avec cette grille sur 5 points :

\begin{center}
\begin{tabularx}{\linewidth}{|X|X|X|}
\hline
\rowcolor{popblueL} \textbf{Critère} & \textbf{Points} & \textbf{Je vérifie...} \\
\hline
Vocabulaire & 2 pts & 5 mots du lexique, correctement orthographiés, avec majuscule \\
\hline
Grammaire & 2 pts & \all{obwohl} + verbe final ; \all{trotzdem} + verbe 2\textsuperscript{e} ; subjonctif II correct \\
\hline
Contenu & 1 pt & le paragraphe répond aux deux questions de la consigne \\
\hline
\end{tabularx}
\end{center}

\begin{experience}[Débat de classe : Toleranz im Alltag]
Par groupes de quatre, préparez en cinq minutes deux arguments pour et deux arguments contre la phrase : \all{Minderheiten müssen ihre Sprache in der Schule lernen dürfen.} « Les minorités doivent pouvoir apprendre leur langue à l'école. » Chaque élève prend la parole au moins une fois et doit utiliser un connecteur d'opinion (\all{meiner Meinung nach}, \all{ich finde, dass...}) et une concession (\all{obwohl} ou \all{trotzdem}). Un rapporteur par groupe présente ensuite la conclusion commune en trois phrases.
\end{experience}

\begin{aretenir}[L'essentiel de la section]
Un commentaire suit toujours le schéma \textbf{Einleitung $\rightarrow$ Analyse $\rightarrow$ Meinung $\rightarrow$ Schluss}. Les connecteurs \all{zuerst, dann, außerdem, schließlich} structurent l'analyse ; \all{meiner Meinung nach} (verbe en 2\textsuperscript{e} position !) et \all{ich finde, dass...} (verbe à la fin) expriment l'opinion. Un résumé reprend les idées avec ses propres mots, sans avis personnel. Une production argumentée au Bac gagne des points quand elle intègre volontairement les outils du chapitre : lexique, concession et subjonctif II du souhait.
\end{aretenir}

\newpage

\coursec{Activité d'intégration}

\coursec{Le problème des minorités}

\courssub{Texte support : « Ein Brief aus Bautzen »}
\begin{textebox}[Source : Simulé — Lettre d'un lycéen sorabe (Allemagne, Saxe) au « Kamerun-Kurier »]
\large \textbf{Ein Brief aus Bautzen / Budyšin}

\medskip
\noindent
Liebe Redaktion,

\medskip
\noindent
mein Name ist Jan Nowak, ich bin 18 Jahre alt und besuche das \textbf{sorbische Gymnasium} in Bautzen. Ich gehöre zur \textbf{sorbischen Minderheit} in der Lausitz. Bei uns in der Schule wird auf \textbf{sorbisch} und auf \textbf{deutsch} unterrichtet. Das ist unser \textbf{Recht}, das im Grundgesetz und im Sorbengesetz geschützt ist.

\medskip
\noindent
\textbf{Obwohl} wir diese Rechte haben, erlebe ich oft \textbf{Vorurteile}. Manche sagen: „Sorbisch ist eine nutzlose Sprache, das braucht niemand.“ \textbf{Trotzdem} spreche ich sorbisch zu Hause und in der Freizeit mit \textbf{großer Freude}. Es ist ein Teil meiner \textbf{kulturellen Identität}.

\medskip
\noindent
In meinem Viertel in Bautzen gibt es \textbf{keine sorbischen Straßenschilder} mehr, nur deutsche. Das finde ich \textbf{schade}. \textbf{Ich wünschte, die Gesellschaft wäre offener} für unsere Sprache. \textbf{Man sollte} die \textbf{Vielfalt} als \textbf{Reichtum} sehen und nicht als Problem.

\medskip
\noindent
Mit freundlichen Grüßen aus der Lausitz \\
Jan Nowak
\end{textebox}

\begin{definition}[Glossaire du texte (Mots nouveaux / difficiles)]
\begin{center}
\begin{tabularx}{\linewidth}{|l|l|X|}
\hline
\rowcolor{poppurpleL} \textbf{Deutsch (Art. + Pl.)} & \textbf{Français} & \textbf{Contexte / Collocation} \\ \hline
\all{das Gymnasium, -en} & le lycée (général) & \all{besuchen} (fréquenter), \all{sorbisches Gymnasium} \\ \hline
\all{die Lausitz (kein Pl.)} & la Lusace (région) & \all{in der Lausitz leben}, \all{ober-/niedersorbisch} \\ \hline
\all{das Sorbengesetz, -e} & la loi sur les Sorabes & \all{geschützt sein durch} \\ \hline
\all{nutzlos} & inutile & \all{nutzlose Sprache}, \all{keinen Nutzen haben} \\ \hline
\all{die Freizeit, -n} & le temps libre & \all{in der Freizeit}, \all{Freizeitaktivität} \\ \hline
\all{der Reichtum, -e} & la richesse & \all{kultureller Reichtum}, \all{Reichtum an Sprachen} \\ \hline
\all{offen} & ouvert (d'esprit) & \all{offen für etw. (Akku) sein}, \all{offene Gesellschaft} \\ \hline
\all{sehen} & voir, considérer & \all{als Reichtum sehen}, \all{als Problem sehen} \\ \hline
\end{tabularx}
\end{center}
\end{definition}

\courssub{Compréhension du texte}
\begin{exercice}[Questions de compréhension globale et détaillée]
\textbf{Consigne :} Répondez en allemand (phrases complètes) ou en français selon la question.
\begin{enumerate}
    \item \textbf{Global :} Wer schreibt den Brief? Woher kommt er? An wen schreibt er? (Qui écrit ? D'où ? À qui ?)
    \item \textbf{Détail :} Welche zwei Sprachen werden an Jans Schule unterrichtet? (Quelles deux langues ?)
    \item \textbf{Grammaire / Sens :} Was bedeutet „Obwohl wir diese Rechte haben, erlebe ich oft Vorurteile“? Erklären Sie den logischen Zusammenhang. (Expliquez le lien logique).
    \item \textbf{Vocabulaire :} Finden Sie im Text zwei Adjektive, die auf \textbf{-los} und \textbf{-reich} enden, und übersetzen Sie sie. (Trouvez deux adjectifs en \textbf{-los} et \textbf{-reich}).
    \item \textbf{Expression :} Formulieren Sie Jans zwei Wünsche (Konjunktiv II) auf Französisch.
    \item \textbf{Analyse :} Warum nennt Jan das Fehlen sorbischer Straßenschilder ein Beispiel für mangelnde Integration? (Argumentation).
\end{enumerate}
\end{exercice}

\begin{corrige}[Exercice Compréhension]
\begin{enumerate}
    \item \all{Jan Nowak, 18 Jahre alt, aus Bautzen in der Lausitz, schreibt an die Redaktion des Kamerun-Kuriers.} / Jan Nowak, 18 ans, de Bautzen en Lusace, écrit à la rédaction du Kamerun-Kurier.
    \item \all{An der Schule wird auf Sorbisch und auf Deutsch unterrichtet.} / Le sorabe et l'allemand.
    \item \all{Obwohl} drückt eine Konzession aus: Trotz der vorhandenen Rechte (Konzessivsatz) folgt die Hauptsatz-Aussage, dass Vorurteile trotzdem existieren. / « Bien que » introduit la concession : malgré l'existence des droits, les préjugés persistent.
    \item \all{nutzlos} (inutile) ; \all{Reichtum} (Richesse — nom, mais adjectif \all{reich} dans \all{reichhaltig} ou ici substantivé). \textit{Note :} Le texte a \all{Reichtum} (nom). L'adjectif \all{reich} n'apparaît pas directement, mais \all{Reichtum} en dérive. L'adjectif en \textbf{-los} est \all{nutzlos}.
    \item 1. Il souhaite que la société soit plus ouverte (\all{Ich wünschte, die Gesellschaft wäre offener}). 2. Il souhaite/recommande qu'on voie la diversité comme une richesse (\all{Man sollte die Vielfalt als Reichtum sehen}).
    \item L'absence de panneaux bilingues montre que la langue minoritaire n'est pas visible/valorisée dans l'espace public, ce qui nie la présence et les droits de la minorité (manque de \all{Integration} et de \all{Respekt}).
\end{enumerate}
\end{corrige}

\courssub{Lexique thématique : Minderheiten und Integration}
\begin{propriete}[Familles de mots et collocations clés (à mémoriser)]
Le vocabulaire du thème s'organise autour de noyaux nominaux, verbaux et adjectivaux. Apprenez \textbf{l'article, le pluriel} et \textbf{une collocation} par mot.
\begin{center}
\begin{tabularx}{\linewidth}{|l|l|X|}
\hline
\rowcolor{poppurpleL} \textbf{Nom (Art. + Pl.)} & \textbf{Verbe / Adjectif} & \textbf{Collocations / Structures utiles} \\ \hline
\all{die Minderheit, -en} & \all{angehören} (dat), \all{minderheitlich} & \all{zu einer Minderheit gehören}, \all{ethnische / sprachliche / nationale Minderheit}, \all{Minderheitenrechte} \\ \hline
\all{die Diskriminierung, -en} & \all{diskriminieren}, \all{benachteiligen}, \all{diskriminierend} & \all{jemanden diskriminieren / benachteiligen}, \all{Opfer der Diskriminierung sein}, \all{gegen Diskriminierung kämpfen} \\ \hline
\all{die Integration, -en} & \all{integrieren}, \all{integriert}, \all{integrativ} & \all{in die Gesellschaft integrieren}, \all{gelungene Integration}, \all{Integrationskurs}, \all{Integrationspolitik} \\ \hline
\all{das Vorurteil, -e} & \all{vorurteilsvoll}, \all{vorurteilsfrei} & \all{Vorurteile haben / abbauen / bekämpfen / hegen}, \all{von Vorurteilen geprägt sein} \\ \hline
\all{die Sprache, -n} & \all{sprachlich}, \all{mehrsprachig}, \all{einsprachig} & \all{Muttersprache / Erstsprache / Zweitsprache}, \all{Sprachbarriere überwinden}, \all{Sprachrechte fordern} \\ \hline
\all{das Recht, -e} & \all{rechtlich}, \all{berechtigt}, \all{sich berechtigt fühlen} & \all{Rechte haben / fordern / einfordern / wahren}, \all{Grundrechte}, \all{Menschenrechte}, \all{gleichberechtigt} \\ \hline
\all{die Identität, -en} & \all{identisch}, \all{identitätsstiftend} & \all{kulturelle Identität bewahren / verlieren / stärken}, \all{Identitätskrise}, \all{Identitätssuche} \\ \hline
\all{die Toleranz, -en} & \all{tolerant}, \all{intolerant}, \all{tolerieren} & \all{Toleranz zeigen / fordern / üben}, \all{tolerant gegenüber jdm. (Dat) sein} \\ \hline
\all{die Ausgrenzung, -en} & \all{ausgrenzen}, \all{ausgegrenzt}, \all{ausgrenzend} & \all{jemanden ausgrenzen}, \all{sozial ausgegrenzt sein}, \all{ausgrenzen aufgrund von ...} \\ \hline
\all{der Respekt, -e} & \all{respektieren}, \all{respektvoll}, \all{respektlos} & \all{Respekt vor jdm./etw. (Dat) haben}, \all{respektvoll behandeln}, \all{Respektlosigkeit} \\ \hline
\end{tabularx}
\end{center}
\end{propriete}

\begin{exercice}[Entraînement lexical — Mots en contexte]
Complétez les phrases avec le mot juste (forme correcte : genre, nombre, cas).
\begin{enumerate}
    \item Viele \_\_\_\_\_\_\_\_\_\_ (Minderheit) in Europa kämpfen um ihre Sprachrechte.
    \item \_\_\_\_\_\_\_\_\_\_ (Diskriminierung) am Arbeitsplatz ist verboten.
    \item Die \_\_\_\_\_\_\_\_\_\_ (Integration) von Neuzugewanderten braucht Zeit und Geld.
    \item Er hat viele \_\_\_\_\_\_\_\_\_\_ (Vorurteil) gegen diese Gruppe.
    \item Meine \_\_\_\_\_\_\_\_\_\_ (Muttersprache) ist Bassa, aber ich spreche auch Französisch.
    \item Jeder Mensch hat \_\_\_\_\_\_\_\_\_\_ (Recht) auf Bildung.
    \item Wir müssen die \_\_\_\_\_\_\_\_\_\_ (kulturelle Identität) der Baka \_\_\_\_\_\_\_\_\_\_ (bewahren).
    \item Mehr \_\_\_\_\_\_\_\_\_\_ (Toleranz) würde das Zusammenleben verbessern.
    \item Niemand soll \_\_\_\_\_\_\_\_\_\_ (ausgrenzen) werden.
    \item Ich habe großen \_\_\_\_\_\_\_\_\_\_ (Respekt) vor deiner Leistung.
\end{enumerate}
\end{exercice}

\begin{corrige}[Exercice Lexique]
\begin{enumerate}
    \item \all{Minderheiten} (Pluriel Nominatif/Accusatif)
    \item \all{Diskriminierung} (Singulier Nominatif)
    \item \all{Integration} (Singulier Nominatif)
    \item \all{Vorurteile} (Pluriel Accusatif après \all{viele})
    \item \all{Muttersprache} (Singulier Nominatif)
    \item \all{Rechte} (Pluriel Accusatif après \all{hat} — \all{jeder Mensch hat Rechte})
    \item \all{kulturelle Identität} (Accusatif fém. sg. faible \all{die kulturelle Identität}) — \all{bewahren} (infinitif)
    \item \all{Toleranz} (Singulier Accusatif après \all{Mehr})
    \item \all{ausgegrenzt} (Participe II passif : \all{werden ausgegrenzt})
    \item \all{Respekt} (Accusatif masc. sg. après \all{großen} — \all{großen Respekt})
\end{enumerate}
\end{corrige}

\courssub{Grammaire 1 : La déclinaison de l'adjectif}
\begin{propriete}[Règle fondamentale : Le déterminant décide]
En allemand, l'adjectif épithète (attributif) se décline \textbf{obligatoirement}. Sa finale dépend \textbf{exclusivement du déterminant qui le précède} (article défini, indéfini, possessif, négatif, indéfini \all{kein}, ou absence de déterminant). Il y a \textbf{trois déclinaisons}.

\cle{Règle d'or :} \textbf{Un seul marqueur de cas/genre/nombre par groupe nominal (si possible).} L'article défini (\all{der, die, das, die}) porte la marque forte $\rightarrow$ adjectif \textbf{faible} (\all{-e} ou \all{-en}). Pas d'article / article sans marque (\all{kein, mein} au masc./neutre nom/acc) $\rightarrow$ adjectif \textbf{forte} (porte la marque). Article indéfini / possessif (avec marque) $\rightarrow$ adjectif \textbf{mixte} (\all{-e} nom/acc fém/neutre, sinon \all{-en}).
\end{propriete}

\begin{exemplebox}[Observation dans le texte support]
\begin{itemize}
    \item \all{das \textbf{sorbisch\emph{e}} Gymnasium} : Article défini \all{das} (neutre, nom/acc) $\rightarrow$ \textbf{Faible} $\rightarrow$ \textbf{-e}.
    \item \all{die \textbf{sorbisch\emph{e}} Minderheit} : Article défini \all{die} (fém, nom/acc) $\rightarrow$ \textbf{Faible} $\rightarrow$ \textbf{-e}.
    \item \all{unser \textbf{Recht}} : Possessif \all{unser} (neutre, nom/acc, sans \textbf{-s} visible mais marque de cas) $\rightarrow$ \textbf{Mixte} $\rightarrow$ \textbf{-e} (neutre nom/acc).
    \item \all{groß\emph{er} Freude} : Préposition \all{mit} + Dat. Zéro article (ou article défini \all{der} masqué ? Non, \all{mit großer Freude} $\rightarrow$ \all{großer} = \textbf{Forte} fém. dat. \textit{Attention :} Souvent \all{mit großer Freude} figé. Si \all{mit der großen Freude} $\rightarrow$ Faible \all{großen}).
    \item \all{kein\emph{e} sorbisch\emph{e} Straßenschilder} : \all{keine} (Pluriel, marque \textbf{-e}) $\rightarrow$ Adjectif \textbf{Mixte} pluriel $\rightarrow$ \textbf{-n} (\all{sorbischen}). \textit{Correction texte :} \all{keine sorbischen Straßenschilder}.
    \item \all{offener} (prédicatif après \all{sein}) : Pas de déclinaison (adjectif attribut du sujet).
\end{itemize}
\end{exemplebox}

\begin{center}
\begin{tabularx}{\linewidth}{|X|X|X|X|}
\hline
\rowcolor{popblueL} \textbf{Cas / Genre} & \textbf{Faible (Art. déf. \textit{der, die, das, die})} & \textbf{Forte (Zéro art. / \textit{kein, mein, dieser} masc./neutre N/A)} & \textbf{Mixte (Art. indéf. \textit{ein, eine, kein, mein, dein} + poss.)} \\ \hline
\textbf{Nom. Mask.} & der \textbf{gut\emph{e}} Mann & \textbf{gut\emph{er}} Mann & ein \textbf{gut\emph{er}} Mann \\ \hline
\textbf{Acc. Mask.} & den \textbf{gut\emph{en}} Mann & \textbf{gut\emph{en}} Mann & einen \textbf{gut\emph{en}} Mann \\ \hline
\textbf{Dat. Mask.} & dem \textbf{gut\emph{en}} Mann & \textbf{gut\emph{em}} Mann & einem \textbf{gut\emph{en}} Mann \\ \hline
\textbf{Gen. Mask.} & des \textbf{gut\emph{en}} Mannes & \textbf{gut\emph{en}} Mannes & eines \textbf{gut\emph{en}} Mannes \\ \hline
\textbf{Nom. Fém.} & die \textbf{gut\emph{e}} Frau & \textbf{gut\emph{e}} Frau & eine \textbf{gut\emph{e}} Frau \\ \hline
\textbf{Acc. Fém.} & die \textbf{gut\emph{e}} Frau & \textbf{gut\emph{e}} Frau & une \textbf{gut\emph{e}} Frau \\ \hline
\textbf{Dat. Fém.} & der \textbf{gut\emph{en}} Frau & \textbf{gut\emph{er}} Frau & einer \textbf{gut\emph{en}} Frau \\ \hline
\textbf{Gen. Fém.} & der \textbf{gut\emph{en}} Frau & \textbf{gut\emph{er}} Frau & einer \textbf{gut\emph{en}} Frau \\ \hline
\textbf{Nom. Neut.} & das \textbf{gut\emph{e}} Kind & \textbf{gut\emph{es}} Kind & ein \textbf{gut\emph{es}} Kind \\ \hline
\textbf{Acc. Neut.} & das \textbf{gut\emph{e}} Kind & \textbf{gut\emph{es}} Kind & ein \textbf{gut\emph{es}} Kind \\ \hline
\textbf{Dat. Neut.} & dem \textbf{gut\emph{en}} Kind & \textbf{gut\emph{em}} Kind & einem \textbf{gut\emph{en}} Kind \\ \hline
\textbf{Gen. Neut.} & des \textbf{gut\emph{en}} Kindes & \textbf{gut\emph{en}} Kindes & eines \textbf{gut\emph{en}} Kindes \\ \hline
\textbf{Pluriel (tous)} & die \textbf{gut\emph{en}} Leute & \textbf{gut\emph{e}} Leute & keine \textbf{gut\emph{en}} Leute \\ \hline
\end{tabularx}
\end{center}

\begin{attention}[Pièges fréquents pour francophones]
\begin{itemize}
    \item \textbf{Oubli du -n au datif pluriel :} \all{mit jung\emph{en} Leuten} (pas *\all{jung\emph{e} Leuten}). \textbf{Tous les cas du pluriel prennent -en en faible/mixte, -e en forte.}
    \item \textbf{Confusion Forte / Mixte au Masculin Nominatif/Accusatif :} \all{ein gut\emph{er}} Mann (Mixte, \all{ein} porte pas de marque $\rightarrow$ adj. forte \all{-er}) vs \all{kein gut\emph{er}} Mann (Mixte, \all{kein} porte \all{-n}? Non, \all{kein} comme \all{ein} $\rightarrow$ \all{kein guter Mann}). \textbf{En réalité :} Après \all{ein, kein, mein, dein, sein, ihr, unser, euer, Ihr} au Masc. Nom. et Neutre Nom./Acc. : l'article n'a \textbf{pas de finale visible} (ou \all{-r}, \all{-s} pour \all{unser/euer}) $\rightarrow$ L'adjectif prend la finale \textbf{FORTE} (\all{-er, -es}). C'est la déclinaison \textbf{MIXTE} qui \textbf{emprunte les finales fortes} là où l'article est "nu".
    \item \textbf{Majuscules :} \all{der Mann}, \all{die Frau} — L'adjectif reste minuscule (\all{der gut\emph{e} Mann}), sauf noms propres (\all{der deutsch\emph{e} Bundestag}).
\end{itemize}
\end{attention}

\begin{exercice}[Entraînement déclinaison — Complétez les finales]
Mettez l'adjectif entre parenthèses à la bonne forme.
\begin{enumerate}
    \item Ich sehe \_\_\_\_\_\_\_\_\_\_ (jung) Mann. (Accusatif)
    \item Sie hilft \_\_\_\_\_\_\_\_\_\_ (alt) Frau. (Datif)
    \item Das ist \_\_\_\_\_\_\_\_\_\_ (gut) Buch. (Nominatif, indéfini)
    \item Wir sprechen über \_\_\_\_\_\_\_\_\_\_ (neu) Auto. (Accusatif, défini)
    \item Er ist stolz auf \_\_\_\_\_\_\_\_\_\_ (kulturell) Identität. (Accusatif, défini)
    \item Ohne \_\_\_\_\_\_\_\_\_\_ (eigen) Sprache verliert man Identität. (Accusatif, possessif)
    \item \_\_\_\_\_\_\_\_\_\_ (viel) Menschen kommen. (Nominatif, indéfini \all{viel} $\rightarrow$ forte)
    \item Die \_\_\_\_\_\_\_\_\_\_ (sorbisch) Sprache ist schön. (Nominatif, défini)
\end{enumerate}
\end{exercice}

\begin{corrige}[Exercice Déclinaison]
\begin{enumerate}
    \item \all{einen jung\emph{en}} Mann (Mixte Acc. Masc. \all{-en})
    \item \all{einer alt\emph{en}} Frau (Mixte Dat. Fém. \all{-en}) / \all{der alt\emph{en}} Frau (Faible)
    \item \all{ein gut\emph{es}} Buch (Mixte Nom. Neut. \all{-es} — article \all{ein} sans marque)
    \item \all{das neu\emph{e}} Auto (Faible Acc. Neut. \all{-e})
    \item \all{seine kulturell\emph{e}} Identität (Mixte Acc. Fém. \all{-e} — \all{seine} comme \all{eine})
    \item \all{eigen\emph{e}} Sprache (Forte Acc. Fém. \all{-e} — pas d'article devant \all{eigen})
    \item \all{Vi\emph{e}le} Menschen (Forte Nom. Plur. \all{-e} — \all{viel} se décline fort)
    \item \all{sorbisch\emph{e}} Sprache (Faible Nom. Fém. \all{-e})
\end{enumerate}
\end{corrige}

\courssub{Grammaire 2 : La concession (\textit{obwohl} / \textit{trotzdem}) et le Subjonctif II (Souhait)}

\courssub{La concession : \cle{obwohl} vs \cle{trotzdem}}
\begin{propriete}[Structure et place du verbe — Tableau récapitulatif]
\begin{center}
\begin{tabularx}{\linewidth}{|X|X|X|}
\hline
\rowcolor{poporangeL} \textbf{Moyen} & \textbf{Statut syntaxique} & \textbf{Place du verbe conjugué} \\ \hline
\all{obwohl} & Subordonnante (introduit une \textbf{Nebensatz}) & \textbf{À LA FIN} de la subordonnée. \\ \hline
\all{trotzdem} & Adverbe / Conjonction de coordination (lie 2 \textbf{Hauptsätze}) & \textbf{2e POSITION} (V2) dans la principale qui suit. \\ \hline
\end{tabularx}
\end{center}

\cle{Sens commun :} Opposition, concession (« bien que / malgré cela »).
\cle{Nuance :} \all{Obwohl} met l'accent sur le fait concédé (la cause) ; \all{trotzdem} met l'accent sur la conséquence, le maintien du fait principal.

\end{propriete}

\begin{exemplebox}[Exemples authentiques (extraits / adaptés du texte)]
\begin{itemize}
    \item \textbf{Obwohl} \all{wir diese Rechte \textbf{haben},} erlebe ich oft Vorurteile. (Verbe \all{haben} à la fin).
    \item Man sagt: „Sorbisch ist nutzlos.“ \textbf{Trotzdem} spreche \all{ich} sorbisch. (\all{Trotzdem} = Place 1 $\rightarrow$ Verbe \all{spreche} = Place 2 $\rightarrow$ Sujet \all{ich} = Place 3).
    \item \textbf{Alternative :} Ich habe Vorurteile, \textbf{trotzdem} kämpfe \all{ich} für meine Rechte. (\all{trotzdem} en Place 3, Verbe \all{kämpfe} en Place 2).
\end{itemize}
\end{exemplebox}

\begin{attention}[Erreurs types : \textbf{*Obwohl ich habe Rechte} (Verbe pas à la fin) ; \textbf{*Trotzdem ich kämpfe} (Verbe pas en 2e). Astuce : \all{Obwohl} = « Queue » (verbe à la fin) ; \all{Trotzdem} = « Deux » (verbe en 2e).]
\end{attention}

\courssub{Le Subjonctif II (Konjunktiv II) : Exprimer un souhait irréel du présent}
\begin{propriete}[Formation et emploi — Niveau Terminale A]
Le \textbf{Konjunktiv II} (Subjonctif II) exprime l'irréel du présent (souhait, hypothèse, politesse). Pour le souhait (\all{Wunsch}), on utilise souvent \all{ich wünschte, ...} + subordonnée (verbe à la fin) ou \all{wenn} + KII.

\begin{center}
\begin{tabular}{|l|l|l|l|l|l|l|}
\hline
\rowcolor{popgreenL} \textbf{Infinitif} & \textbf{ich} & \textbf{du} & \textbf{er/sie/es} & \textbf{wir} & \textbf{ihr} & \textbf{sie/Sie} \\ \hline
\all{sein} & \all{wäre} & \all{wärest} & \all{wäre} & \all{wären} & \all{wäret} & \all{wären} \\ \hline
\all{haben} & \all{hätte} & \all{hättest} & \all{hätte} & \all{hätten} & \all{hättet} & \all{hätten} \\ \hline
\all{können} & \all{könnte} & \all{könntest} & \all{könnte} & \all{könnten} & \all{könntet} & \all{könnten} \\ \hline
\all{sollen} & \all{sollte} & \all{solltest} & \all{sollte} & \all{sollten} & \all{solltet} & \all{sollten} \\ \hline
\all{wollen} & \all{wollte} & \all{wolltest} & \all{wollte} & \all{wollten} & \all{wolltet} & \all{wollten} \\ \hline
\all{müssen} & \all{müsste} & \all{müsstest} & \all{müsste} & \all{müssten} & \all{müsstet} & \all{müssten} \\ \hline
\all{dürfen} & \all{dürfte} & \all{dürftest} & \all{dürfte} & \all{dürften} & \all{dürftet} & \all{dürften} \\ \hline
\all{werden} & \all{würde} & \all{würdest} & \all{würde} & \all{würden} & \all{würdet} & \all{würden} \\ \hline
\end{tabular}
\end{center}

\cle{Autres verbes :} \textbf{\all{würde} + Infinitif} (forme standard, évite les formes simples archaïques comme \all{ich käme, ich gäbe}).
\cle{Marqueurs de souhait :} \all{Ich wünschte, ...} (+\all{dass} ou sans \all{dass} $\rightarrow$ verbe fin), \all{Wenn ich nur ... könnte!}, \all{Hätte ich doch ...!} (inversion sans \all{wenn}).
\end{propriete}

\begin{exemplebox}[Exemples du texte et variations]
\begin{itemize}
    \item \all{Ich wünschte, die Gesellschaft \textbf{wäre} offener.} (\all{wäre} = KII \all{sein}, 3e pers. sg., fin de subordonnée).
    \item \all{Man \textbf{sollte} die Vielfalt \textbf{sehen}.} (\all{sollte} = KII \all{sollen}, 3e pers. sg. \all{man} + infinitif \all{sehen} à la fin).
    \item \all{Ich \textbf{wäre} gerne zweisprachig.} (\all{wäre} + adjectif).
    \item \all{Wenn ich \textbf{könnte}, \textbf{würde} ich sorbisch lernen.} (Conditionnel hypothétique).
\end{itemize}
\end{exemplebox}

\begin{attention}[Piège majeur : \textbf{\all{würde} vs \all{wurde}}]
\begin{itemize}
    \item \all{Ich \textbf{würde} gerne helfen.} (KII de \all{werden} = « je voudrais / je ferais »).
    \item \all{Ich \textbf{wurde} gestern krank.} (Präteritum de \all{werden} = « je suis devenu / j'ai été » malade). \textbf{Ne jamais écrire *\all{Ich würde krank} pour « je suis devenu malade » !}
\end{itemize}
\end{attention}

\begin{exercice}[Entraînement Grammaire 2 — Transformation]
Transformez selon la contrainte.
\begin{enumerate}
    \item Ich habe Rechte. Man hört mich nicht. (\textit{Concession avec \cle{obwohl}})
    \item Die Sprache ist wichtig. Man verbietet sie. (\textit{Concession avec \cle{trotzdem}})
    \item Die Gesellschaft ist intolerant. (\textit{Souhait : \cle{Ich wünschte, ... wäre ...}})
    \item Man respektiert die Minderheiten nicht. (\textit{Recommandation KII : \cle{Man sollte ...}})
    \item Ich lerne die Sprache. (\textit{Souhait d'habileté : \cle{Ich könnte ...}})
\end{enumerate}
\end{exercice}

\begin{corrige}[Exercice Grammaire 2]
\begin{enumerate}
    \item \all{Obwohl ich Rechte habe, hört man mich nicht.}
    \item \all{Die Sprache ist wichtig, trotzdem verbietet man sie.} (Ou : \all{... . Trotzdem verbietet man sie.})
    \item \all{Ich wünschte, die Gesellschaft wäre toleranter.}
    \item \all{Man sollte die Minderheiten respektieren.}
    \item \all{Ich könnte die Sprache lernen.}
\end{enumerate}
\end{corrige}

\courssub{Intégration : Production guidée — « Ein Leserbrief für Toleranz »}
\courssub{Objectifs de l'activité}
Cette activité d'intégration vise à mobiliser l'ensemble des acquis de la leçon AL04 dans une tâche complexe de production écrite et orale, conforme aux exigences du programme de Terminale A (LV2). Les élèves devront :
\begin{itemize}
    \item \textbf{Réinvestir le lexique thématique} (\cle{die Minderheit}, \cle{die Diskriminierung}, \cle{die Integration}, \cle{das Vorurteil}, \cle{die Sprache}, \cle{die Rechte}, \cle{die Identität}, \cle{benachteiligen}, \cle{akzeptieren}, \cle{der Respekt}) de manière pertinente et correcte (genre, pluriel).
    \item \textbf{Appliquer la déclinaison de l'adjectif} (forte, faible, mixte) dans des groupes nominaux variés (sujet, objet, prépositionnel) sans erreur d'orthographe ni de finale.
    \item \textbf{Maîtriser les structures de la concession} : subordonnée avec \all{obwohl} (verbe en fin) et parataxe avec \all{trotzdem} (place 1 ou 3), pour nuancer l'argumentation.
    \item \textbf{Utiliser le subjonctif II (Konjunktiv II)} pour formuler des souhaits irréels du présent (\all{ich wünschte, ich hätte, ich wäre, ich könnte, man sollte}) dans une visée persuasive.
    \item \textbf{Produire un texte cohérent et authentique} (lettre au courrier des lecteurs, 10--12 lignes), respectant les conventions du genre (salutation, objet, corps, formule de politesse, signature).
    \item \textbf{Développer des compétences transversales} : travail collaboratif (négociation du sens, relecture par les pairs), expression orale (lecture expressive, débat), esprit critique (vote argumenté).
\end{itemize}

\courssub{Situation}
\begin{textebox}[Contexte : Appel à contributions du « Kamerun-Kurier »]
\large \textbf{Der Kamerun-Kurier – Ihre Stimme für Vielfalt und Toleranz}

\medskip
\noindent
Liebe Leserinnen und Leser,

Kamerun ist ein Land der Vielfalt: über 250 Ethnien, zwei offizielle Amtssprachen (Französisch und Englisch) und unzählige lokale Sprachen prägen unseren Alltag. Doch Vielfalt bedeutet nicht automatisch Gleichberechtigung. Viele junge Menschen, die zu sprachlichen, kulturellen oder sozialen Minderheiten gehören, erleben täglich Vorurteile, Diskriminierung oder Ausgrenzung – in der Schule, auf dem Arbeitsmarkt, im Viertel oder sogar in der eigenen Familie.

Wir möchten Ihnen eine Plattform bieten. Schreiben Sie uns einen Leserbrief (10--12 Zeilen)! Schildern Sie eine konkrete Situation, in der Sie sich als Angehörige(r) einer Minderheit benachteiligt gefühlt haben. Nennen Sie das Problem beim Namen, fordern Sie Ihre Rechte ein und formulieren Sie zwei konkrete Wünsche für eine bessere Zukunft.

Die besten Zuschriften veröffentlichen wir in unserer Sonderausgabe zum „Tag der Minderheiten“ und die Gewinnergruppe erhält ein Praktikum in unserer Redaktion in Yaoundé.

Einsendeschluss: Freitag, 15. November. Adresse: redaktion@kameroon-kurier.cm

Mit solidarischen Grüßen

\emph{Die Redaktion des Kamerun-Kuriers}
\end{textebox}

\noindent
\textbf{Mise en situation pour la classe :} Vous êtes élèves en Terminale A dans un lycée de Yaoundé (Lycée Général Leclerc, Collège de la Retraite, etc.). Vous maîtrisez l'allemand (LV2) et vous vous sentez concernés par cette thématique. En groupes de 3 à 4, vous décidez de participer au concours. Vous devez choisir un profil de « jeune appartenant à une minorité au Cameroun » (ex. : jeune anglophone à Yaoundé, jeune Baka de l'Est, jeune Mbororo de l'Adamaoua, jeune germanophone dans un environnement francophone, jeune en situation de handicap, etc.) et rédiger collectivement cette lettre.

\courssub{Consignes}
Les consignes sont conçues pour guider pas à pas le travail de groupe, de la compréhension de la consigne à la production finale, en passant par la planification linguistique.

\begin{enumerate}
    \item \textbf{Analyse de la situation de communication (10 min).}
    \begin{itemize}
        \item Lisez le texte d'appel (ci-dessus) dans la \texttt{textebox}. Répondez collectivement en allemand : \all{Wer schreibt an wen?} (Qui écrit à qui ?), \all{Was ist das Thema?} (Quel est le thème ?), \all{Welche Textsorte?} (Quel genre textuel ?), \all{Welche formalen Vorgaben gibt es?} (Quelles contraintes formelles ? : longueur, date, adresse, 2 souhaits, etc.).
        \item Choisissez le profil de votre « persona » (identité, minorité, lieu de vie, problème concret). Notez 3 mots-clés en allemand pour ce profil.
    \end{itemize}
    
    \item \textbf{Planification linguistique \& Recherche lexicale (15 min).}
    \begin{itemize}
        \item Sur une feuille de brouillon, dressez la liste des \textbf{6 mots minimum du lexique de la leçon} que vous allez imposer dans votre texte. Vérifiez l'article et le pluriel (cf. \texttt{Ressources à mobiliser}).
        \item Construisez \textbf{2 groupes nominaux avec adjectifs déclinés} (différents cas : nominatif, accusatif, datif ; différentes déclinaisons : forte, faible, mixte). Ex. : \all{ein junger Mann} (faible, nom. — \textit{correction: mixte/forte}), \all{keine faire Chancen} (forte, acc. — \textit{correction: \all{keine fairen Chancen}}), \all{mit altem Vorurteil} (mixte, dat. — \textit{correction: \all{mit altem Vorurteil} forte ou \all{mit dem alten Vorurteil} faible}). Écrivez-les à part pour les vérifier.
        \item Rédigez \textbf{1 phrase avec \cle{obwohl}} (subordonnée, verbe à la fin) et \textbf{1 phrase avec \cle{trotzdem}} (conjonction/adverbe, verbe en 2e position). Ex. : \all{Obwohl ich gute Noten habe, werde ich abgelehnt.} / \all{Ich habe gute Noten, werde trotzdem abgelehnt.}
        \item Formulez \textbf{2 souhaits au Subjonctif II} (\cle{Konjunktiv II}) : un avec \all{wäre/hätte}, un avec \all{könnte/sollte/würde + infinitif}. Ex. : \all{Ich wünschte, ich wäre akzeptiert.} / \all{Man sollte mehr Toleranz zeigen.}
    \end{itemize}
    
    \item \textbf{Rédaction collaborative de la lettre (25 min).}
    \begin{itemize}
        \item Rédigez la version finale sur une grande feuille (format A3 ou paperboard) pour lisibilité lors de la lecture.
        \item Structure obligatoire :
        \begin{enumerate}
            \item Salutation : \all{Sehr geehrte Damen und Herren,}
            \item Objet : \all{Betreff: Leserbrief zum Thema „Minderheiten in Kamerun“}
            \item Introduction : Présentation du « je » (qui ? où ? quelle minorité ?).
            \item Développement : Description du problème concret (préjugé, discrimination), utilisation de \cle{obwohl} / \cle{trotzdem}.
            \item Revendication : Droits, intégration, respect.
            \item Deux souhaits (Subjonctif II).
            \item Conclusion + Formule de politesse : \all{Mit freundlichen Grüßen} + Prénom + Ville.
        \end{enumerate}
        \item \textbf{Autocontrôle (grille) :} Cochez chaque contrainte une fois remplie (6 mots vocab, 2 adj. déclinés, 1 \cle{obwohl}, 1 \cle{trotzdem}, 2 KII, 10-12 lignes).
    \end{itemize}
    
    \item \textbf{Relecture croisée \& Correction (10 min).}
    \begin{itemize}
        \item Échangez la feuille avec un autre groupe. Lisez la lettre à voix basse. Vérifiez : orthographe (majuscules noms, Umlaute, ß), déclinaisons (tableau de référence), ordre des mots (obwohl = verbe fin, trotzdem = verbe 2e), formation KII (würde + infinitif / forme simple). Annotez les erreurs suspectes au crayon.
        \item Récupérez votre production, corrigez les erreurs validées.
    \end{itemize}
    
    \item \textbf{Présentation orale \& Vote (15 min).}
    \begin{itemize}
        \item Un porte-parole par groupe lit la lettre avec l'intonation appropriée (ton engagé, respectueux).
        \item La classe note sur un bulletin : \textit{Contenu (convaincant ?) / Langue (correct ?) / Respect contraintes / Originalité}.
        \item Dépouillement et annonce de la lettre gagnante (simulation publication / stage).
    \end{itemize}
\end{enumerate}

\courssub{Ressources à mobiliser}
Cette section regroupe les outils grammaticaux, lexicaux et méthodologiques vus dans la leçon. Les élèves doivent les avoir sous les yeux (affichage au tableau, fiche distribuée ou manuel ouvert).

\begin{propriete}[Rappel : La déclinaison de l'adjectif (Terminale A)]
L'adjectif attribut (épithète) se décline selon le \textbf{déterminant qui le précède} (article défini, indéfini, possessif, négatif, ou zéro article). Trois types de déclinaison : 
\begin{center}
\begin{tabularx}{\linewidth}{|X|X|X|X|}
\hline
\rowcolor{popblueL} \textbf{Cas} & \textbf{Faible (après art. défini \textit{der, die, das, die})} & \textbf{Forte (zéro article / après \textit{kein, mein, dieser} au n. s. m./n.)} & \textbf{Mixte (après art. indéfini \textit{ein, eine, kein, mein, dein} + poss. au n. s. m./n.)} \\ \hline
\textbf{Nom. m.} & der \textbf{gut\emph{e}} Mann & \textbf{gut\emph{er}} Mann & ein \textbf{gut\emph{er}} Mann \\ \hline
\textbf{Acc. m.} & den \textbf{gut\emph{en}} Mann & \textbf{gut\emph{en}} Mann & einen \textbf{gut\emph{en}} Mann \\ \hline
\textbf{Dat. m.} & dem \textbf{gut\emph{en}} Mann & \textbf{gut\emph{em}} Mann & einem \textbf{gut\emph{en}} Mann \\ \hline
\textbf{Gen. m.} & des \textbf{gut\emph{en}} Mannes & \textbf{gut\emph{en}} Mannes & eines \textbf{gut\emph{en}} Mannes \\ \hline
\textbf{Nom. f.} & die \textbf{gut\emph{e}} Frau & \textbf{gut\emph{e}} Frau & eine \textbf{gut\emph{e}} Frau \\ \hline
\textbf{Acc. f.} & die \textbf{gut\emph{e}} Frau & \textbf{gut\emph{e}} Frau & eine \textbf{gut\emph{e}} Frau \\ \hline
\textbf{Dat. f.} & der \textbf{gut\emph{en}} Frau & \textbf{gut\emph{er}} Frau & einer \textbf{gut\emph{en}} Frau \\ \hline
\textbf{Gen. f.} & der \textbf{gut\emph{en}} Frau & \textbf{gut\emph{er}} Frau & einer \textbf{gut\emph{en}} Frau \\ \hline
\textbf{Nom. n.} & das \textbf{gut\emph{e}} Kind & \textbf{gut\emph{es}} Kind & ein \textbf{gut\emph{es}} Kind \\ \hline
\textbf{Acc. n.} & das \textbf{gut\emph{e}} Kind & \textbf{gut\emph{es}} Kind & ein \textbf{gut\emph{es}} Kind \\ \hline
\textbf{Dat. n.} & dem \textbf{gut\emph{en}} Kind & \textbf{gut\emph{em}} Kind & einem \textbf{gut\emph{en}} Kind \\ \hline
\textbf{Gen. n.} & des \textbf{gut\emph{en}} Kindes & \textbf{gut\emph{en}} Kindes & eines \textbf{gut\emph{en}} Kindes \\ \hline
\textbf{Pluriel (tous genres)} & die \textbf{gut\emph{en}} Leute & \textbf{gut\emph{e}} Leute & keine \textbf{gut\emph{en}} Leute \\ \hline
\end{tabularx}
\end{center}
\cle{Attention francophones :} L'adjectif \textbf{s'accorde toujours} en genre, nombre et cas avec le nom qu'il qualifie, même si le déterminant porte déjà la marque (ex. \all{dem gut\emph{en}} Mann, pas *dem gut Mann). Le \textbf{e muet} final (Schwa) s'écrit \textbf{-e} à tous les cas du féminin, du pluriel, et au nominatif/accusatif neutre (faible/mixte) ; il devient \textbf{-en} aux autres cas (masc. acc./dat./gen., neutre dat./gen., pluriel toutes déclinaisons).
\end{propriete}

\begin{propriete}[Rappel : La concession — \cle{obwohl} vs \cle{trotzdem}]
\begin{center}
\begin{tabularx}{\linewidth}{|X|X|}
\hline
\rowcolor{poporangeL} \textbf{Structure} & \textbf{Règle de place du verbe} \\ \hline
\all{obwohl} (subordonnante) & Introduit une \textbf{proposition subordonnée}. Le verbe conjugué \textbf{va à la fin} de la proposition. \\ 
\all{Beispiel:} \all{Obwohl ich die Sprache gut spreche, werde ich nicht verstanden.} & \all{..., \textbf{spreche}} (fin). La proposition principale qui suit commence par le verbe (V2) : \all{werde ich ...} \\ \hline
\all{trotzdem} (adverbe / conjonction de coordination) & Relie deux \textbf{propositions principales} (ou commence une phrase). Le verbe reste en \textbf{2e position} (place 1 = \cle{trotzdem} ou sujet). \\ 
\all{Beispiel:} \all{Ich spreche die Sprache gut, trotzdem werde ich nicht verstanden.} & \all{trotzdem} = place 1 $\rightarrow$ \all{werde} (V2). Ou : \all{Ich spreche die Sprache gut. Trotzdem werde ich nicht verstanden.} \\ \hline
\end{tabularx}
\end{center}
\cle{Sens :} Les deux expriment une opposition, une concession (« bien que / malgré cela »). \all{Obwohl} insiste sur le fait concessionnel ; \all{trotzdem} insiste sur la conséquence/mainien du fait principal.
\end{propriete}

\begin{propriete}[Rappel : Subjonctif II (Konjunktiv II) — Wunsch (Souhait irréel du présent)]
Pour exprimer un souhait non réalisé dans le présent, on utilise le \textbf{Konjunktiv II}.
\begin{itemize}
    \item \textbf{Verbes auxiliaires \& modaux (formes simples, à connaître par cœur) :}
    \begin{center}
    \begin{tabular}{|l|l|l|l|l|l|l|}
    \hline
    \rowcolor{popgreenL} \textbf{Infinitif} & \textbf{ich} & \textbf{du} & \textbf{er/sie/es} & \textbf{wir} & \textbf{ihr} & \textbf{sie/Sie} \\ \hline
    \all{sein} & \all{wäre} & \all{wärest} & \all{wäre} & \all{wären} & \all{wäret} & \all{wären} \\ \hline
    \all{haben} & \all{hätte} & \all{hättest} & \all{hätte} & \all{hätten} & \all{hättet} & \all{hätten} \\ \hline
    \all{können} & \all{könnte} & \all{könntest} & \all{könnte} & \all{könnten} & \all{könntet} & \all{könnten} \\ \hline
    \all{sollen} & \all{sollte} & \all{solltest} & \all{sollte} & \all{sollten} & \all{solltet} & \all{sollten} \\ \hline
    \all{wollen} & \all{wollte} & \all{wolltest} & \all{wollte} & \all{wollten} & \all{wolltet} & \all{wollten} \\ \hline
    \all{müssen} & \all{müsste} & \all{müsstest} & \all{müsste} & \all{müssten} & \all{müsstet} & \all{müssten} \\ \hline
    \all{dürfen} & \all{dürfte} & \all{dürftest} & \all{dürfte} & \all{dürften} & \all{dürftet} & \all{dürften} \\ \hline
    \all{werden} & \all{würde} & \all{würdest} & \all{würde} & \all{würden} & \all{würdet} & \all{würden} \\ \hline
    \end{tabular}
    \end{center}
    \item \textbf{Autres verbes :} Forme \textbf{würde + infinitif} (standard à l'oral et à l'écrit pour la plupart des verbes). \all{Ich würde gerne helfen.} \all{Man würde es besser machen.}
    \item \textbf{Marqueurs de souhait :} \all{Ich wünschte, ...} (subordonnée avec verbe au KII), \all{Wenn ich nur ... könnte!}, \all{Hätte ich doch ...!} (inversion sujet-verbe sans \all{wenn}).
\end{itemize}
\cle{Piège :} Ne pas confondre \all{würde} (KII de \all{werden}) et \all{wurde} (Präteritum de \all{werden} = « est devenu / a été »). \all{Ich würde} (je voudrais/ferais) vs \all{Ich wurde} (je suis devenu / j'ai été).
\end{propriete}

\begin{definition}[Lexique thématique obligatoire — Liste de référence (Article + Pluriel)]
\begin{center}
\begin{tabularx}{\linewidth}{|l|l|X|}
\hline
\rowcolor{poppurpleL} \textbf{Deutsch (Art. + Pl.)} & \textbf{Français} & \textbf{Famille de mots / Collocations utiles} \\ \hline
\all{die Minderheit, -en} & la minorité & \all{angehören} (faire partie), \all{ethnische / sprachliche / religiöse Minderheit}, \all{Minderheitenrechte} \\ \hline
\all{die Diskriminierung, -en} & la discrimination & \all{diskriminieren} (verbe), \all{benachteiligen} (désavantager), \all{Opfer der Diskriminierung sein} \\ \hline
\all{die Integration, -en} & l'intégration & \all{integrieren} (verbe), \all{Integrationspolitik}, \all{gelungene / gescheiterte Integration} \\ \hline
\all{das Vorurteil, -e} & le préjugé & \all{Vorurteile haben / abbauen / bekämpfen}, \all{vorurteilsfrei}, \all{klischeehaft} \\ \hline
\all{die Sprache, -n} & la langue & \all{Muttersprache}, \all{Amts- / Verkehrssprache}, \all{Sprachbarriere}, \all{mehrsprachig} \\ \hline
\all{das Recht, -e} & le droit & \all{Rechte haben / fordern / einfordern}, \all{Grundrechte}, \all{Menschenrechte}, \all{gleichberechtigt} \\ \hline
\all{der Respekt, -e} & le respect & \all{Respekt vor jdm. (Dat) haben}, \all{respektvoll}, \all{Anerkennung} \\ \hline
\all{die Toleranz, -en} & la tolérance & \all{tolerant}, \all{Intoleranz}, \all{Toleranz zeigen / fordern} \\ \hline
\all{die Ausgrenzung, -en} & l'exclusion & \all{ausgrenzen}, \all{ausgegrenzt sein}, \all{sozial ausgeschlossen} \\ \hline
\all{die Identität, -en} & l'identité & \all{Identitätsverlust}, \all{kulturelle Identität bewahren} \\ \hline
\end{tabularx}
\end{center}
\cle{Conseil :} Cochez les 6 mots que vous intégrez. Vérifiez l'accord de l'adjectif épithète qui les accompagne (ex. \all{ethnische Minderheit} f. nom. sg. faible ; \all{keine faire Chancen} $\rightarrow$ \all{keine \textbf{fairen} Chancen} f. pl. acc. forte).
\end{definition}

\begin{methode}[Rédiger un Leserbrief (Lettre au courrier des lecteurs) — Structure type]
\begin{enumerate}
    \item \textbf{Kopf} : Expéditeur (optionnel ici), Destinataire (\all{An die Redaktion des Kamerun-Kuriers}), Date, Lieu (\all{Yaoundé, den 10. November 2025}).
    \item \textbf{Betreffzeile} : Clair, précis (\all{Betreff: Leserbrief: Meine Erfahrung als ...}).
    \item \textbf{Anrede} : \all{Sehr geehrte Damen und Herren,} (formel standard).
    \item \textbf{Einleitung} : « Je » + situation + minorité concernée. Accroche.
    \item \textbf{Hauptteil} : Argumentation (Problème concret $\rightarrow$ \cle{obwohl}/\cle{trotzdem} $\rightarrow$ Revendication droits). \textbf{Contraintes grammaticales ici.}
    \item \textbf{Schluss / Appell} : Les \textbf{2 souhaits (KII)} + appel à l'action collective.
    \item \textbf{Grußformel} : \all{Mit freundlichen Grüßen} (standard formel).
    \item \textbf{Unterschrift} : Prénom + Ville/Quartier (\all{Amadou, Yaoundé-Mvog-Ada}).
\end{enumerate}
\end{methode}

\begin{popfigure}
\begin{tikzpicture}[
    mindmap, 
    grow cyclic, 
    every node/.style={concept, execute at begin node=\hskip0pt}, 
    root concept/.append style={concept color=popblue, font=\large\bfseries, text width=3.5cm, align=center},
    level 1 concept/.append style={concept color=popblue!70!popgreen, font=\normalsize\bfseries, text width=2.8cm, align=center, sibling angle=90, level distance=4.5cm},
    level 2 concept/.append style={concept color=poporange!80!poppink, font=\small, text width=2.2cm, align=center, sibling angle=45, level distance=3.5cm},
    text=white
]
\node[align=center]{Leserbrief\\Checkliste}
    child [clockwise from=90] {node[align=center]{Lexique\\ (6+ mots)}
        child {node {Minderheit}}
        child {node {Diskriminierung}}
        child {node {Integration}}
        child {node {Vorurteil}}
        child {node {Rechte}}
        child {node {Sprache}}
    }
    child [clockwise from=0] {node[align=center]{Grammaire\\ (Contraintes)}
        child {node {2 Adj. déclinés}}
        child {node {1 $\times$ obwohl}}
        child {node {1 $\times$ trotzdem}}
        child {node {2 $\times$ KII Wunsch}}
    }
    child [clockwise from=-90] {node[align=center]{Structure\\ (Forme)}
        child {node {Betreff}}
        child {node {Anrede}}
        child {node {Ich-Perspektive}}
        child {node {Argumentation}}
        child {node {Appell}}
        child {node {Grußformel}}
    }
    child [clockwise from=180] {node[align=center]{Stratégie\\ (Groupe)}
        child {node {Persona choix}}
        child {node {Brouillon commun}}
        child {node {Relecture croisée}}
        child {node {Lecture expressive}}
        child {node {Vote classe}}
    };
\end{tikzpicture}
\legende{Carte mentale : Check-list de l'activité « Ein Leserbrief für Toleranz »}
\end{popfigure}

\courssub{Proposition de production et corrigé commenté}
\begin{exoresolu}[Modèle de Leserbrief — Profil : Jeune Baka (Est-Cameroun)]
\textbf{Énoncé :} Rédiger la lettre au « Kamerun-Kurier » (10--12 lignes) en respectant toutes les contraintes.

\tcblower

\textbf{Production modèle (allemand) :}

\medskip
\noindent
\textbf{An die Redaktion des Kamerun-Kuriers} \\
\textbf{Yaoundé, den 12. November 2025} \\
\textbf{Betreff: Leserbrief – Meine Stimme als junge Baka-Frau} \\

\medskip
\noindent
Sehr geehrte Damen und Herren,

\medskip
\noindent
ich heiße Mireille, bin 19 Jahre alt und gehöre zur \textbf{ethnischen Minderheit} der Baka im Osten Kameruns. \textbf{Obwohl} ich das Abitur mit Erfolg \textbf{bestanden habe}, werde ich bei Bewerbungen oft \textbf{wegen meiner Herkunft} abgelehnt. Man hat \textbf{Vorurteile}: „Baka können nicht lernen“, hört man. \textbf{Trotzdem} kämpfe ich um meinen Platz an der Universität. \textbf{Diskriminierung} verletzt meine \textbf{Rechte} und zerstört \textbf{Integration}. \textbf{Ich wünschte, die Gesellschaft wäre toleranter.} \textbf{Man sollte} die \textbf{kulturelle Identität} jedes Volkes \textbf{respektieren} und fördern.

\medskip
\noindent
Mit freundlichen Grüßen \\
Mireille, Yokadouma

\medskip
\textbf{Commentaire détaillé des choix linguistiques (Corrigé commenté) :}

\begin{enumerate}
    \item \textbf{Lexique thématique (7 mots validés) :}
    \begin{itemize}
        \item \all{die Minderheit} (l. 5 : \all{zur ethnischen Minderheit} — \textbf{Faible, Datif, Féminin, Singulier} : \all{zur} = \all{zu der} + \all{ethnisch\emph{en}} + \all{Minderheit}).
        \item \all{das Vorurteil} (l. 7 : \all{Vorurteile} — \textbf{Pluriel Accusatif} après \all{hat}).
        \item \all{die Diskriminierung} (l. 8 : \all{Diskriminierung} — \textbf{Nominatif Féminin Singulier} sujet de \all{verletzt}).
        \item \all{die Rechte} (l. 8 : \all{meine Rechte} — \textbf{Accusatif Pluriel} poss. \all{meine} + nom \all{Rechte}).
        \item \all{die Integration} (l. 8 : \all{Integration} — \textbf{Accusatif F. Sg.} objet de \all{zerstört}).
        \item \all{die Identität} (l. 10 : \all{die kulturelle Identität} — \textbf{Accusatif F. Sg.} \all{die kulturell\emph{e} Identität} — \textbf{Faible} article défini \all{die} $\rightarrow$ \textbf{-e}). 
        \item \all{der Respekt} (l. 10 : verbe \all{respektieren} utilisé). \cle{Total : 7 items lexicaux distincts. Contrainte respectée.}
    \end{itemize}
    
    \item \textbf{Déclinaison de l'adjectif (3 occurrences validées, types variés) :}
    \begin{itemize}
        \item \textbf{Occurrence 1 (Faible, Datif, Féminin, Singulier) :} \all{zur \textbf{ethnisch\emph{en}} Minderheit}. Analyse : \all{zur} = \all{zu} (préposition datif) + \all{der} (art. déf. f. dat.). Déclinaison \textbf{faible} $\rightarrow$ \textbf{-en}. Correct.
        \item \textbf{Occurrence 2 (Faible, Accusatif, Féminin, Singulier) :} \all{die \textbf{kulturell\emph{e}} Identität} (l. 10). Article défini \all{die} $\rightarrow$ faible $\rightarrow$ \textbf{-e}. Correct.
        \item \textbf{Occurrence 3 (Forte, Datif, Féminin, Singulier - implicite dans "mit großer Freude" du texte support, ici ajouté pour l'exemple) :} \all{mit \textbf{groß\emph{er}} Freude}. Zéro article $\rightarrow$ Forte fém. dat. $\rightarrow$ \textbf{-er}. L'élève doit en présenter 2 distincts et corrects (ex: 1 faible, 1 forte ou 1 mixte).
    \end{itemize}
    
    \item \textbf{Concession (2 structures distinctes) :}
    \begin{itemize}
        \item \all{\textbf{Obwohl} ich das Abitur \textbf{bestanden habe}, ...} (l. 6). \textbf{Subordonnée} : verbe \all{habe} en \textbf{position finale}. Participe \all{bestanden} avant. Structure parfaite.
        \item \all{Man hat Vorurteile. \textbf{Trotzdem} kämpfe ich ...} (l. 7-8). \textbf{Proposition principale} : \all{Trotzdem} en \textbf{position 1}, verbe \all{kämpfe} en \textbf{position 2 (V2)}. Sujet \all{ich} en position 3. Correct. \cle{Alternative acceptée :} \all{Ich habe Vorurteile, trotzdem kämpfe ich ...} (\all{trotzdem} pos 3, verbe V2 \all{kämpfe}).
    \end{itemize}
    
    \item \textbf{Subjonctif II — Souhaits (2 occurrences) :}
    \begin{itemize}
        \item \all{\textbf{Ich wünschte, die Gesellschaft \textbf{wäre} toleranter.}} (l. 9). 
        \begin{itemize}
            \item \all{wünschte} = KII de \all{wünschen} (forme simple Präteritum/KII).
            \item \all{wäre} = KII de \all{sein} (forme simple, 3e pers. sg.). 
            \item Structure \all{ich wünschte,} [Hauptsatz], \all{die Gesellschaft wäre...} [Nebensatz sans \all{dass}, verbe \all{wäre} à la fin]. Correct.
        \end{itemize}
        \item \all{\textbf{Man sollte} die kulturelle Identität ... \textbf{respektieren} und fördern.} (l. 10). 
        \begin{itemize}
            \item \all{sollte} = KII de \all{sollen} (forme simple, 3e pers. sg. \all{man}).
            \item \all{respektieren / fördern} = infinitifs en position finale. 
            \item Sens : « On devrait respecter... » (recommandation morale / souhait de changement). Correct.
        \end{itemize}
    \end{itemize}
    
    \item \textbf{Orthographe, Ponctuation, Majuscules :} Tous les noms (\all{Minderheit, Abitur, Herkunft, Vorurteile, Diskriminierung, Rechte, Integration, Gesellschaft, Identität, Volk}) portent une majuscule. Guillemets allemands \all{„ „} utilisés pour la citation orale. Umlauts (\all{ü, ö, ä}) et \all{ß} (\all{Grüßen}) respectés.
    
    \item \textbf{Registre \& Cohérence :} Ton formel (\all{Sehr geehrte Damen und Herren}, \all{Mit freundlichen Grüßen}). Perspective subjective (\all{ich}) légitime. Argumentation : fait $\rightarrow$ concession $\rightarrow$ revendication $\rightarrow$ souhaits. Longueur : ~11 lignes manuscrites.
\end{enumerate}

\cle{Note pour l'enseignant :} Ce modèle est une \textbf{proposition de référence}. Les productions des élèves varieront selon le persona choisi (anglophone, handicapé, germaniste, etc.). L'évaluation portera sur le \textbf{respect strict des 6 contraintes formelles} (check-list) et la \textbf{correction morphosyntaxique} (déclinaisons, ordre mots, KII), plus que sur l'originalité du contenu, bien que celle-ci soit valorisée au vote.
\end{exoresolu}

\courssub{Bilan de l'activité}
\begin{aretenir}[Bilan compétences — Leçon AL04 : Intégration]
À l'issue de cette activité, l'élève de Terminale A doit être capable de :
\begin{itemize}
    \item \textbf{Comprendre} un appel à contributions journalistique authentique (genre \textit{Leserbriefaufruf}) et en extraire les contraintes formelles et thématiques.
    \item \textbf{Mobiliser} un lexique spécialisé (\cle{Minderheit, Diskriminierung, Integration, Vorurteil, Rechte, Sprache, Identität, Toleranz, Ausgrenzung, Respekt}) avec articles, pluriels et collocations correctes dans un texte personnel.
    \item \textbf{Contrôler} la déclinaison de l'adjectif (faible, forte, mixte) dans au moins deux groupes nominaux distincts (différents cas/genres), en justifiant son choix par l'analyse du déterminant.
    \item \textbf{Articuler} une argumentation nuancée grâce aux outils de la concession (\cle{obwohl} + verbe fin / \cle{trotzdem} + V2), en distinguant clairement subordination et coordination.
    \item \textbf{Exprimer} l'irréel du présent (souhait, recommandation) via le \cle{Konjunktiv II} (formes simples \all{wäre/hätte/sollte/könnte} ou \all{würde + infinitif}), en évitant la confusion avec le Präteritum (\all{wurde}).
    \item \textbf{Rédiger} un texte cohérent, formaté (Leserbrief), de 10-12 lignes, en respectant les conventions sociolinguistiques (formules d'appel/de politesse, objet, signature).
    \item \textbf{Travailler en équipe} (négociation du sens, relecture par les pairs, répartition des rôles : scribe, correcteur, porte-parole) et \textbf{s'exprimer en public} (lecture expressive, écoute critique, vote argumenté).
\end{itemize}
\end{aretenir}

\begin{attention}[Erreurs fréquentes observées lors de l'activité — Pistes de remédiation]
\begin{itemize}
    \item \textbf{Déclinaison :} Oubli du \textbf{-n} au datif pluriel (\all{mit jung\emph{en}} Leuten, pas *\all{jung\emph{e} Leuten}). Confusion forte/mixte au masculin nominatif/accusatif (\all{ein gut\emph{er}} Mann vs \all{kein gut\emph{er}} Mann — les deux sont \textbf{-er} car article "nu" $\rightarrow$ forte ; \all{den gut\emph{en}} Mann acc. faible). \rightarrow \textbf{Exercice ciblé :} « Declension-Sprint » sur 10 GN imposés.
    \item \textbf{Obwohl / Trotzdem :} Inversion du verbe (\all{*Obwohl ich habe gute Noten} ; \all{*Trotzdem ich kämpfe}). \rightarrow \textbf{Astuce mnémotechnique :} \all{Obwohl} = « Verbe à la fin » (queue) ; \all{Trotzdem} = « Verbe en 2 » (deux doigts).
    \item \textbf{Konjunktiv II :} \all{*Ich würde sein} au lieu de \all{Ich wäre} ; \all{*Man wurde} (passé) au lieu de \all{Man würde} / \all{Man sollte}. \rightarrow \textbf{Tableau de référence KII} affiché en permanence ; drilling oral \all{ich wäre / ich hätte / ich könnte / ich sollte / ich würde}.
    \item \textbf{Lexique :} Confusion \all{die Minderheit} (sg) / \all{die Minderheiten} (pl) ; \all{das Vorurteil} (n) / \all{die Vorurteile} (pl) ; \all{die Diskriminierung} (f) vs \all{diskriminieren} (v). \rightarrow \textbf{Fiche vocabulaire} à apprendre par cœur (article + pluriel + 1 collocation).
    \item \textbf{Majuscules :} Noms communs écrits en minuscule (*\all{die sprache, die rechte}). \rightarrow \textbf{Règle d'or :} « En allemand, \textbf{tout nom} a une majuscule. Point. »
\end{itemize}
\end{attention}

\begin{savaistu}[Contexte culturel : Minorités au Cameroun \& en Allemagne]
\textbf{Cameroun :} La Constitution reconnaît le bilinguisme officiel (Français/Anglais). Les minorités « autochtones » (Baka/Bagyeli au Sud-Est, Mbororo/Fulani nomades au Nord/Adamaoua, Montagnards de l'Ouest) font l'objet de politiques d'intégration spécifiques (MINAS, ONG). Le problème de l'exclusion linguistique (anglophones vs francophones) est au cœur de la crise du NOSO (North-West/South-West) depuis 2016. L'allemand, langue d'enseignement (LV2), crée une « minorité germanophone » scolaire dynamique (clubs, jumelages, DAAD, Goethe-Institut Yaoundé/Douala).

\textbf{Allemagne / Autriche / Suisse (DACH) :} Minorités nationales reconnues : Danois (Schleswig-Holstein), Sorabes/Wendes (Lusace/Brandebourg-Saxe), Frisons (Côte Nord), Roms/Sintis, Allemands de langue des signes. Loi-cadre sur les minorités (2005). En Suisse : 4 langues nationales (allemand, français, italien, romanche) — le romanche est langue officielle partielle (minorité). Autriche : Minorités slovène, croate, hongroise, tchèque, slovaque, rom, langue des signes (loi 1976/2000). Le concept de \all{Leitkultur} (culture directrice) vs \all{Multikulturalismus} anime le débat public.
\end{savaistu}

\begin{exercice}[Entraînement individuel — Consolidation post-activité]
\textbf{Transformez les phrases suivantes en intégrant la contrainte indiquée entre parenthèses. Écrivez la phrase complète en allemand.}
\begin{enumerate}
    \item Ich bin eine gute Schülerin. Man hört mich nicht. (\textit{Concession avec \cle{obwohl}})
    \item Er hat ein starkes Vorurteil. Er kennt mich nicht. (\textit{Concession avec \cle{trotzdem}})
    \item Die Gesellschaft ist intolerant. (\textit{Souhait KII : \cle{Ich wünschte, ...}})
    \item Man respektiert die Rechte der Minderheiten nicht. (\textit{Souhait KII avec modal : \cle{Man sollte ...}})
    \item Ich schreibe einen Brief an die Redaktion. (\textit{Ajoutez un adjectif décliné \textbf{mixte accusatif masc.} : \textit{einen ... Brief}})
    \item Wir brauchen mehr Toleranz. (\textit{Ajoutez un adjectif décliné \textbf{faible nominatif fém.} : \textit{die ... Toleranz}})
\end{enumerate}
\end{exercice}

\begin{corrige}[Exercice consolidation]
\begin{enumerate}
    \item \all{Obwohl ich eine gute Schülerin bin, hört man mich nicht.} (Verbe \all{bin} à la fin).
    \item \all{Er hat ein starkes Vorurteil, trotzdem kennt er mich nicht.} (Ou : \all{... . Trotzdem kennt er mich nicht.} Verbe \all{kennt} en 2e pos).
    \item \all{Ich wünschte, die Gesellschaft wäre toleranter.} (KII \all{wäre} à la fin).
    \item \all{Man sollte die Rechte der Minderheiten respektieren.} (KII \all{sollte} + infinitif \all{respektieren} fin).
    \item \all{Ich schreibe einen \textbf{lange\emph{n}} Brief an die Redaktion.} (Mixte acc. masc. : \all{einen} + \all{lang\emph{en}}).
    \item \all{Wir brauchen die \textbf{notwendig\emph{e}} Toleranz.} (Faible nom. fém. : \all{die} + \all{notwendig\emph{e}}).
\end{enumerate}
\end{corrige}

\newpage

\coursec{Méthodes et exercices résolus}

\coursec{Le problème des minorités}

\begin{textebox}[Ein Brief aus Bautzen — Texte support]
\all{„Liebe Marie, seit einem Jahr wohne ich in Bautzen. Hier leben viele Sorben, eine kleine slawische Minderheit. Sie haben ihre eigene Sprache und ihre eigenen Schulen. Obwohl sie oft diskriminiert wurden, haben sie ihre Kultur nie vergessen. Trotzdem fühlen sich manche junge Sorben heute als Fremde im eigenen Land. Wenn ich mehr Zeit hätte, würde ich Sorbisch lernen! Dein Paul.“}
\end{textebox}

\begin{savaistu}[Les Sorbes (Sorben) — minorité slave en Allemagne]
Les Sorbes (\all{die Sorben}) sont une minorité slave occidentale reconnue en Allemagne (Brandebourg, Saxe). Ils parlent le sorbe (\all{Sorbisch}), ont leurs propres écoles, médias et institutions culturelles. Malgré une histoire marquée par l'assimilation forcée (époque nazie, RDA), ils bénéficient aujourd'hui de droits constitutionnels (\all{Minderheitenrechte}). Bautzen (\all{Budyšin}) est leur centre culturel. Cette situation illustre les thèmes du cours : \cle{die Minderheit}, \cle{die Sprache}, \cle{die Rechte}, \cle{die Diskriminierung}, \cle{die Integration}, \cle{das Vorurteil}.
\end{savaistu}

\courssub{Méthode 1 : Comprendre un texte sur les minorités (global puis détaillé)}

\begin{methode}[Lire et comprendre un texte sur les Minderheiten]
\begin{enumerate}
\item \textbf{Première lecture globale} : repérer le type de texte (ici \all{Brief} personnel), l'émetteur (\all{Paul}), le destinataire (\all{Marie}), le thème (minorités, Sorbes). Souligner les mots-clés du champ lexical : \all{die Minderheit, die Diskriminierung, das Vorurteil, die Integration, die Rechte, die Sprache}.
\item \textbf{Deuxième lecture détaillée} : pour chaque paragraphe (ou phrase complexe), formuler mentalement l'idée principale en une phrase française.
\item \textbf{Répondre aux questions de compréhension} : reformuler avec ses propres mots (allemand), jamais recopier de longs passages. Citer un court extrait entre guillemets pour justifier : \all{Im Text steht: „...“}.
\item \textbf{Questions personnelles} (\all{Was meinen Sie dazu?}) : donner son avis avec \all{Ich meine, dass... / Meiner Meinung nach...} et un argument.
\item \textbf{Vérifier} : chaque réponse doit être une phrase complète avec le verbe conjugué en 2\textsuperscript{e} position.
\end{enumerate}
\end{methode}

\begin{exoresolu}[Compréhension de texte — type Bac]
\textbf{Texte (extrait) :} \all{„Liebe Marie, seit einem Jahr wohne ich in Bautzen. Hier leben viele Sorben, eine kleine slawische Minderheit. Sie haben ihre eigene Sprache und ihre eigenen Schulen. Obwohl sie oft diskriminiert wurden, haben sie ihre Kultur nie vergessen. Trotzdem fühlen sich manche junge Sorben heute als Fremde im eigenen Land. Wenn ich mehr Zeit hätte, würde ich Sorbisch lernen! Dein Paul.“}

\textbf{Questions :}
1) Wo wohnt Paul?
2) Welche Minderheit lebt in Bautzen?
3) Haben die Sorben ihre Kultur vergessen?
4) Was würde Paul gern lernen?
\tcblower
\textbf{Raisonnement :} on localise chaque information dans le texte, puis on reformule en phrase complète allemande.

1) \all{Paul wohnt seit einem Jahr in Bautzen.} (Paul habite à Bautzen depuis un an.)

2) \all{In Bautzen leben viele Sorben, eine kleine slawische Minderheit.} (À Bautzen vivent de nombreux Sorabes, une petite minorité slave.)

3) \all{Nein, obwohl sie oft diskriminiert wurden, haben sie ihre Kultur nie vergessen.} Justification : le texte dit explicitement \all{„haben sie ihre Kultur nie vergessen“}.

4) \all{Wenn er mehr Zeit hätte, würde Paul Sorbisch lernen.} Attention : la question porte sur le souhait de Paul (3\textsuperscript{e} personne), on reprend le Konjunktiv II (\all{hätte, würde lernen}).

\textbf{Conclusion :} réponses complètes, verbe en 2\textsuperscript{e} position, citations courtes à l'appui : la méthode est respectée.
\end{exoresolu}

\courssub{Méthode 2 : Employer le vocabulaire du thème (Minderheiten und Integration)}

\begin{methode}[Maîtriser le lexique thématique]
\begin{enumerate}
\item Apprendre chaque nom \textbf{avec son article et son pluriel} :
      \all{die Minderheit, -en} ; \all{die Diskriminierung, -en} ; \all{das Recht, -e} ; \all{das Vorurteil, -e} ; \all{die Integration} (souvent au singulier) ; \all{die Sprache, -n}.
\item Construire des \textbf{familles de mots} (dérivation) :
      \all{diskriminieren (Verb) → die Diskriminierung (Nom) → diskriminiert (Partizip II / Adj.)} ;
      \all{integrieren → die Integration → integriert} ;
      \all{vorurteilen (selten) → das Vorurteil → vorurteilsfrei (Adj.)}.
\item Mémoriser les \textbf{collocations} (associations figées) :
      \all{Rechte haben / verteidigen / einfordern} ;
      \all{gegen Vorurteile kämpfen / ankämpfen} ;
      \all{eine Sprache sprechen / lernen / pflegen} ;
      \all{in die Gesellschaft integriert sein / werden} (préposition \all{in} + accusatif).
\item Réutiliser au moins trois mots du thème dans chaque production écrite ou orale.
\end{enumerate}
\end{methode}

\begin{popfigure}
\begin{tikzpicture}[mindmap, grow cyclic, every node/.style=concept, concept color=popblue!80,
    level 1/.append style={level distance=3.5cm, sibling angle=120, font=\small},
    level 2/.append style={level distance=2.8cm, sibling angle=50, font=\footnotesize}]
\node[concept, align=center]{Minderheiten\\(die Minderheiten)}
    child[concept color=popgreen!70] { node[concept, align=center]{Rechte\\(die Rechte)} }
    child[concept color=poporange!70] { node[concept, align=center]{Sprache\\(die Sprachen)} }
    child[concept color=poppurple!70] { node[concept, align=center]{Integration\\(die Integration)} };
\node[concept, align=center] at (0,-4.5){Diskriminierung\\(die Diskriminierung)} [clockwise from=60]
    child[concept color=poppink!70] { node[concept, align=center]{Vorurteile\\(die Vorurteile)} }
    child[concept color=popgold!70] { node[concept, align=center]{kämpfen gegen\\(kämpfen gegen)} };
\end{tikzpicture}
\legende{Carte mentale : champ lexical des minorités et de l'intégration (noms avec articles, verbes avec prépositions)}
\end{popfigure}

\begin{exoresolu}[Vocabulaire — phrases à compléter (cas et nombre)]
Complétez avec la forme correcte : \all{Minderheit, Vorurteile, Rechte, Integration, Sprache}.

1) Die Sorben sind eine kleine \dots\ in Deutschland.
2) Viele Menschen kämpfen gegen \dots\ .
3) Alle Menschen haben die gleichen \dots\ .
4) Die \dots\ der Ausländer in die Gesellschaft ist wichtig.
5) Die Sorben sprechen ihre eigene \dots\ .
\tcblower
\textbf{Solutions et justifications cas par cas :}

1) \all{Die Sorben sind eine kleine \textbf{Minderheit}.} (Sujet \all{Die Sorben} → nominatif pluriel ; attribut \all{eine kleine Minderheit} → nominatif singulier féminin après article indéfini).

2) \all{Viele Menschen kämpfen gegen \textbf{Vorurteile}.} (Verbe \all{kämpfen gegen} + accusatif ; \all{Vorurteile} = pluriel de \all{das Vorurteil}).

3) \all{Alle Menschen haben die gleichen \textbf{Rechte}.} (Verbe \all{haben} + accusatif ; \all{Rechte} = pluriel de \all{das Recht} ; adjectif \all{gleichen} : déclinaison faible après \all{die}).

4) \all{Die \textbf{Integration} der Ausländer in die Gesellschaft ist wichtig.} (Sujet \all{Die Integration} → nominatif singulier féminin ; pas de pluriel usuel pour \all{Integration} dans ce sens).

5) \all{Die Sorben sprechen ihre eigene \textbf{Sprache}.} (Verbe \all{sprechen} + accusatif ; \all{Sprache} = singulier féminin ; adjectif \all{eigene} : déclinaison faible après possessif \all{ihre}).

\textbf{Rappel :} vérifier à chaque fois le cas exigé par le verbe ou la préposition et le nombre (singulier/pluriel).
\end{exoresolu}

\courssub{Grammaire 1 : La déclinaison de l'adjectif épithète}

\begin{propriete}[Règle générale : trois déclinaisons selon le déterminant]
L'adjectif épithète (placé entre le déterminant et le nom) s'accorde en \textbf{genre, nombre, cas}. La terminaison dépend de ce qui le précède :
\begin{enumerate}
\item \textbf{Déclinaison faible} (après article \textbf{défini} \all{der, die, das, den, dem, des} ou \all{dieser, jener, jeder, welcher, mancher, solcher} ou \all{alle} au pluriel) : l'article marque le cas → l'adjectif prend \textbf{-e} (nom/acc sg f/n, nom sg m) ou \textbf{-en} (autres cas).
\item \textbf{Déclinaison mixte} (après article \textbf{indéfini} \all{ein, eine, ein}, possessif \all{mein, dein, sein, ihr, unser, euer, Ihr}, négatif \all{kein}) : l'article ne marque pas toujours le genre/cas → l'adjectif \textbf{porte la marque forte} là où l'article est neutre (\all{-er} masc nom, \all{-es} neut nom/acc, \all{-e} fém nom/acc) ; sinon \textbf{-en}.
\item \textbf{Déclinaison forte} (\textbf{sans article} ni déterminant) : l'adjectif porte \textbf{toutes les marques} de l'article défini (\all{-er, -e, -es, -en}).
\end{enumerate}
\textbf{Règle d'or au pluriel} : après \all{die / alle / diese / meine / keine / viele / wenige / einige} → toujours \textbf{-en} (sauf nominatif/accusatif sans article → \textbf{-e}).
\end{propriete}

\begin{exemplebox}[Tableau récapitulatif — Adjectif \all{klein} + \all{Minderheit (f)} / \all{Lehrer (m)} / \all{Kind (n)} / \all{Minderheiten (pl)}]
\begin{center}
\begin{tabular}{|l|c|c|c|c|}
\hline
\rowcolor{popblueL}
\textbf{Cas} & \textbf{faible (\all{die/der/das})} & \textbf{mixte (\all{ein/kein/mein})} & \textbf{forte (\all{Ø})} & \textbf{Pluriel (faible/mixte)} \\
\hline
Nominatif & \all{die kleine Minderheit} & \all{eine kleine Minderheit} & \all{kleine Minderheit} & \all{die kleinen Minderheiten} \\
Accusatif & \all{die kleine Minderheit} & \all{eine kleine Minderheit} & \all{kleine Minderheit} & \all{die kleinen Minderheiten} \\
Datif & \all{der kleinen Minderheit} & \all{einer kleinen Minderheit} & \all{kleiner Minderheit} & \all{den kleinen Minderheiten} \\
Génitif & \all{der kleinen Minderheit} & \all{einer kleinen Minderheit} & \all{kleiner Minderheit} & \all{der kleinen Minderheiten} \\
\hline
Nominatif & \all{der kleine Lehrer} & \all{ein \textbf{kleiner} Lehrer} & \all{\textbf{kleiner} Lehrer} & \all{die kleinen Lehrer} \\
Accusatif & \all{den kleinen Lehrer} & \all{einen kleinen Lehrer} & \all{\textbf{kleinen} Lehrer} & \all{die kleinen Lehrer} \\
\hline
\end{tabular}
\end{center}
\end{exemplebox}

\begin{exoresolu}[Grammaire — adjectifs à décliner (contexte minorités)]
Complétez avec la terminaison correcte :
1) Die deutsch\dots\ Minderheit in Polen ist klein.
2) Ein jung\dots\ Sorbischlehrer arbeitet in Bautzen.
3) Die Sorben haben ihre eigen\dots\ Sprache behalten.
4) Viele diskriminier\dots\ Menschen brauchen Hilfe.
5) Er kämpft für die Rechte der klein\dots\ Minderheiten.
\tcblower
\textbf{Solutions et justifications détaillées :}

1) \all{Die \textbf{deutsche} Minderheit} : article défini \all{die} + nominatif féminin singulier → déclinaison faible \textbf{-e}.

2) \all{Ein \textbf{junger} Sorbischlehrer} : \all{ein} ne marque pas le masculin au nominatif → l'adjectif porte la marque forte \textbf{-er} (déclinaison mixte).

3) \all{ihre \textbf{eigene} Sprache} : possessif \all{ihre} (marque déjà le fém. acc.) + accusatif féminin → déclinaison faible \textbf{-e}.

4) \all{Viele \textbf{diskriminierte} Menschen} : \all{viele} (quantifieur) au pluriel \textbf{sans article défini} devant, sujet (nominatif) → déclinaison forte pluriel nominatif/accusatif \textbf{-e}. (Participe II \all{diskriminiert} utilisé comme adjectif).
\emph{Attention :} \all{Vielen diskriminierten Menschen} serait correct seulement au datif (\all{Helfen vielen diskriminierten Menschen}).

5) \all{die Rechte der \textbf{kleinen} Minderheiten} : génitif pluriel après article défini \all{der} → déclinaison faible \textbf{-en}.

\textbf{Conclusion :} toujours se poser la question « \textbf{Quel déterminant devant ?} » (défini, indéfini/possessif/négatif, rien) et « \textbf{Quel cas ?} » avant de choisir la terminaison.
\end{exoresolu}

\courssub{Grammaire 2 : La concession (\all{obwohl} / \all{trotzdem} / \all{trotz})}

\begin{propriete}[Concession : subordonnée vs phrase indépendante]
\begin{itemize}
\item \textbf{\all{obwohl}} (bien que) = conjonction de subordination → verbe conjugué \textbf{à la fin} de la subordonnée.
      \all{Obwohl die Sorben eine kleine Minderheit \textbf{sind}, haben sie ihre eigenen Schulen.}
      Si la subordonnée est en tête, la principale commence par le verbe (inversion) : \all{..., \textbf{haben} sie ...}.
\item \textbf{\all{trotzdem}} (pourtant, malgré cela) = adverbe de liaison → phrase \textbf{indépendante}, verbe en \textbf{2\textsuperscript{e} position}.
      \all{Die Sorben sind eine kleine Minderheit. \textbf{Trotzdem} \textbf{haben} sie ihre eigenen Schulen.}
\item \textbf{\all{trotz}} + \textbf{génitif} (langage soutenu/écrit) ou datif (parlé) : préposition.
      \all{Trotz \textbf{der} Diskriminierung (Gen.) haben sie ihre Sprache behalten.}
\item \textbf{Interdiction} : ne jamais combiner *\all{obwohl ... trotzdem} dans la même phrase.
\item \textbf{Variantes} : \all{obgleich, obzwar} (soutenu, littéraire) = \all{obwohl}.
\end{itemize}
\end{propriete}

\begin{exoresolu}[Thème — traduction avec la concession]
Traduisez en allemand :
1) Bien que les Sorabes soient une petite minorité, ils ont leurs propres écoles.
2) Ils ont été souvent discriminés. Pourtant, ils n'ont jamais oublié leur culture.
3) Malgré les préjugés, beaucoup de jeunes Sorabes sont fiers de leur langue.
\tcblower
\textbf{Solutions détaillées :}

1) \all{Obwohl die Sorben eine kleine Minderheit \textbf{sind}, \textbf{haben} sie ihre eigenen Schulen.}
Justification : \all{obwohl} renvoie \all{sind} en fin ; principale en inversion (\all{haben} en 1\textsuperscript{re}) ; \all{eigenen} : possessif \all{ihre} + pluriel accusatif → \all{-en}.

2) \all{Sie wurden oft diskriminiert. \textbf{Trotzdem} \textbf{haben} sie ihre Kultur nie \textbf{vergessen}.}
Justification : \all{trotzdem} ouvre phrase indépendante, verbe \all{haben} en 2\textsuperscript{e} position ; \all{vergessen} verbe fort : \all{vergessen – vergaß – vergessen} (Partizip II \all{vergessen}).

3) \all{\textbf{Trotz der Vorurteile} \textbf{sind} viele junge Sorben stolz \textbf{auf} ihre Sprache.}
Justification : \all{trotz} + génitif (\all{der Vorurteile}) ; \all{stolz auf} + accusatif (faux ami : « fier de » ≠ *\all{stolz von}) ; \all{viele junge Sorben} : pluriel sans article défini → déclinaison forte nominatif \all{junge} (\textbf{-e}).
\end{exoresolu}

\courssub{Grammaire 3 : Le subjonctif II (Konjunktiv II) pour exprimer un souhait}

\begin{propriete}[Formation et emploi du Konjunktiv II (souhait / irréel du présent)]
\begin{enumerate}
\item \textbf{Formation} :
      \begin{itemize}
      \item Verbes auxiliaires/modaux : radical Präteritum + Umlaut + terminaisons KII :
            \all{haben → hätte}, \all{sein → wäre}, \all{können → könnte}, \all{müssen → müsste}, \all{dürfen → dürfte}, \all{wollen → wollte}, \all{sollen → sollte}, \all{werden → würde}.
      \item Verbes forts : radical Präteritum + Umlaut (si possible) + terminaisons : \all{geben → gäbe}, \all{finden → fände}, \all{gehen → ginge}.
      \item Verbes faibles / la plupart : \textbf{\all{würde} + Infinitif} à la fin (forme de substitution courante) : \all{ich würde lernen, ich würde machen}.
      \end{itemize}
\item \textbf{Souhait réalisable / poli} : \all{Ich \textbf{hätte} gern ... / Ich \textbf{würde} gern ... + Infinitif}.
      \all{Ich würde gern Sorbisch lernen.}
\item \textbf{Souhait irréel / regret} (si seulement...) : \all{Wenn ich \textbf{doch / nur / bloß} mehr Zeit \textbf{hätte}!} (Particules \all{doch, nur, bloß} renforcent le regret).
\item \textbf{Période hypothétique (Si...)} : \all{Wenn} + KII (subordonnée) → Principal aussi au KII.
      \all{Wenn ich reich \textbf{wäre}, \textbf{würde} ich eine Schule \textbf{bauen}.}
\item \textbf{Réflexe Bac} : Repérer \all{wenn} + verbe à la fin (KII) → s'attendre à KII dans la principale (\all{würde/hätte/wäre/gäbe}).
\end{enumerate}
\end{propriete}

\begin{exoresolu}[Expression écrite guidée — Production modèle]
\textbf{Consigne :} Rédigez un court paragraphe (5–6 phrases) : « \all{Wenn ich die Lage der Minderheiten verbessern könnte...} » (Si je pouvais améliorer la situation des minorités...). Utilisez au moins \textbf{deux formes de Konjunktiv II}, \textbf{une concession} (\all{obwohl} ou \all{trotzdem}) et \textbf{trois mots du vocabulaire du thème}.
\tcblower
\textbf{Raisonnement pas à pas :}
1. Ouvrir par l'hypothèse principale (\all{Wenn} + KII).
2. Développer 2 actions concrètes au KII (\all{würde + Inf.} ou \all{hätte gern}).
3. Insérer une concession (\all{Obwohl} + indicatif réel ou \all{Trotzdem} + phrase principale).
4. Conclure par une conséquence générale (\all{Wenn} + KII → KII) ou un appel à l'action (indicatif présent).

\textbf{Production modèle :}

\all{Wenn ich die Lage der Minderheiten \textbf{verbessern könnte}, \textbf{würde} ich zuerst gegen die \textbf{Vorurteile} \textbf{kämpfen}. Ich \textbf{hätte} gern mehr Schulen, wo die Kinder ihre eigene \textbf{Sprache} lernen \textbf{könnten}. \textbf{Obwohl} die \textbf{Integration} oft schwer \textbf{ist}, \textbf{würde} ich allen Menschen die gleichen \textbf{Rechte} \textbf{geben}. \textbf{Wenn} alle Menschen tolerant \textbf{wären}, \textbf{gäbe} es keine \textbf{Diskriminierung} mehr. \textbf{Trotzdem} \textbf{müssen} wir jeden Tag für diese Ideale \textbf{arbeiten}.}

\textbf{Traduction :} Si je pouvais améliorer la situation des minorités, je lutterais d'abord contre les préjugés. J'aimerais avoir plus d'écoles où les enfants pourraient apprendre leur propre langue. Bien que l'intégration soit souvent difficile, je donnerais à tous les mêmes droits. Si tout le monde était tolérant, il n'y aurait plus de discrimination. Pourtant, nous devons travailler chaque jour pour ces idéaux.

\textbf{Vérification grammaticale (check-list Bac) :}
\begin{itemize}
\item KII corrects : \all{könnte, würde, hätte, könnten, wäre, gäbe}.
\item Concession : \all{Obwohl ... ist} (verbe à la fin, indicatif car fait réel) ; \all{Trotzdem} + verbe 2\textsuperscript{e} pos.
\item Vocabulaire thème (5 mots) : \all{Vorurteile, Sprache, Integration, Rechte, Diskriminierung}.
\item Ordre des mots : verbe 2\textsuperscript{e} pos. principale, verbe fin subordonnée.
\item Majuscules à tous les noms : \all{Lage, Minderheiten, Vorurteile, Schulen, Kinder, Sprache, Integration, Rechte, Diskriminierung, Menschen, Ideale}.
\end{itemize}
\end{exoresolu}

\courssub{Entraînement : Exercices d'application (Bac blanc)}

\begin{exercice}[Grammaire \& Vocabulaire — Entraînement complet]
\textbf{A. Déclinaison :} Mettez l'adjectif entre parenthèses à la forme correcte.
1) \all{Die (klein) \_\_\_\_\_\_ Minderheit hat (eigen) \_\_\_\_\_\_ Schulen.}
2) \all{Ein (jung) \_\_\_\_\_\_ Lehrer unterrichtet (sorbisch) \_\_\_\_\_\_ Kinder.}
3) \all{Wir helfen (diskriminiert) \_\_\_\_\_\_ Menschen.}
4) \all{Ohne (eigen) \_\_\_\_\_\_ Sprache verliert man (kulturell) \_\_\_\_\_\_ Identität.}

\textbf{B. Concession :} Reliez les deux phrases par \all{obwohl} (subordonnée en tête) ou \all{trotzdem}.
1) \all{Die Sorben sind wenige. Sie haben ihre Kultur bewahrt.}
2) \all{Es gibt viele Vorurteile. Die jungen Sorben sind stolz auf ihre Herkunft.}

\textbf{C. Konjunktiv II :} Mettez le verbe au KII (forme \all{würde} ou forme simple).
1) \all{Wenn ich Zeit (haben) \_\_\_\_\_\_, ich (lernen) \_\_\_\_\_\_ Sorbisch.}
2) \all{Ich (wünschen) \_\_\_\_\_\_, dass alle Menschen gleich (sein) \_\_\_\_\_\_.}
3) \all{Was (machen / Sie) \_\_\_\_\_\_ an meiner Stelle?}

\textbf{D. Expression écrite :} Répondez en 3 phrases minimum : \all{Was würden Sie tun, wenn Sie Minister für Minderheitenfragen wären?} (Utilisez KII, un mot de concession, 2 mots du thème).
\end{exercice}

\begin{corrige}[Exercice Entraînement complet]
\textbf{A. Déclinaison :}
1) \all{Die \textbf{kleine} Minderheit hat \textbf{eigene} Schulen.} (\all{die} faible fém. sg nom/acc → -e ; \all{eigene} : pluriel acc. après Ø → forte -e).
2) \all{Ein \textbf{junger} Lehrer unterrichtet \textbf{sorbische} Kinder.} (\all{ein} masc nom → mixte -er ; \all{sorbische} : pluriel acc. sans art. → forte -e).
3) \all{Wir helfen \textbf{diskriminierten} Menschen.} (Datif pluriel après \all{helfen} : \all{den} (faible) ou \all{diskriminierten} (mixte/forte) → toujours \textbf{-en}).
4) \all{Ohne \textbf{eigene} Sprache verliert man \textbf{kulturelle} Identität.} (\all{ohne} + acc. sg fém sans art. → forte \textbf{-e} ; \all{kulturelle} : acc. sg fém sans art. → forte \textbf{-e}).

\textbf{B. Concession :}
1) \all{\textbf{Obwohl die Sorben wenige sind}, haben sie ihre Kultur bewahrt.} (ou : \all{Die Sorben sind wenige. \textbf{Trotzdem} haben sie ihre Kultur bewahrt.})
2) \all{\textbf{Obwohl es viele Vorurteile gibt}, sind die jungen Sorben stolz auf ihre Herkunft.} (ou : \all{Es gibt viele Vorurteile. \textbf{Trotzdem} sind die jungen Sorben stolz auf ihre Herkunft.})

\textbf{C. Konjunktiv II :}
1) \all{Wenn ich Zeit \textbf{hätte}, \textbf{würde} ich Sorbisch \textbf{lernen}.} (ou \all{lernte}).
2) \all{Ich \textbf{wünschte}, dass alle Menschen gleich \textbf{wären}.} (Subjonctif I \all{seien} possible en discours rapporté, mais KII \all{wären} pour souhait irréel).
3) \all{Was \textbf{würden Sie} an meiner Stelle \textbf{machen}?} (Politesse / conseil).

\textbf{D. Expression écrite (proposition) :}
\all{Wenn ich Minister für Minderheitenfragen \textbf{wäre}, \textbf{würde} ich ein Gesetz gegen \textbf{Diskriminierung} erlassen. \textbf{Obwohl} das Geld knapp \textbf{ist}, \textbf{würde} ich mehr Geld für \textbf{Sprachschulen} geben. \textbf{Trotzdem} braucht man auch die Unterstützung der Gesellschaft.}
\end{corrige}

\begin{aretenir}[Récapitulatif de la leçon AL04 — Points clés pour le Bac]
\begin{itemize}
\item \textbf{Texte} : Lettre personnelle (\all{Brief}) sur les Sorbes (minorité slave en Saxe). Structure : situation → faits culturels → concession (\all{obwohl/trotzdem}) → souhait personnel (\all{KII}).
\item \textbf{Vocabulaire} : \all{die Minderheit(-en), die Diskriminierung(-en), das Recht(-e), das Vorurteil(-e), die Integration, die Sprache(-n)} + familles (\all{diskriminieren, integrieren, vorurteilsfrei}) + collocations (\all{kämpfen gegen, stolz auf, Rechte haben}).
\item \textbf{Adjectif} : 3 déclinaisons (Faible \all{die} → -e/-en ; Mixte \all{ein/kein} → marques fortes si article neutre ; Forte \all{Ø} → marques de l'article). \textbf{Pluriel : après article/déterminant → toujours -en.}
\item \textbf{Concession} : \all{obwohl} + verbe \textbf{à la fin} (inversion principale) ; \all{trotzdem} + verbe \textbf{2\textsuperscript{e} pos.} ; \all{trotz} + \textbf{Génitif}. Jamais \all{obwohl ... trotzdem}.
\item \textbf{Konjunktiv II (Souhait)} : \all{hätte, wäre, könnte, würde + Inf.} ; \all{Wenn} + KII → Principal KII. Particules \all{doch/nur/bloß} pour regret.
\item \textbf{Production} : Réponses complètes (verbe 2\textsuperscript{e}), citations courtes, KII pour hypothèse, concession pour nuancer, vocabulaire thématique imposé.
\end{itemize}
\end{aretenir}

\begin{attention}[Les 5 erreurs qui coûtent des points au Bac — Check-list finale]
\begin{enumerate}
\item \textbf{Majuscules oubliées} : *\all{die minderheit} → \all{\textbf{M}inderheit}. \textbf{Tout nom} prend une majuscule en allemand.
\item \textbf{Verbe mal placé après \all{obwohl}} : *\all{Obwohl sie sind diskriminiert...} → \all{Obwohl sie diskriminiert \textbf{wurden}, ...}. Verbe conjugué \textbf{à la fin}. Et inversion en principale : \all{..., \textbf{haben} sie ...}.
\item \textbf{Confusion \all{obwohl} / \all{trotzdem}} : *\all{Obwohl sie arm sind, trotzdem sind sie stolz.} → Soit \all{Obwohl sie arm sind, sind sie stolz.} Soit \all{Sie sind arm. Trotzdem sind sie stolz.}
\item \textbf{Faux amis / Calques} : « fier de » = \all{stolz \textbf{auf} (+ Akk.)}, pas *\all{stolz von} ; « lutter contre » = \all{kämpfen \textbf{gegen}}, pas *\all{kämpfen mit} ; \all{die Sprache} = la langue, \all{die Rede} = le discours ; \all{das Vorurteil} = le préjugé, \all{die Vorurteile} (pl.).
\item \textbf{Déclinaison au hasard} : Après \all{ein} masc. nom. → \all{-er} (\all{ein \textbf{junger} Lehrer}) ; Après \all{die} pluriel → \all{-en} (\all{die \textbf{kleinen} Minderheiten}) ; Sans article pluriel nom/acc → \all{-e} (\all{\textbf{viele diskriminierte} Menschen}). \textbf{Vérifier le déterminant AVANT d'écrire la terminaison.}
\end{enumerate}
\end{attention}

\newpage

\coursec{Exercices}

\courssub{Je vérifie mes connaissances}

\begin{exercice}[QCM : Compréhension globale du texte]
Le texte \guillemotleft Ein Brief aus Bautzen \guillemotright\ est une lettre écrite par un jeune Sorabe. Choisissez la bonne réponse.
\begin{enumerate}
    \item Quel est le sujet principal de la lettre ?
    \begin{enumerate}
        \item La cuisine traditionnelle sorabe.
        \item La situation de la minorité sorabe en Lusace (Brandebourg/Saxe).
        \item L'histoire de la ville de Bautzen au Moyen Âge.
    \end{enumerate}
    \item Quel sentiment exprime l'auteur face à la perte de la langue sorabe ?
    \begin{enumerate}
        \item De l'indifférence.
        \item De la colère et de la tristesse.
        \item De la joie car l'allemand est plus pratique.
    \end{enumerate}
    \item Quelle mesure concrète l'auteur souhaite-t-il pour l'avenir ?
    \begin{enumerate}
        \item L'interdiction de l'allemand dans les écoles.
        \item Le bilinguisme obligatoire dans l'administration et les écoles.
        \item Le déménagement de tous les Sorabes en Tchéquie.
    \end{enumerate}
\end{enumerate}
\end{exercice}

\begin{exercice}[Vrai ou Faux ? Justifiez par une phrase du texte]
Déterminez si les affirmations suivantes sont vraies (V) ou fausses (F). Recopiez la phrase du texte qui justifie votre choix.
\begin{enumerate}
    \item Les Sorabes sont une minorité germanique reconnue en Allemagne.
    \item L'auteur pense que l'intégration signifie l'assimilation totale à la culture majoritaire.
    \item Il existe des écoles où l'enseignement se fait en sorabe.
    \item L'auteur utilise le subjonctif II pour exprimer un souhait irréel du présent.
\end{enumerate}
\end{exercice}

\begin{exercice}[Vocabulaire thématique : Familles de mots et prépositions]
Complétez le tableau suivant avec les formes manquantes (nom, verbe, adjectif) et l'article/le pluriel pour les noms. Traduisez en français.
\begin{center}
\begin{tabularx}{\linewidth}{|X|X|X|X|}
\hline
\rowcolor{popblueL} \textbf{Nom (Art. Pl.)} & \textbf{Verbe} & \textbf{Adjectif / Participe} & \textbf{Français} \\
\hline
die Diskriminierung, -en & diskriminieren & diskriminierend & la discrimination / discriminer / discriminatoire \\
\hline
die Integration, -en &  & integriert &  \\
\hline
 & vorurteilen (reflexiv) &  & le préjugé / avoir des préjugés / préjugé \\
\hline
das Recht, -e &  & rechtlich &  \\
\hline
die Minderheit, -en &  & minderheitlich &  \\
\hline
 & assimilieren &  & l'assimilation / assimiler / assimilé \\
\hline
\end{tabularx}
\end{center}
\end{exercice}

\begin{exercice}[Conjugaison : Subjonctif II (Konjunktiv II) du souhait]
Mettez les verbes entre parenthèses au \textbf{Konjunktiv II} (forme du souhait : \emph{würde} + infinitif ou forme simple \emph{hätte, wäre, gäbe, könnte, wollte}). Attention à la place du verbe dans la subordonnée introduite par \emph{dass} ou \emph{wenn}.
\begin{enumerate}
    \item Ich wünschte, alle Menschen \_\_\_\_\_\_\_\_\_\_ (haben) mehr Respekt voreinander.
    \item Wenn ich Bundeskanzler \_\_\_\_\_\_\_\_\_\_ (sein), ich \_\_\_\_\_\_\_\_\_\_ (geben) den Minderheiten mehr Rechte.
    \item Mein Vater \_\_\_\_\_\_\_\_\_\_ (können) gern sorbisch sprechen, aber er hat es nie gelernt.
    \item Wäre es nicht besser, wenn wir alle \_\_\_\_\_\_\_\_\_\_ (verstehen) uns besser?
    \item Ich \_\_\_\_\_\_\_\_\_\_ (möchte) nicht, dass die Sprache \_\_\_\_\_\_\_\_\_\_ (sterben).
\end{enumerate}
\end{exercice}

\courssub{Je m'entraîne}

\begin{exercice}[Traduction thématique : Français $\rightarrow$ Allemand]
Traduisez les phrases suivantes en allemand en mobilisant le vocabulaire de la leçon (minorités, droits, intégration, préjugés) et la grammaire (déclinaison adjectivale, concession).
\begin{enumerate}
    \item Bien que les minorités aient des droits, elles sont souvent discriminées.
    \item L'intégration ne veut pas dire que l'on doit oublier sa propre langue.
    \item Beaucoup de jeunes Sorabes parlent couramment l'allemand et le sorabe.
    \item Malgré les préjugés, la culture sorabe est vivante en Lusace.
    \item Je souhaite que le gouvernement fasse plus pour la protection des langues minoritaires.
\end{enumerate}
\end{exercice}

\begin{exercice}[Transformation : Concession \emph{obwohl} $\leftrightarrow$ \emph{trotzdem}]
Réécrivez les phrases en remplaçant \emph{obwohl} par \emph{trotzdem} (ou inversement) sans changer le sens. Attention à l'ordre des mots (position du verbe) !
\begin{enumerate}
    \item Obwohl die Sorben eine kleine Minderheit sind, haben sie eigene Schulen.
    \item Die Sprache ist bedroht, trotzdem lernen viele Kinder sie in der Schule.
    \item Obwohl er kein Sorbe ist, lernt er die Sprache.
    \item Die Gesetze schützen die Minderheit, trotzdem gibt es Probleme im Alltag.
    \item Obwohl sie müde war, ging sie zur Demonstration für Minderheitenrechte.
\end{enumerate}
\end{exercice}

\begin{exercice}[Texte à trous : Déclinaison de l'adjectif (faible, mixte, forte)]
Complétez le texte sur les Sorabes en mettant les adjectifs entre parenthèses à la bonne terminaison. Indiquez entre crochets le type de déclinaison utilisée (F=Faible, M=Mixte, S=Forte).
\begin{textebox}[Lückentext: Die Sorben in der Lausitz]
Die Sorben sind \_\_\_\_\_\_\_\_\_\_ (klein) slawische Minderheit in Deutschland. [ \_\_\_\_ ]
Sie leben in \_\_\_\_\_\_\_\_\_\_ (östlich) Teil von Sachsen und Brandenburg. [ \_\_\_\_ ]
\_\_\_\_\_\_\_\_\_\_ (viel) Sorben sprechen noch ihre Muttersprache. [ \_\_\_\_ ]
In \_\_\_\_\_\_\_\_\_\_ (sorbisch) Schulen wird zweisprachig unterrichtet. [ \_\_\_\_ ]
Die Regierung gibt \_\_\_\_\_\_\_\_\_\_ (finanziell) Unterstützung für die Kultur. [ \_\_\_\_ ]
Ohne \_\_\_\_\_\_\_\_\_\_ (politisch) Wille würde die Sprache sterben. [ \_\_\_\_ ]
Man sieht oft \_\_\_\_\_\_\_\_\_\_ (zweisprachig) Straßenschilder in Bautzen. [ \_\_\_\_ ]
\_\_\_\_\_\_\_\_\_\_ (jung) Sorben sind oft stolz auf ihre Identität. [ \_\_\_\_ ]
Trotz \_\_\_\_\_\_\_\_\_\_ (schwierig) Geschichte haben sie überlebt. [ \_\_\_\_ ]
Ich wünsche mir \_\_\_\_\_\_\_\_\_\_ (gut) Zukunft für die Minderheit. [ \_\_\_\_ ]
\end{textebox}
\end{exercice}

\begin{exercice}[Appariement : Structures du souhait (Konjunktiv II)]
Reliez chaque début de phrase (1--5) à la fin logique (a--e) et conjuguez le verbe entre parenthèses au \emph{Konjunktiv II}.
\begin{center}
\begin{tabular}{|p{5cm}|p{5cm}|}
\hline
1. Ich wünschte, die Regierung \_\_\_\_\_\_\_\_\_\_ (geben) & a. \_\_\_\_\_\_\_\_\_\_ (sein) mehr zweisprachige Kitas. \\
2. Wenn alle Menschen \_\_\_\_\_\_\_\_\_\_ (haben) & b. \_\_\_\_\_\_\_\_\_\_ (können) wir friedlicher zusammenleben. \\
3. Wäre ich Minister, ich \_\_\_\_\_\_\_\_\_\_ (machen) & c. \_\_\_\_\_\_\_\_\_\_ (haben) mehr Respekt vor Minderheiten. \\
4. Es wäre schön, wenn es \_\_\_\_\_\_\_\_\_\_ (geben) & d. \_\_\_\_\_\_\_\_\_\_ (fördern) die Sprachen stärker. \\
5. Die Kinder \_\_\_\_\_\_\_\_\_\_ (lernen) gern sorbisch, & e. wenn die Eltern sie unterstützen. \\
\hline
\end{tabular}
\end{center}
\end{exercice}

\courssub{Je me prépare au Bac}

\begin{exercice}[Compréhension écrite type Baccalauréat (Texte original)]
Lisez le texte suivant, puis répondez aux questions \textbf{en allemand} (phrases complètes). Le nombre de mots indicatif est donné entre parenthèses.

\begin{textebox}[Ein Brief aus Bautzen -- Version intégrale pour l'examen]
Lieber Lukas,

ich schreibe dir aus Bautzen, der Hauptstadt der Sorben. Hier in der Lausitz leben wir Sorben schon seit über tausend Jahren, aber unsere Sprache und Kultur sind heute stark gefährdet. Obwohl die Verfassung uns Minderheitenrechte garantiert, ist der Alltag oft anders. In den Behörden wird fast nur Deutsch gesprochen, und viele junge Leute ziehen in den Westen, weil es hier wenig Arbeit gibt.

Mein Großvater hat mir immer gesagt: „Wer seine Sprache verliert, verliert seine Seele.“ Deshalb besuche ich das sorbische Gymnasium, wo wir zweisprachig unterrichtet werden. Ich wünschte, alle Kinder in der Region hätten die Chance, Sorbisch zu lernen, nicht nur wir. Wenn die Politik mehr Geld für zweisprachige Kindergärten gäbe, wäre die Zukunft sicherer.

Trotzdem bin ich optimistisch. Letes Jahr habe ich an einem Festival teilgenommen, wo junge Sorben moderne Musik auf Sorbisch gesungen haben. Das hat mir gezeigt: Die Sprache lebt, wenn wir sie nutzen. Ich hoffe, du besuchst mich bald, damit ich dir unsere Welt zeigen kann.

Dein Freund,
Jan
\end{textebox}

\begin{enumerate}
    \item \textbf{Globalverständnis :} Worum geht es in dem Brief? Nennen Sie drei Hauptthemen. (ca. 30 Wörter)
    \item \textbf{Detailverständnis :} Warum sind Sprache und Kultur der Sorben nach Jan „stark gefährdet“? Nennen Sie zwei Gründe aus dem Text. (ca. 40 Wörter)
    \item \textbf{Reformulierung / Grammatik :} Erklären Sie den Satz: „Wer seine Sprache verliert, verliert seine Seele.“ Was meint der Großvater damit? (ca. 40 Wörter)
    \item \textbf{Recht / Falsch -- Begründung :} Entscheiden Sie, ob die Aussage richtig oder falsch ist. Begründen Sie mit einem Satz aus dem Text: „Jan besucht eine deutsche Schule, in der nur Deutsch gesprochen wird.“ (Richtig / Falsch + Begründung)
    \item \textbf{Analyse linguistique (Konjunktiv II) :} Finden Sie im Text drei Formen des Konjunktiv II. Nennen Sie den Infinitiv und erklären Sie kurz, welchen Nuancen (Wunsch, Irrealität, Höflichkeit) sie hier ausdrücken. (ca. 50 Wörter)
\end{enumerate}
\end{exercice}

\begin{exercice}[Thème et Version : Traduction mixte]
\begin{enumerate}
    \item \textbf{Thème (Français $\rightarrow$ Allemand) -- ca. 60 Wörter :}
    Traduisez en allemand :
    \guillemotleft Bien que les minorités soient protégées par la loi, elles sont souvent discriminées dans la vie quotidienne. Pourtant, beaucoup de jeunes restent optimistes. S'ils avaient plus de droits, ils pourraient préserver leur culture. Je souhaite que la société montre plus de respect. \guillemotright
    \item \textbf{Version (Allemand $\rightarrow$ Français) :}
    Traduisez en français le passage suivant extrait du texte de Jan (lignes 5--9) :
    \guillemotleft Mein Großvater hat mir immer gesagt: „Wer seine Sprache verliert, verliert seine Seele.“ Deshalb besuche ich das sorbische Gymnasium, wo wir zweisprachig unterrichtet werden. Ich wünschte, alle Kinder in der Region hätten die Chance, Sorbisch zu lernen, nicht nur wir. \guillemotright
\end{enumerate}
\end{exercice}

\begin{exercice}[Expression écrite type Baccalauréat -- Sujet de rédaction]
Traitez le sujet suivant en allemand. \textbf{Nombre de mots attendu : 80--100 mots.}
Respectez la structure : Introduction (thèse) -- Arguments (2--3 avec connecteurs logiques) -- Conclusion (ouverture).

\begin{textebox}[Sujet]
\textbf{Welche Rolle spielen Sprachen für die Integration? Vergleichen Sie die Situation in Deutschland (z. B. Sorben, Einwanderer) und in Kamerun (offizielle Sprachen, langues nationales).}

\textit{Indications :}
\begin{itemize}
    \item Utilisez le vocabulaire de la leçon : \emph{die Integration, die Muttersprache, die Zweisprachigkeit, das Vorurteil, die Identität, fördern, bewahren}.
    \item Utilisez au moins une phrase avec \emph{obwohl} et une avec \emph{trotzdem}.
    \item Utilisez au moins une forme de \emph{Konjunktiv II} (souhait / irrealité).
    \item Soignez la déclinaison des adjectifs.
\end{itemize}
\end{textebox}
\end{exercice}

\newpage

\coursec{Corrigés des exercices}

\begin{corrige}[Exercice 1 — QCM : Compréhension globale du texte]
\textbf{Réponses correctes :}
\begin{enumerate}
    \item \textbf{b)} La situation de la minorité sorabe en Lusace (Brandebourg/Saxe).\\
    \emph{Justification :} toute la lettre porte sur la langue, la culture et les droits des Sorabes ; la cuisine et le Moyen Âge ne sont pas évoqués.
    \item \textbf{b)} De la colère et de la tristesse.\\
    \emph{Justification :} l'auteur dénonce la discrimination quotidienne et craint la disparition de sa langue (\all{stark gefährdet}).
    \item \textbf{b)} Le bilinguisme obligatoire dans l'administration et les écoles.\\
    \emph{Justification :} l'auteur souhaite que tous les enfants puissent apprendre le sorabe (\all{Ich wünschte, alle Kinder in der Region hätten die Chance, Sorbisch zu lernen}) ; il ne demande jamais d'interdire l'allemand ni de quitter le pays.
\end{enumerate}
\textbf{Point de méthode :} au QCM, éliminez d'abord les réponses qui exagèrent ou déforment le texte (a et c sont ici des caricatures).
\end{corrige}

\begin{corrige}[Exercice 2 — Vrai ou Faux ?]
\begin{enumerate}
    \item \textbf{Faux.} Les Sorabes sont une minorité \emph{slave}, pas germanique. Justification : \all{Die Sorben sind eine kleine slawische Minderheit in Deutschland.} (Ils sont bien \emph{reconnus}, mais l'adjectif « germanique » rend l'affirmation fausse.)
    \item \textbf{Faux.} L'auteur distingue intégration et assimilation : il veut garder sa langue. Justification : \all{Wer seine Sprache verliert, verliert seine Seele.} — pour lui, s'intégrer ne signifie pas oublier sa culture.
    \item \textbf{Vrai.} Justification : \all{Deshalb besuche ich das sorbische Gymnasium, wo wir zweisprachig unterrichtet werden.}
    \item \textbf{Vrai.} Justification : \all{Ich wünschte, alle Kinder in der Region hätten die Chance, Sorbisch zu lernen.} — \all{hätten} est un Konjunktiv II exprimant un souhait non réalisé dans le présent.
\end{enumerate}
\textbf{Point de vigilance :} une affirmation n'est vraie que si \emph{tous} ses éléments sont confirmés par le texte ; un seul détail faux (ici « germanique ») suffit à la rendre fausse.
\end{corrige}

\begin{corrige}[Exercice 3 — Vocabulaire : familles de mots]
\begin{center}
\begin{tabularx}{\linewidth}{|X|X|X|X|}
\hline
\rowcolor{popblueL} \textbf{Nom (Art. Pl.)} & \textbf{Verbe / expression} & \textbf{Adjectif} & \textbf{Français} \\
\hline
die Diskriminierung, -en & diskriminieren (jmdn. diskriminieren) & diskriminierend & la discrimination / discriminer / discriminatoire \\
\hline
die Integration, -en & integrieren (sich integrieren) & integriert & l'intégration / (s')intégrer / intégré \\
\hline
das Vorurteil, -e & Vorurteile haben (gegen + Akk.) & voreingenommen & le préjugé / avoir des préjugés / préjugé, partial \\
\hline
das Recht, -e & ein Recht haben / gewähren & rechtlich & le droit / accorder un droit / juridique \\
\hline
die Minderheit, -en & in der Minderheit sein & minderheitlich & la minorité / être en minorité / minoritaire \\
\hline
die Assimilation, -en & assimilieren (sich) & assimiliert & l'assimilation / (s')assimiler / assimilé \\
\hline
\end{tabularx}
\end{center}
\textbf{Remarques :}
\begin{itemize}
    \item \all{das Vorurteil} : attention au genre (\emph{das}), fréquente erreur des francophones qui disent *\all{die Vorurteil}. Pluriel : \all{die Vorurteile}.
    \item Le verbe *\all{sich vorurteilen} n'existe pas en allemand standard : on dit \all{Vorurteile haben gegen + Akk.} (\all{Sie hat Vorurteile gegen Ausländer.} « Elle a des préjugés contre les étrangers. ») ou l'adjectif \all{voreingenommen sein} (« être partial, avoir des préjugés »).
    \item \all{integrieren} est régulier : \all{ich integriere, integrierte, hat integriert}. Prononciation : « in-té-gri-ren ».
    \item \all{diskriminieren} se construit avec l'accusatif de la personne : \all{Man diskriminiert die Minderheit.} « On discrimine la minorité. »
\end{itemize}
\end{corrige}

\begin{corrige}[Exercice 4 — Konjunktiv II du souhait]
\begin{enumerate}
    \item \all{Ich wünschte, alle Menschen \textbf{hätten} mehr Respekt voreinander.}\\
    « Je souhaiterais que tous les hommes aient plus de respect les uns pour les autres. » — \all{hätten} = Konjunktiv II de \all{haben} (3\textsuperscript{e} pers. plur.).
    \item \all{Wenn ich Bundeskanzler \textbf{wäre}, \textbf{gäbe} ich den Minderheiten mehr Rechte.}\\
    « Si j'étais chancelier, je donnerais plus de droits aux minorités. » — \all{wäre} (sein), \all{gäbe} (geben, forme simple préférable à \all{würde geben}). Après la subordonnée en tête, le verbe de la principale est en position 1.
    \item \all{Mein Vater \textbf{könnte} gern Sorbisch sprechen, aber er hat es nie gelernt.}\\
    « Mon père aimerait savoir parler sorabe, mais il ne l'a jamais appris. » — \all{könnte} (können) + infinitif. Remarque : les noms de langues prennent la majuscule (\all{Sorbisch, Deutsch}).
    \item \all{Wäre es nicht besser, wenn wir alle uns besser \textbf{verstünden}?}\\
    « Ne serait-il pas mieux que nous nous comprenions tous mieux ? » — \all{verstünden} (verstehen, irrégulier : \all{verstand} $\rightarrow$ \all{verstände/verstünde}). Dans la subordonnée \all{wenn}, le verbe conjugué est en fin de phrase.
    \item \all{Ich \textbf{möchte} nicht, dass die Sprache \textbf{stürbe} (ou : \textbf{sterben würde}).}\\
    « Je ne voudrais pas que la langue meure. » — \all{möchte} est déjà la forme Konjunktiv II de \all{mögen} ; \all{stürbe} (forme simple littéraire de \all{sterben} : \all{starb} $\rightarrow$ \all{stürbe}) ou \all{würde sterben} (forme courante).
\end{enumerate}
\textbf{Point de vigilance :} ne jamais écrire *\all{würde hätten} ; \all{haben}, \all{sein} et les verbes modaux s'emploient toujours à la forme simple (\all{hätte, wäre, könnte}).
\end{corrige}

\begin{corrige}[Exercice 5 — Traduction thématique (Français $\rightarrow$ Allemand)]
\begin{enumerate}
    \item \all{Obwohl die Minderheiten Rechte haben, werden sie oft diskriminiert.}\\
    \emph{Grammaire :} \all{obwohl} + verbe en fin de subordonnée ; principale avec verbe en position 1 après la subordonnée.
    \item \all{Integration bedeutet nicht, dass man seine eigene Sprache vergessen muss.}\\
    \emph{Vocabulaire :} \all{eigen} = propre ; \all{bedeuten} = signifier (faux ami : « bénéficier » = \all{profitieren}).
    \item \all{Viele junge Sorben sprechen fließend Deutsch und Sorbisch.}\\
    \emph{Déclinaison :} \all{viele junge Sorben} — déclinaison forte au pluriel nominatif (\all{-e} / \all{-e}).
    \item \all{Trotz der Vorurteile ist die sorbische Kultur in der Lausitz lebendig.}\\
    \emph{Grammaire :} \all{trotz} + génitif (\all{der Vorurteile}) ; \all{sorbische} : déclinaison faible après l'article défini \all{die}.
    \item \all{Ich wünsche mir, dass die Regierung mehr für den Schutz der Minderheitensprachen tut.}\\
    \emph{Vocabulaire :} \all{der Schutz} (protection) + génitif ; variante avec souhait marqué : \all{..., dass die Regierung mehr für den Schutz der Minderheitensprachen tun würde.}
\end{enumerate}
\textbf{Point de vigilance :} « bien que » = \all{obwohl} (jamais *\all{während}, qui signifie « pendant que / tandis que ») ; « malgré » = \all{trotz} + génitif.
\end{corrige}

\begin{corrige}[Exercice 6 — Transformation \emph{obwohl} $\leftrightarrow$ \emph{trotzdem}]
\begin{enumerate}
    \item \all{Die Sorben sind eine kleine Minderheit, \textbf{trotzdem haben} sie eigene Schulen.}\\
    \emph{Règle :} \all{trotzdem} occupe la position 1, le verbe suit immédiatement (inversion).
    \item \all{\textbf{Obwohl die Sprache bedroht ist}, lernen viele Kinder sie in der Schule.}\\
    \emph{Règle :} \all{obwohl} renvoie le verbe conjugué en fin de subordonnée ; la principale commence par son verbe.
    \item \all{Er ist kein Sorbe, \textbf{trotzdem lernt} er die Sprache.}
    \item \all{\textbf{Obwohl die Gesetze die Minderheit schützen}, gibt es Probleme im Alltag.}
    \item \all{Sie war müde, \textbf{trotzdem ging} sie zur Demonstration für Minderheitenrechte.}
\end{enumerate}
\textbf{Point de vigilance :} l'erreur classique du francophone est *\all{trotzdem sie lernen...} (calque de « pourtant elles apprennent »). Après \all{trotzdem}, inversion obligatoire : \all{trotzdem lernen sie...}. Ne jamais combiner les deux : *\all{obwohl..., trotzdem...} est impossible en allemand.
\end{corrige}

\begin{corrige}[Exercice 7 — Texte à trous : déclinaison de l'adjectif]
\begin{enumerate}
    \item \all{Die Sorben sind \textbf{eine kleine} slawische Minderheit.} [M] — après \all{ein-} (nom. fém.), déclinaison mixte : \all{-e}.
    \item \all{Sie leben im (\textbf{in dem} \textbf{östlichen}) Teil von Sachsen und Brandenburg.} [F] — après \all{dem} (datif), déclinaison faible : \all{-en}.
    \item \all{\textbf{Viele} Sorben sprechen noch ihre Muttersprache.} [S] — sans article, pluriel nominatif, déclinaison forte : \all{-e}.
    \item \all{In \textbf{sorbischen} Schulen wird zweisprachig unterrichtet.} [S] — sans article, datif pluriel : \all{-en}.
    \item \all{Die Regierung gibt \textbf{finanzielle} Unterstützung für die Kultur.} [S] — sans article, accusatif féminin : \all{-e}.
    \item \all{Ohne \textbf{politischen} Willen würde die Sprache sterben.} [S] — sans article, accusatif masculin (\all{der Wille}) : \all{-en}.
    \item \all{Man sieht oft \textbf{zweisprachige} Straßenschilder in Bautzen.} [S] — sans article, accusatif pluriel : \all{-e}.
    \item \all{\textbf{Junge} Sorben sind oft stolz auf ihre Identität.} [S] — sans article, nominatif pluriel : \all{-e}.
    \item \all{Trotz \textbf{schwieriger} Geschichte haben sie überlebt.} [S] — sans article, génitif féminin : \all{-er}.
    \item \all{Ich wünsche mir \textbf{eine gute} Zukunft für die Minderheit.} [M] — après \all{eine} (acc. fém.), déclinaison mixte : \all{-e}.
\end{enumerate}
\textbf{Moyen mnémotechnique :} si l'article porte déjà la marque du cas (\all{der, die, das, dem, den...}), l'adjectif reste « faible » (\all{-e} ou \all{-en}) ; sans article, c'est l'adjectif qui porte la marque (déclinaison forte).
\end{corrige}

\begin{corrige}[Exercice 8 — Appariement : structures du souhait]
\textbf{Associations correctes : 1--d, 2--c, 3--b, 4--a, 5--e.}
\begin{enumerate}
    \item \all{Ich wünschte, die Regierung \textbf{förderte} (ou : \textbf{fördern würde}) die Sprachen stärker.}\\
    « Je souhaiterais que le gouvernement soutienne davantage les langues. »
    \item \all{Wenn alle Menschen mehr Respekt vor Minderheiten \textbf{hätten}, ...}\\
    « Si tous les hommes avaient plus de respect envers les minorités... »
    \item \all{Wäre ich Minister, ... \textbf{könnten} wir friedlicher zusammenleben.}\\
    « Si j'étais ministre, nous pourrions vivre ensemble plus paisiblement. » — \all{könnten} (können, 1\textsuperscript{re} pers. plur.).
    \item \all{Es wäre schön, wenn es mehr zweisprachige Kitas \textbf{gäbe}.}\\
    « Ce serait bien s'il y avait plus de crèches bilingues. » — \all{es gäbe} = « il y aurait » (Konjunktiv II de \all{es gibt}).
    \item \all{Die Kinder \textbf{würden} gern Sorbisch \textbf{lernen}, wenn die Eltern sie \textbf{unterstützten}.}\\
    « Les enfants apprendraient volontiers le sorabe si les parents les soutenaient. » — ici \all{würde} + infinitif est la forme la plus naturelle pour un verbe faible ; dans la subordonnée irréelle, \all{unterstützten} (Konjunktiv II de \all{unterstützen}).
\end{enumerate}
\textbf{Point de vigilance :} \all{die Kita} (die Kindertagesstätte) = la crèche ; \all{zweisprachig} = bilingue. Ne pas confondre \all{gäbe} (il y aurait) et \all{gäbe es} en tête de phrase interrogative. Dans une phrase irréelle, \emph{les deux} verbes sont au Konjunktiv II : *\all{wenn die Eltern sie unterstützen} (indicatif) est ici une faute.
\end{corrige}

\begin{corrige}[Exercice 9 — Compréhension écrite type Baccalauréat]
\textbf{Barème indicatif (sur 20) :} Q1 : 3 pts ; Q2 : 4 pts ; Q3 : 4 pts ; Q4 : 3 pts ; Q5 : 6 pts. On évalue : exactitude du contenu, correction grammaticale, richesse du lexique, respect du nombre de mots.

\textbf{Proposition de réponses modèles :}
\begin{enumerate}
    \item \all{In dem Brief geht es um die gefährdete sorbische Sprache und Kultur in der Lausitz, um die Diskriminierung im Alltag und um Jans Hoffnung auf eine zweisprachige Zukunft.} (29 mots)\\
    « La lettre traite de la langue et de la culture sorabes menacées, de la discrimination quotidienne et de l'espoir d'un avenir bilingue. »
    \item \all{Erstens wird in den Behörden fast nur Deutsch gesprochen. Zweitens ziehen viele junge Leute in den Westen, weil es in der Lausitz wenig Arbeit gibt. Deshalb sprechen immer weniger Menschen Sorbisch.} (33 mots)\\
    « D'abord, on ne parle presque qu'allemand dans les administrations. Ensuite, beaucoup de jeunes partent vers l'Ouest (les Länder de l'ouest) parce qu'il y a peu de travail en Lusace. »
    \item \all{Der Großvater meint: Die Sprache ist das Zentrum der Identität. Wenn ein Volk seine Muttersprache verliert, verliert es auch seine Kultur, seine Geschichte und seine Seele, also alles, was es ausmacht.} (31 mots)\\
    « Le grand-père veut dire que la langue est le centre de l'identité : perdre sa langue maternelle, c'est perdre sa culture, son histoire, son âme. »
    \item \textbf{Falsch.} \all{Jan besucht das sorbische Gymnasium, wo die Schüler zweisprachig unterrichtet werden, nicht eine rein deutsche Schule.}\\
    « Faux. Jan fréquente le lycée sorabe, où l'enseignement est bilingue, et non une école purement allemande. »
    \item \all{Im Text findet man: \textbf{hätten} (Infinitiv: haben) — ein Wunsch für die Gegenwart; \textbf{gäbe} (geben) — eine irreale Bedingung; \textbf{wäre} (sein) — Folge einer irrealen Bedingung. Alle drücken Wünsche aus, die jetzt nicht real sind.} (35 mots)\\
    « On trouve : \emph{hätten} (avoir) — souhait présent ; \emph{gäbe} (donner) — condition irréelle ; \emph{wäre} (être) — conséquence d'une condition irréelle. Toutes expriment des souhaits non réalisés. »
\end{enumerate}
\textbf{Points de vigilance :} répondre par des phrases complètes (verbe conjugué en position 2) ; ne pas recopier de longs passages sans reformulation ; pour la question 5, toujours donner l'infinitif du verbe demandé.
\end{corrige}

\begin{corrige}[Exercice 10 — Thème et Version]
\textbf{1. Thème (proposition de traduction) :}\\
\all{Obwohl die Minderheiten durch das Gesetz geschützt sind, werden sie oft im Alltag diskriminiert. Trotzdem bleiben viele junge Menschen optimistisch. Wenn sie mehr Rechte hätten, könnten sie ihre Kultur bewahren. Ich wünsche mir, dass die Gesellschaft mehr Respekt zeigt.} (38 mots)

\textbf{Points de grammaire vérifiés :}
\begin{itemize}
    \item \all{Obwohl ... geschützt sind} : verbe en fin de subordonnée (passif : \all{geschützt sind}).
    \item \all{Trotzdem bleiben} : inversion après \all{trotzdem}.
    \item \all{hätten ... könnten} : Konjunktiv II dans la phrase conditionnelle irréelle.
    \item \all{die Gesellschaft} = la société (faux ami : « la société » commerciale = \all{die Firma}).
\end{itemize}

\textbf{2. Version (proposition de traduction) :}\\
« Mon grand-père m'a toujours dit : "Celui qui perd sa langue perd son âme." C'est pourquoi je fréquente le lycée sorabe, où nous recevons un enseignement bilingue. Je souhaiterais que tous les enfants de la région aient la chance d'apprendre le sorabe, pas seulement nous. »

\textbf{Points de vigilance :} \all{deshalb} = « c'est pourquoi » (et non « donc » seul) ; \all{wo} est ici un relatif (« où / dans lequel ») ; \all{hätten die Chance} = « aient la chance » (Konjunktiv II rendu par le subjonctif français).
\end{corrige}

\begin{corrige}[Exercice 11 — Expression écrite type Baccalauréat]
\textbf{Barème indicatif (sur 20) :} respect du sujet et de la structure : 5 pts ; correction grammaticale (déclinaisons, ordre des mots) : 6 pts ; vocabulaire thématique : 4 pts ; connecteurs et contraintes (\all{obwohl}, \all{trotzdem}, Konjunktiv II) : 3 pts ; orthographe (majuscules, Umlaute) : 2 pts.

\textbf{Proposition de rédaction modèle (environ 95 mots) :}

\all{Sprachen spielen eine große Rolle für die Integration. In Deutschland lernen die Sorben Sorbisch und Deutsch, obwohl ihre Sprache gefährdet ist. Auch Einwanderer brauchen Deutsch, um Arbeit zu finden, trotzdem sollten sie ihre Muttersprache bewahren. In Kamerun sind Französisch und Englisch die offiziellen Sprachen, aber die nationalen Sprachen wie Ewondo oder Duala sind wichtig für die Identität. Ich wünschte, alle Schulen förderten die Zweisprachigkeit. Wer zwei Sprachen spricht, hat zwei Heimaten. Integration bedeutet nicht Assimilation, sondern Respekt vor der kulturellen Vielfalt.}

\textbf{Traduction :} « Les langues jouent un grand rôle dans l'intégration. En Allemagne, les Sorabes apprennent le sorabe et l'allemand, bien que leur langue soit menacée. Les immigrés aussi ont besoin de l'allemand pour trouver du travail ; pourtant, ils devraient préserver leur langue maternelle. Au Cameroun, le français et l'anglais sont les langues officielles, mais les langues nationales comme l'ewondo ou le douala sont importantes pour l'identité. Je souhaiterais que toutes les écoles encouragent le bilinguisme. Qui parle deux langues a deux patries. L'intégration ne signifie pas l'assimilation, mais le respect de la diversité culturelle. »

\textbf{Points de vigilance :}
\begin{itemize}
    \item Vérifier chaque adjectif : \all{die offiziellen Sprachen} (faible), \all{nationale Sprachen} (forte), \all{der kulturellen Vielfalt} (faible, datif fém. après \all{vor}).
    \item Une seule phrase avec \all{obwohl} et une avec \all{trotzdem} suffisent : inutile d'en mettre partout.
    \item \all{die Heimat} (plur. \all{Heimaten}) = la patrie, le pays natal ; beau mot de conclusion.
    \item Respecter la fourchette 80--100 mots : compter avant de recopier au propre.
\end{itemize}
\end{corrige}

\newpage

\coursec{Fiche bilan}

\begin{popfigure}
\begin{tikzpicture}[mindmap, grow cyclic, every node/.style={concept, text=white, font=\small\bfseries},
root concept/.append style={concept color=popblue, font=\normalsize\bfseries},
level 1/.append style={level distance=3.4cm, sibling angle=51.4},
level 2/.append style={level distance=1.9cm, sibling angle=45, font=\scriptsize}]
\node[root concept, align=center]{AL04\\ Das Problem\\ der Minderheiten}
child[concept color=poporange]{ node[align=center]{Texte\\ \all{Ein Brief aus Bautzen}}
  child{ node[align=center]{Sorben:\\ Sprache, Kultur} }
  child{ node[align=center]{Alltag,\\ Diskriminierung} }
}
child[concept color=popgreen]{ node[align=center]{Lexique\\ Minderheiten}
  child{ node[align=center]{\all{die Minderheit,}\\ \all{das Vorurteil}} }
  child{ node[align=center]{\all{die Integration,}\\ \all{die Rechte}} }
}
child[concept color=poppurple]{ node[align=center]{Déclinaison\\ de l'adjectif}
  child{ node[align=center]{faible, mixte,\\ forte} }
  child{ node[align=center]{4 cas,\\ 3 genres} }
}
child[concept color=popgold]{ node{Concession}
  child{ node[align=center]{\all{obwohl}\\ verbe en fin} }
  child{ node[align=center]{\all{trotzdem}\\ verbe en 2\up{e}} }
}
child[concept color=poppink]{ node[align=center]{Konjunktiv II\\ du souhait}
  child{ node[align=center]{\all{hätte, wäre,}\\ \all{könnte}} }
  child{ node[align=center]{\all{Wenn ich}\\ \all{doch...!}} }
}
child[concept color=popdark]{ node{Production}
  child{ node[align=center]{résumé,\\ réponses} }
  child{ node[align=center]{court\\ paragraphe} }
};
\end{tikzpicture}
\legende{Carte mentale de la leçon AL04 : le problème des minorités}
\end{popfigure}

\begin{bilan}
\textbf{1. Lexique clé à connaître par cœur}
\begin{itemize}
\item \all{die Minderheit, -en} : la minorité ; \all{die Mehrheit, -en} : la majorité
\item \all{die Diskriminierung, -en} : la discrimination ; \all{diskriminieren} : discriminer
\item \all{das Vorurteil, -e} : le préjugé ; \all{das Recht, -e} : le droit ; \all{die Rechte} : les droits
\item \all{die Integration, -en} : l'intégration ; \all{sich integrieren} : s'intégrer
\item \all{die Sprache, -n} : la langue ; \all{die Muttersprache, -n} : la langue maternelle
\item \all{die Ausländerfeindlichkeit} : la xénophobie ; \all{tolerant} : tolérant ; \all{die Toleranz} : la tolérance
\item \all{die Gleichberechtigung} : l'égalité des droits ; \all{schützen} : protéger
\end{itemize}

\textbf{2. Grammaire à connaître par cœur}
\begin{center}
\begin{tabularx}{\linewidth}{|X|X|}
\hline
\rowcolor{popblueL} \textbf{Règle} & \textbf{Exemple} \\
\hline
Adjectif après \all{der/die/das} (faible) : \textbf{-e} au nominatif singulier, \textbf{-en} presque partout ailleurs & \all{das kleine Dorf} ; \all{in dem kleinen Dorf} \\
\hline
Adjectif après \all{ein/kein/mein} (mixte) : \textbf{-er/-es/-e} au nominatif (l'adjectif « porte » la marque du genre) & \all{ein kleines Dorf} ; \all{keine deutsche Schule} \\
\hline
Adjectif sans article (forte) : l'adjectif prend la terminaison de l'article & \all{guter Wein} ; \all{frisches Brot} \\
\hline
\all{obwohl} + verbe conjugué \textbf{à la fin} (subordonnée) & \all{Obwohl sie deutsch spricht, fühlt sie sich fremd.} \\
\hline
\all{trotzdem} : adverbe, verbe en \textbf{2\up{e} position} (inversion) & \all{Trotzdem bleibt sie optimistisch.} \\
\hline
Konjunktiv II du souhait : \all{hätte, wäre, könnte, gäbe} + \all{doch/nur} & \all{Hätte ich doch mehr Zeit!} ; \all{Wenn es doch keine Vorurteile gäbe!} \\
\hline
\end{tabularx}
\end{center}

\textbf{3. Méthodes express}
\begin{enumerate}
\item \textbf{Comprendre le texte} : lire le titre et la source, repérer le thème (minorité, lieu, problème), puis lecture détaillée avec le glossaire.
\item \textbf{Répondre aux questions} : reformuler avec ses propres mots, citer un élément du texte, verbe en 2\up{e} position.
\item \textbf{Résumé} : 3 à 5 phrases, présent de narration, connecteurs \all{zuerst, dann, außerdem, schließlich}.
\item \textbf{Production guidée} : donner son avis avec \all{Ich finde, dass...} / \all{Meiner Meinung nach...} + une concession avec \all{obwohl} + un souhait au Konjunktiv II.
\end{enumerate}
\end{bilan}

\begin{aretenir}[Checklist avant le Bac]
\begin{itemize}
\item[$\square$] Je connais le vocabulaire des minorités avec article et pluriel (\all{die Minderheit, -en} ; \all{das Vorurteil, -e}).
\item[$\square$] Je sais décliner l'adjectif dans les trois déclinaisons (faible, mixte, forte) aux quatre cas.
\item[$\square$] Je place correctement le verbe après \all{obwohl} (fin de phrase) et après \all{trotzdem} (inversion).
\item[$\square$] Je ne confonds pas \all{obwohl} (conjonction) et \all{trotzdem} (adverbe).
\item[$\square$] Je forme le Konjunktiv II : \all{ich hätte, ich wäre, ich könnte, es gäbe}.
\item[$\square$] J'exprime un souhait avec \all{Wenn ... doch/nur ...!} ou \all{Hätte ich doch...!}.
\item[$\square$] Je mets la majuscule à tous les noms allemands (\all{die Sprache, das Recht}).
\item[$\square$] Je sais résumer le texte \all{Ein Brief aus Bautzen} en quelques phrases au présent.
\item[$\square$] Je sais donner mon avis sur l'intégration avec \all{Meiner Meinung nach} et une phrase de concession.
\item[$\square$] Je vérifie la place du verbe (2\up{e} position en principale, fin en subordonnée) dans chaque phrase produite.
\end{itemize}
\end{aretenir}

\findecours

\end{document}
